%!TEX TS-program = xelatex
%!TEX TS-options = -shell-escape
\newcommand{\verfasser}{Phil Steinhorst}					% Verfasser des Dokuments
\newcommand{\fach}{Elementare Zahlentheorie}				% Titel der Vorlesung								
\newcommand{\shortFach}{EZT}								% Titel kurz
\newcommand{\prof}{Prof.\,Dr.\,Falko Lorenz}					% Dozent
\newcommand{\untertitel}{gelesen von \prof}					% Untertitel
\newcommand{\semester}{Wintersemester 2014/2015}					% Semester
\newcommand{\homepage}{http://wwwmath.uni-muenster.de/u/karin.halupczok/elZTWiSe14/}	% Vorlesungshomepage
% % % % % % % % % % % % % % % % % % % % % % % % % % % % % % % % %
%		PhistScript.tex											%
%		XeLaTeX-Konfiguration für Skripte						%
%																%
%		Author: Phil Steinhorst									%
% % % % % % % % % % % % % % % % % % % % % % % % % % % % % % % % %

\RequirePackage[thinlines]{easybmat} %-- muss aufgrund eines Bugs vor etex und tikz geladen werden
\documentclass[a4paper, twoside, headsepline, index=totoc,toc=listof, fontsize=9pt, cleardoublepage=empty, headinclude, DIV=13, BCOR=13mm, titlepage]{scrartcl}

\usepackage{scrtime} 	% KOMA, Uhrzeit ermoeglicht
\usepackage{scrpage2} 	% wie fancyhdr, nur optimiert auf KOMA-Skript, leicht andere Syntax
\usepackage{etoolbox}

\usepackage[usenames, table, x11names]{xcolor}	% Farben. Vor TikZ laden!

% % % TikZ-Pakete % % %
\usepackage{tikz}
\usepackage{tikz-cd}
\usetikzlibrary{external}
\tikzset{>=latex}
\usetikzlibrary{shapes,arrows,intersections}
\usetikzlibrary{calc,3d}
\usetikzlibrary{decorations.pathreplacing,decorations.markings}
\usetikzlibrary{angles}
\tikzexternalize[prefix=tikz/, up to date check=diff]	% -shell-escape-Flag nötig!

% % % verwende LuaLaTeX, wegen dynamischer Speicherallokation
\tikzset{external/system call={xelatex \tikzexternalcheckshellescape 
    -halt-on-error -interaction=batchmode -jobname "\image" "\texsource"}}
    
% % % tikzexternalize für tikzcd deaktivieren, da inkompatibel
\AtBeginEnvironment{tikzcd}{\tikzexternaldisable}
\AtEndEnvironment{tikzcd}{\tikzexternalenable}

% % % um Inkompatibilitaeten von quotes und polyglossia bzw. babel zu vermeiden
\tikzset{
  every picture/.append style={
    execute at begin picture={\shorthandoff{"}},
    execute at end picture={\shorthandon{"}}
  }
}
\usetikzlibrary{quotes}
\usepackage{pgfplots}

% % % Mathe-Zeugs
\usepackage{mathtools} 						% beinhaltet amsmath
\mathtoolsset{showonlyrefs, centercolon}
\usepackage{wasysym}
\usepackage{amssymb} 						% zusätzliche Symbole
\usepackage{latexsym} 						% zusätzliche Symbole
\usepackage{stmaryrd} 						% für Blitz
\usepackage{nicefrac} 						% schräge Brüche
\usepackage{cancel} 						% Befehle zum Durchstreichen
\usepackage{extarrows}						% mehr Pfeile
\usepackage{mathdots}
\DeclareSymbolFont{bbold}{U}{bbold}{m}{n}
\DeclareSymbolFontAlphabet{\mathbbold}{bbold}
\newcommand{\mathds}[1]{\mathbb{#1}} 		% um Kompatibilitaet mit frueheren Benutzung von dsfont herzustellen
\def\mathul#1#2{\color{#1}\underline{{\color{black}#2}}\color{black}} %farbiges Untersteichen im Mathe-Modus

% % % XeTeX & Fonts
\usepackage{mathspec} 						% beinhaltet fontspec 
\usepackage{polyglossia} 					% babel-Ersatz
\setmainlanguage[spelling=new,babelshorthands=false]{german}
\newcommand\glqq{"}
\newcommand\grqq{"}
\defaultfontfeatures{Mapping=tex-text, WordSpace={1.4},Extension=.otf} %
\setmainfont{SourceSansPro}[%
	UprightFont=*-Regular,
	BoldFont=*-Semibold,
	ItalicFont=*-LightIt,
	BoldItalicFont=*-SemiboldIt,
	SmallCapsFont={LinLibertine_R.otf},
	SmallCapsFeatures={Letters=SmallCaps}]
\setsansfont{SourceSansPro}[%
	UprightFont=*-Regular,
	BoldFont=*-Semibold,
	ItalicFont=*-LightIt,
	BoldItalicFont=*-SemiboldIt,
	SmallCapsFont={LinLibertine_R.otf},
	SmallCapsFeatures={Letters=SmallCaps}]
\setmonofont{SourceCodePro}[Scale=0.9,UprightFont=*-Regular, BoldFont=*-Semibold, ItalicFont=*-Light] 
\usepackage{xltxtra}
\usepackage{fontawesome}

% % % Misc.
\usepackage[neverdecrease]{paralist}
\usepackage[german=quotes]{csquotes}
\usepackage{booktabs}
\usepackage{wrapfig}
\usepackage{float}
\usepackage[margin=10pt, font=small, labelfont=sf, format=plain, indention=1em]{caption}
\captionsetup[wrapfigure]{name=Abb. }
\usepackage{stackrel}
\usepackage{ifthen}
\usepackage{multicol}
%\flushbottom

% % % Unterstreichung
\usepackage[normalem]{ulem}
\setlength{\ULdepth}{1.8pt}

% % % Indexverarbeitung
\usepackage{makeidx}
\newcommand{\bet}[1]{\textbf{#1}}
\newcommand{\Index}[1]{\textbf{#1}\index{#1}}
\makeindex
\renewcommand{\indexpagestyle}{scrheadings}

% % % Marginnote & todonotes
\usepackage{marginnote}
\renewcommand*{\marginfont}{\color{gray} \footnotesize }
\usepackage[textsize=small]{todonotes}
\makeatletter
\renewcommand{\todo}[2][]{\tikzexternaldisable\@todo[#1]{#2}\tikzexternalenable}
\makeatother

% % % Konfiguration von Hyperref pdfstartview=FitH, 
\usepackage[hidelinks, pdfpagelabels,  bookmarksopen=true, bookmarksnumbered=true, linkcolor=black, urlcolor=RoyalBlue3, plainpages=false, hypertexnames=false, citecolor=black, hypertexnames=true, pdfauthor={\verfasser}, pdfborderstyle={/S/U}, linkbordercolor=RoyalBlue3, colorlinks=true, unicode, pdfencoding=auto]{hyperref}

\newcommand{\RM}[1]{\MakeUppercase{\romannumeral #1{}}} 	% Römische Zahlen
\usepackage{ellipsis}										% Punkte

% % % % % % % % % % % % % % % % % % % % % % % % % % % % % % % % %
%		PhistMath.tex											%
%		Weitere Mathe-Befehle									%
%																%
%		Author: Phil Steinhorst									%
% % % % % % % % % % % % % % % % % % % % % % % % % % % % % % % % %

% % % Buchstaben und Zahlen
\newcommand{\aff}{\mathbb{A}}
\newcommand{\CC}{\mathbb{C}}
\newcommand{\FF}{\mathbb{F}}
\newcommand{\HH}{\mathbb{H}}
\newcommand{\KK}{\mathbb{K}}
\newcommand{\NN}{\mathbb{N}}
\newcommand{\OO}{\mathbb{O}}
\newcommand{\QQ}{\mathbb{Q}}
\newcommand{\RR}{\mathbb{R}}
\newcommand{\ZZ}{\mathbb{Z}}
\newcommand{\bigO}{\mathcal{O}}					% Landau-O
\newcommand{\ind}{1\hspace{-0,8ex}1} 			% Indikatorfunktion (Doppeleins)

% % % Abk�rzungen
\newcommand{\ab}[1]{\overline{#1}}					% Abschluss
\newcommand{\assoz}{\mathrel{\hat=}}				% assoziiert
\newcommand{\bewrueck}{\glqq$\Leftarrow$\grqq:} 	% Beweis R�ckrichtung
\newcommand{\bewhin}{\glqq$\Rightarrow$\grqq:}		% Beweis Hinrichtung
\newcommand{\borel}{\mathfrak{B}}					% Borelsche Sigma-Algebra
\newcommand{\maxid}{\mathfrak{m}}					% maximales Ideal
\newcommand{\setone}{\{1\}}							% Einsmenge
\newcommand{\leb}{\lambda \hspace{-0,95ex}\lambda}	% Lebesgue-Ma� (Doppel-Lambda)
\newcommand{\Lp}{\mathcal{L}}						% L^p-R�ume
\newcommand{\NT}{\trianglelefteq}					% Normalteiler
\newcommand{\setnull}{\{0\}}						% Nullmenge
\newcommand{\weak}{\rightharpoonup}					% schwache Konvergenz
\newcommand{\weaks}{\overset{*}{\rightharpoonup}}	% schwache *-Konvergenz
\newcommand{\salg}{\mathfrak{A}}					% Sigma-Algebra (Skript-A)
\newcommand{\zyklot}[1]{#1^{(\infty)}}				% zyklotomische Erweiterung


% % % Operatoren
\DeclareMathOperator{\Alt}{Alt} 					% Alternierende n-Linearform
\DeclareMathOperator{\Aut}{Aut} 					% Automorphismen
\DeclareMathOperator{\Bil}{Bil} 					% Bilinearformen
\DeclareMathOperator{\bild}{Bild} 					% Bild
\DeclareMathOperator{\Char}{char} 
\DeclareMathOperator{\dom}{dom} 					% Domain
\DeclareMathOperator{\diam}{diam}					% Durchmesser
\DeclareMathOperator{\dist}{dist} 					% Distanz
\DeclareMathOperator{\eqs}{\mathrel{\widehat{=}}}	% entspricht
\DeclareMathOperator{\diver}{div} 					% Gradient
\DeclareMathOperator{\EPK}{EPK} 					% Einpunktkompaktifizierung
\DeclareMathOperator{\End}{End} 					% Endomorphismen
\DeclareMathOperator{\esssup}{esssup}				% essentielles Supremum
\DeclareMathOperator{\Gal}{Gal}	 					% Galoisgruppe
\DeclareMathOperator{\ggT}{ggT} 					% ggT
\DeclareMathOperator{\GL}{GL}						% allgemeine lineare Gruppe
\DeclareMathOperator{\grad}{grad} 					% Gradient
\DeclareMathOperator{\Grad}{Grad} 					% Grad
\DeclareMathOperator{\Hess}{Hess} 					% Hesse-Matrix
\DeclareMathOperator{\Hom}{Hom} 					% Homomorphismen
\DeclareMathOperator{\id}{id} 						% identische Abbildung
\DeclareMathOperator{\im}{im} 						% image
\DeclareMathOperator{\Jac}{Jac} 					% Jacobson-Radikal
\DeclareMathOperator{\Kern}{Kern}					% Kern
\DeclareMathOperator{\kgV}{kgV} 					% kgV
\DeclareMathOperator{\Koker}{Koker} 				% Kokern
\DeclareMathOperator{\Cov}{Cov} 					% Kovarianz
\DeclareMathOperator{\Mod}{Mod} 					% Moduln
\DeclareMathOperator{\modu}{mod} 					% Modulo
\DeclareMathOperator{\ord}{ord} 					% Ordnung
\DeclareMathOperator{\der}{\partial}				% Partielle Ableitung
\DeclareMathOperator{\pot}{\mathcal{P}}				% Potenzmenge
\DeclareMathOperator{\prlim}{\varprojlim\limits}	% projektiver Limes
\DeclareMathOperator{\Quot}{Quot}					% Quotientenring
\DeclareMathOperator{\Rang}{Rang} 					% Rang
\DeclareMathOperator{\rot}{rot} 					% Rotation
\DeclareMathOperator{\sgn}{sgn} 					% Signum
\DeclareMathOperator{\Spec}{Spec} 					% Spektrum
\DeclareMathOperator{\SL}{SL} 						% Spezielle lineare Gruppe
\DeclareMathOperator{\SO}{SO} 						% Spezielle orthogonale Gruppe
\DeclareMathOperator{\SU}{SU} 						% Spezielle unit�re Gruppe
\DeclareMathOperator{\Spur}{Spur} 					% Spur
\DeclareMathOperator{\supp}{supp} 					% Tr�ger
\DeclareMathOperator{\Sym}{Sym} 					% Symmetrische Gruppe
\DeclareMathOperator{\tr}{tr} 						% trace

% % % Klammerungen
\DeclarePairedDelimiter{\abs}{\lvert}{\rvert}		% Betrag
\DeclarePairedDelimiter{\ceil}{\lceil}{\rceil}		% aufrunden
\DeclarePairedDelimiter{\floor}{\lfloor}{\rfloor}	% aufrunden
\DeclarePairedDelimiter{\sprod}{\langle}{\rangle}	% spitze Klammern
\DeclarePairedDelimiter{\enbrace}{(}{)}				% runde Klammern
\DeclarePairedDelimiter{\benbrace}{\lbrack}{\rbrack}			% eckige Klammern
\DeclarePairedDelimiter{\penbrace}{\{}{\}}			% geschweifte Klammern

% % % Norm
\DeclarePairedDelimiter\doppelstrich{\Vert}{\Vert}
\newcommand{\norm}[2][\relax]{
\ifx#1\relax \ensuremath{\doppelstrich*{#2}}
\else \ensuremath{\doppelstrich*{#2}_{#1}}
\fi}							% selbst definierte Mathe-Befehle

\setlength\parindent{0pt}             % ohne Einrueckung
\renewcommand{\baselinestretch}{1.2}
% % % Konfiguration von scrheadings
\setheadsepline{1pt}[\color{black}]
\pagestyle{scrheadings}
\clearscrheadfoot

% % % Kopf-/Fußzeilenlayout
\providecommand{\shortFach}{\fach}
\ohead{\small \leftmark}
\ofoot[{ \small \thepage}]{ \small \thepage} 
\automark{section}

% Metadaten
\title{\fach}
\author{\verfasser}
\subtitle{\untertitel}
\date{\semester}

% Inhaltsverzeichnis
\usepackage[tocindentauto]{tocstyle}
\usetocstyle{KOMAlike}	

% ntheorem package
\usepackage[hyperref]{ntheorem}			% Lade ntheorem mit hyperref-Workaround
\theoremstyle{break}					% Zeilenumbruch nach Theoremkopf
\theoremheaderfont{\usekomafont{disposition}} % Benutze die gleiche Schriftart, wie für \section \subsection etc.
\theorembodyfont{\normalfont}			% nicht kursiv
\theorempreskipamount0.5cm				% Abstand vor Theorem
\theorempostskipamount0.5cm				% Abstand nach Theorem

% Workaround für Tabulatoren im Theorem-Verzeichnis
\makeatletter
\def\thm@@thmline@name#1#2#3#4{%
        \@dottedtocline{-2}{0em}{2.3em}%
                   {\makebox[\widesttheorem][l]{#1 \protect\numberline{#2}}#3}%
                   {#4}}
\@ifpackageloaded{hyperref}{
\def\thm@@thmline@name#1#2#3#4#5{%
    \ifx\\#5\\%
        \@dottedtocline{-2}{0em}{2.3em}%
            {\makebox[\widesttheorem][l]{#1 \protect\numberline{#2}}#3}%
            {#4}
    \else
        \ifHy@linktocpage\relax\relax
            \@dottedtocline{-2}{0em}{2.3em}%
                {\makebox[\widesttheorem][l]{#1 \protect\numberline{#2}}#3}%
                {\hyper@linkstart{link}{#5}{#4}\hyper@linkend}%
        \else
            \@dottedtocline{-2}{0em}{2.3em}%
                {\hyper@linkstart{link}{#5}%
                  {\makebox[\widesttheorem][l]{#1 \protect\numberline{#2}}#3}\hyper@linkend}%
                    {#4}%
        \fi
    \fi}
}
\makeatother
\newlength\widesttheorem
\AtBeginDocument{
  \settowidth{\widesttheorem}{Proposition 10.10\quad}	% Referenzstring für die Breite der ersten Spalte
}

\newcommand{\qed}{\ifmmode \tag*{$\square$} \else \hfill $\square$ \fi} % qed-Symbol.

							% XeLaTeX-Konfiguration für Skripte

% % % Dokumentspezifische Einstellungen % % %
\setcounter{tocdepth}{2}				% Tiefe im Inhaltsverzeichnis (1: nur Sections, 2: auch subsections...)

\numberwithin{equation}{section}			% Section in Equation-Nummerierung einbeziehen
\newcounter{countsatz}						% Zähler für Sätze
\numberwithin{countsatz}{section}			% Section in Satz-Nummerierung einbeziehen
\newcounter{countfalko}						% Zähler für Fs
\numberwithin{countfalko}{section}			% Section in F-Nummerierung einbeziehen
\newcounter{countbsp}						% Zähler für Beispiele
\numberwithin{countbsp}{section}			% Section in Beispiel-Nummerierung einbeziehen
\newcounter{countdef}						% Zähler für Definitionen
\numberwithin{countdef}{section}			% Section in Definition-Nummerierung einbeziehen
\newcounter{countlemma}						% Zähler für Lemmata
\numberwithin{countlemma}{section}			% Section in Lemma-Nummerierung einbeziehen

% ab hier können neue Umgebungen für Sätze definiert und benannt werden!
% Syntax: \newtheorem{interner Umgebungsname}[Zähler]{Name im Dokument}
% Möchte man Sätze und Definitionen getrennt durchnummerieren, benötigt man weitere Zähler.
\newtheorem{satz}[countsatz]{Satz}
\newtheorem{bsp}[countbsp]{Beispiel}
\newtheorem{falko}[countfalko]{F\hspace{-0.25em}}
\newtheorem{defn}[countdef]{Definition}
\newtheorem{lemma}[countlemma]{Lemma}
% % % % % % % % % % % % % % % % % % % % % % %
\newcommand{\PP}{\mathbb{P}}
\DeclareMathOperator{\li}{li}
\newcommand{\kon}{\equiv}
\newcommand{\leg}[2]{\enbrace*{\frac{#1}{#2}}}

% % % % % % % % % % % % % % % % % % % % % % %
\begin{document}
%!TEX root = ../LA2.tex
\begin{titlepage}
	% Nach einer Vorlage von http://www.LaTeXTemplates.com
	\newcommand{\HRule}{\rule{\linewidth}{0.8mm}} % Defines a new command for the horizontal lines, change thickness here

	\center % Center everything on the page
 
	\begin{minipage}{0.4\textwidth}
	\begin{flushleft}
	\includegraphics[height=1.5cm,keepaspectratio]{../!config/Bilder/wwulogo.pdf}\\[1cm]
	\end{flushleft}
	\end{minipage}
	\hfill
	\begin{minipage}{0.4\textwidth}
	\begin{flushright}
	\vspace*{0.3cm}
	\includegraphics[height=1.2cm,keepaspectratio]{../!config/Bilder/fb10logo.pdf} \
	\end{flushright}
	\end{minipage}

	\vspace{2cm}
	


	\HRule \\[0.8cm]
	{ \huge \sffamily\bfseries Lineare Algebra II}\\[0.4cm] % Title of your document
	\HRule \\[1cm]
 
	{\LARGE gelesen von} \\[.7cm]
	\textsc{\LARGE \textbf{Prof. Dr. Siegfried Echterhoff}}\\[.7cm]
	{\LARGE im Sommersemester 2016}\\[2cm]


	\vfill
	{\Large Vorlesungsmitschrift von Phil Steinhorst} \\[.5cm]
	{\large Stand: \today}
	
\end{titlepage}							% Titelseite ausgeben
%!TEX root = ../EKK_SS15.tex
\begin{abstract}
	\section*{Vorwort}
	\label{sec:preface}
	Der vorliegende Text ist eine inhaltliche Aufbereitung zur Vorlesung \textit{Elliptische Kurven und Kryptographie}, gelesen von PD Dr. Karin Halupczok an der WWU Münster im Sommersemester 2015. Der Inhalt entspricht weitestgehend dem handschriftlichen Skript, welches auf der Vorlesungswebsite bereitsgestellt wird. Dieses Werk ist daher keine Eigenleistung des Autors und wird nicht von der Dozentin der Veranstaltung korrekturgelesen. Für die Korrektheit des Inhalts wird keinerlei Garantie übernommen. Bemerkungen, Korrekturen und Ergänzungen kann man folgenderweise loswerden:
	\begin{itemize}
		\item persönlich durch Überreichen von Notizen oder per E-Mail
		\item durch Abändern der entsprechenden \TeX-Dateien und Versand per E-Mail an mich
		\item direktes Mitarbeiten via GitLab. Dieses Skript befindet sich im \texttt{LaTeX-WWU}-Repository von Jannes Bantje:
		\begin{center}
			\url{https://gitlab.com/JaMeZ-B/latex-wwu}
		\end{center}
	\end{itemize}
	
	\subsection*{Literatur}
	\label{sub:lit}
	\begin{itemize}
		\item Blake, Seroussi, Smart: Elliptic curves in cryptography \cite{BlakeSeroussiSmart}
		\item Menezes, van Oorschot, Vanstone: Handbook of applied cryptography \cite{MenezesOorschotVanstone}
		\item Silverman: The arithmetic of elliptic curves \cite{SilvermanEC}
		\item Silverman: A friendly introduction to number theory, chap. 40-45 \cite{SilvermanNT}
		\item Washington: Elliptic curves, number theory and cryptography \cite{Washington}
		\item Werner: Elliptische Kurven in der Kryptographie \cite{Werner}
	\end{itemize}
	
	\subsection*{Kommentar der Dozentin}
	In der Vorlesung beschäftigen wir uns mit den arithmetischen und geometrischen Eigenschaften elliptischer Kurven sowie deren Anwendungen in der Kryptographie. Dabei werden wir auch einen Vergleich mit Anwendungen der elementaren Zahlentheorie in der Kryptographie ziehen. Wir verfolgen eine elementare Herangehensweise, d.h. Kenntnisse der algebraischen Geometrie und der Funktionen- oder Zahlentheorie werden nicht benötigt. Es genügen die Vorkenntnisse aus den Grundvorlesungen.
	\newpage	
	\subsection*{Vorlesungswebsite}
	\label{sub:link}
	Das handgeschriebene Skript sowie weiteres Material findet man unter folgendem Link:
	\begin{center}
		\url{http://wwwmath.uni-muenster.de/u/karin.halupczok/ellKKSoSe15/}
	\end{center}
	
	\vfill
	\begin{flushright}
		Phil Steinhorst \\
		p.st@wwu.de
	\end{flushright}
	\newpage
\end{abstract}
\cleardoubleemptypage
\tableofcontents
\newpage											% Inhaltsverzeichnis ausgeben
\section{Fundamentalsatz der elementaren Arithmetik}
\label{sec:para1}

\minisec{Terminologie}
	Sei $R$ ein kommutativer Ring mit $1 \neq 0$. $R$ heißt \Index{Integritätsring} bzw. \bet{nullteilerfrei}, wenn gilt: \index{Nullteiler} \marginnote{14.10.\\ \ [1]}
	\[ a \cdot b = 0 \quad \Rightarrow \quad a = 0 \text{ oder } b = 0.\]

\begin{bsp} \label{bsp_integritaetsringe}
	\begin{itemize}
		\item $\ZZ$
		\item $\ZZ[\sqrt{2}] := \{a + b\sqrt{2} : a,b \in \ZZ \} \subseteq \RR$ \\
			$\ZZ[i] := \{a + bi : a,b \in \ZZ\} \subseteq \CC$ \\
			$\ZZ[\sqrt{-5}] := \dots$
		\item $K[X]$ für $K$ Körper \\
			$\ZZ[X]$
		\item $K$ Körper
		\item $\CC \sprod{z} := \penbrace*{\text{konvergente Potenzreihen } \sum\limits_{n=0}^{\infty} a_n z^n}$
		\item Nicht nullteilerfrei ist z.B. $\mathcal{C}[0,1] := \{f \colon [0,1] \rightarrow \RR \text{ stetig} \}$
	\end{itemize}
\end{bsp}

\begin{defn}[Teilbarkeit] \label{def_1.1}
	Seien $a,b \in R$. $a$ heißt ein \Index{Teiler} von $b$, wenn ein $q \in R$ existiert mit $b = qa$. Wir schreiben dann:
	\[ a | b \]
	Ist $R$ nullteilerfrei und $a \neq 0$, so ist $q$ eindeutig bestimmt.
\end{defn}

\begin{falko}[Triviale Teilbarkeitsregeln] \label{F1.1}
	\begin{enumerate}[(i)]
		\item $a | 0, 1 | a, a | a$
		\item $a | b, b | c \quad \Rightarrow \quad a | c$
		\item $a | b, a | c \quad \Rightarrow \quad a | b+c, a | b-c$
		\item $a_1 | b_1, a_2 | b_2 \quad \Rightarrow \quad a_1 a_2 | b_1 b_2$
		\item $ac | bc \quad \Rightarrow \quad a | b$, falls $c \neq 0$ und $R$ nullteilerfrei.
	\end{enumerate}	
\end{falko}

\begin{defn}[Einheit, assoziiert] \label{def_1.2}
	\begin{enumerate}[(i)]
		\item $e \in R$ heißt eine \Index{Einheit} in $R$, falls $e | 1$ gilt, d.h. falls ein $f \in R$ existiert mit $ef = 1$. $f$ ist eindeutig bestimmt. Wir setzen $e^{-1} := f$ und schreiben auch $\frac{1}{e}$ für $e^{-1}$. \\
		Wir bezeichnen die \Index{Einheitengruppe} von $R$ mit $R^\times := \{x \in R : x \text{ ist Einheit in } R\}$.
		\item $a \in R$ heißt \Index{assoziiert} zu $b \in R$, falls $a | b$ und $b | a$ gilt. Schreibe: $a \assoz b$.
	\end{enumerate}
\end{defn}

\newpage
\begin{bsp}
	\begin{enumerate}[1)]
		\item Sei $K$ ein Körper, dann ist $K^\times = K \setminus \setnull$. \qquad $\ZZ^\times = \{1,-1\}$, \qquad $K[X]^\times = K^\times$, \\
		$\mathcal{C}[0,1]^\times = \{f \in \mathcal{C}[0,1] : f(x) \neq 0 \text{ für alle } x \in [0,1]\}$, \qquad  $\ZZ[\sqrt{2}]^\times = \{\pm (1+\sqrt{2})^k : k \in \ZZ \}$ \\
		$\ZZ[X]^\times = \{1, -1\}$ \qquad $\CC \sprod{z}^\times = \penbrace*{ \sum a_n z^n \in \CC \sprod{z} : a_0 \neq 0}$
		\item $e \in R^\times \quad \Leftrightarrow \quad e | a$ für jedes $a \in R$.
	\end{enumerate}
\end{bsp}

\begin{falko} \label{F1.2}
	Sei $R$ ein Integritätsring, $a, b \in R$ und $b \neq 0$. Dann gilt:
	\[ a \assoz b \quad \Leftrightarrow \quad \exists e \in R^\times \text{ mit } b = ea \]
\end{falko}

\minisec{Beweis}
	\begin{description}
		\item[\bewrueck] $a | b, e^{-1}b = a, b | a$
		\item[\bewhin] Da $a | b$ und $b | a$, existieren $e, f \in R$, sodass $b = ea$ und $a = fb$. $\Rightarrow b = efb \Rightarrow ef = 1$, da $b \neq 0$ und $R$ nullteilerfrei. \qed
	\end{description}

\setlength{\fboxsep}{10pt}
\setlength{\fboxrule}{3pt}
\begin{center}
	\fbox{\textbf{Ab jetzt ist, wenn nichts anderes gesagt, $R$ ein Integritätsring!}}
\end{center}

\begin{defn}[unzerlegbar, irreduzibel, zusammengesetzt] \label{def_1.3}
	Sei $a \in R \setminus R^\times$. $a$ heißt \Index{unzerlegbar} oder \Index{irreduzibel} in $R$, wenn gilt:
	\[ a = bc \text{ in } R \quad \Rightarrow \quad b \in R^\times \text{ oder } c \in R^\times. \]
	Andernfalls heißt $a$ \bet{zerlegbar}, \bet{zusammengesetzt} oder \bet{reduzibel}.
\end{defn}

\minisec{Bemerkung}
	$a$ unzerlegbar $\quad \Leftrightarrow \quad $ jeder Teiler von $a$ ist Einheit oder assoziiert zu $a$ \\
	$a$ zerlegbar $\quad \Leftrightarrow \quad a$ hat echten Teiler, d.h. einen Teiler, der weder eine Einheit ist noch assoziiert zu $a$

\setcounter{countdef}{2}
\begin{defn}[Primzahl] \label{def_1.3'}
	Ein $p \in \ZZ$ heißt \Index{Primzahl}, wenn $p \in \NN$ und $p$ unzerlegbar in $\ZZ$. Wir bezeichnen mit $\PP$ die Menge der Primzahlen von $\ZZ$. $a$ unzerlegbar in $\ZZ \Leftrightarrow a = p$ oder $a = -p$ mit $p \in \PP$.
\end{defn}

\minisec{Bemerkung}
	$a \in \ZZ$ sei zerlegbar, $a \neq 0$. Dann gibt es eine Primzahl $p$ mit $p | a$ und $p \leq \sqrt{|a|}$.

\begin{defn}[Zerlegung in unzerlegbare Faktoren] \label{def_1.4}
	Wir sagen, $a \in R$ besitzt in $R$ eine \Index{Zerlegung in unzerlegbare Faktoren}, wenn
	\begin{equation}
	\begin{aligned}
		a = ep_1p_2\dots p_r \text{ mit } e \in R^\times \text{ und } p_1,\dots,p_r \text{ unzerlegbar} \label{eq_def_1.4}
	\end{aligned}
	\end{equation}
	\eqref{eq_def_1.4} heißt eine Zerlegung von $a$ in unzerlegbare Faktoren. Auch $r = 0$ ist erlaubt.
\end{defn}

\begin{falko} \label{F1.3}
	In $\ZZ$ besitzt jedes $a \neq 0$ eine Zerlegung in unzerlegbare Faktoren.
\end{falko}	

\setcounter{countfalko}{2}
\begin{falko} \label{F1.3'}
	Jede natürliche Zahl $a > 1$ besitzt eine Zerlegung $a = p_1p_2 \dots p_r$ mit Primzahlen $p_1,\dots,p_r$ und $r \geq 1$.
\end{falko}

\minisec{Bemerkung}
	\begin{enumerate}[1)]
		\item Die Aussage F1.3 gilt auch für die Beispiele \ref{bsp_integritaetsringe}, mit Ausnahme von $\mathcal{C}[0,1]$.
		\item Sei $R$ ein Integritätsring, der die \Index{Teilbarkeitsbedingung für Hauptideale} erfüllt, so besitzt jedes $a \neq 0$ aus $R$ eine Zerlegung in unzerlegbare Faktoren.
		\item Primzahlen sind die multiplikativen Bausteine (Atome) von $\NN$.
		\item Im Beispiel $\CC\sprod{z}$ von oben gibt es (bis auf Assoziiertheit) nur das einzige unzerlegbare Element $z$. Dieses ist ein \Index{Primelement} (der Begriff folgt weiter unten).
	\end{enumerate}
	
\begin{satz}[Existenz unendlich vieler Primzahlen] \label{satz_1.1}
	Es gibt unendlich viele Primzahlen.
\end{satz}

\textbf{Bemerkungen} \\
	Es sei $p_1,p_2,\dots$ die aufsteigend sortierte Folge der Primzahlen.
	\begin{enumerate}[1)]
		\item $a_n := p_1p_2\dots p_n +1$ ist Primzahl für $n \leq 5$, aber z.B. nicht für $n = 6$. Unklar ist, ob unendlich viele $a_n$ Primzahlen oder keine Primzahlen sind.
		\item Für $x \in \RR_{>0}$ definieren wir:
		\[ \pi(x) := \# \penbrace{p \in \PP : p \leq x}\]
	\end{enumerate}

\minisec{Primzahlsatz (Gauß, Legendre)}
	\begin{equation}
		\pi(x) \sim \frac{x}{\log x}, \text{ d.h. } \lim\limits_{x \rightarrow \infty} \frac{\pi(x)}{x / \log x} = 1 \label{eq_primzahlsatz1}
	\end{equation}
	\begin{equation}
		\pi(x) \sim \int_{2}^{x} \frac{1}{\log t} dt =: \li(x) \label{eq_primzahlsatz_2}
	\end{equation}
	\begin{equation}
	\pi(x) > \frac{x}{\log x} \text{ für alle } x \geq 17 \label{eq_primzahlsatz_3}
	\end{equation}
	\begin{equation}
	\pi(n) > \frac{n}{\log n} \text{ für alle } n \in \NN, n \geq 11 \label{eq_primzahlsatz_4}
	\end{equation}
	
\begin{defn}[eindeutige Zerlegung] \label{def_1.5}
	Sei $R$ ein kommutativer Ring mit $1 \neq 0$. Wir sagen, $a \in R \setminus \setnull$ hat eine \bet{eindeutige} \Index{Zerlegung in unzerlegbare Faktoren}, wenn $a$ eine Zerlegung
	\[ a = ep_1p_2 \dots p_r \]
	in unzerlegbare Faktoren besitzt und eine solche im folgendem Sinne eindeutig ist: Ist auch
	\[ a = e'p_1'p_2'\dots p'_{r'} \]
	eine solche Zerlegung, so gilt $r = r'$ und nach Umnummerierung $p_i' \assoz p_i$ für alle $1 \leq i \leq r$.
\end{defn}

\newpage
\begin{falko}\label{F1.4}
	In dem Integritätsring $R$ besitze jedes Element $a \neq 0$ eine Zerlegung in unzerlegbare Faktoren. Dann sind äquivalent: \begin{enumerate}[(i)]
		\item Jedes $a \neq 0$ aus $R$ hat eindeutige Zerlegung in unzerlegbare Faktoren.
		\item Ist $p$ unzerlegbar, so gilt: $p | ab \Rightarrow p | a$ oder $p | b$.
	\end{enumerate}
\end{falko}

\begin{defn}[Primelement] \label{def_1.6}
	Sei $R$ ein kommutativer Ring mit $1 \neq 0$. Ein $p \in R\setminus R^\times$ heißt \Index{Primelement} von $R$, wenn für alle $a, b \in R$ gilt:
	\begin{equation}
		p | ab \quad \Leftrightarrow \quad p | a \text{ oder } p | b \label{eq_def_1.6}
	\end{equation}
\end{defn}

\minisec{Bemerkung}
	\begin{enumerate}[1)]
		\item $0$ ist Primelement in $R \Leftrightarrow R$ ist Integritätsring
		\item In einem Integritätsring $R$ gilt: Jedes Primelement $p \neq 0$ ist unzerlegbar.
	\end{enumerate}

\begin{lemma}
	Seien $a, b \in \NN$. Sei $m = \kgV(a,b) \in \NN$. Dann gilt:
	\[ a|c \text{ und } b|c \quad \Rightarrow \quad m | c \]
	$m$ ist also auch minimal bzgl. der Teilbarkeitsrelation $|$.
\end{lemma}

\begin{falko}[Satz von Euklid] \label{F1.5}
	Jede Primzahl $p$ ist ein Primelement von $\ZZ$, d.h. es gilt stets \eqref{eq_def_1.6}. \marginnote{17.10.\\ \ [2]} (Das gleiche gilt für $-p$, also für jedes unzerlegbare Element von $\ZZ$.) \index{Satz von Euklid}
\end{falko}

\minisec{Fundamentalsatz der elementaren Arithmetik}
	In $\ZZ$ hat jedes $a \neq 0$ eine eindeutige Zerlegung in unzerlegbare Faktoren. \index{Fundamentalsatz der elementaren Arithmetik}

\minisec{Bemerkung}
	Eindeutige Zerlegung in unzerlegbare Faktoren hat man zum Beispiel auch für die Ringe $\ZZ[\sqrt{2}], \ZZ[i], K[X]$ und $K$ für $K$ Körper, $\ZZ[X]$ und $\CC\sprod{z}$, nicht aber für $\ZZ[\sqrt{-5}]$:
	\[ 3 \cdot 3 = 9 = (2+ \sqrt{-5})(2 - \sqrt{-5}) \]
	Dies sind zwei wesentlich verschiedene Zerlegungen in unzerlegbare Faktoren.

\begin{defn}[Exponent] \label{def_1.7}
	Sei $p$ eine Primzahl und $a \in \ZZ \setminus \setnull$. Dann heißt
	\[ w_p(a) :=\max\{k \in \NN_0 : p^k |a \} \]
	der \Index{Exponent} von $p$ in $a$. Wir setzen $w_p(0) := \infty$.
\end{defn}

\begin{falko}[Eigenschaften der Exponentfunktion] \label{F1.6}
	Die Funktion $w_p \colon \ZZ \rightarrow \NN_0 \cup \{\infty\}$ hat folgende Eigenschaften:
	\begin{enumerate}[(i)]
		\item $w_p(a+b) \geq \min(w_p(a),w_p(b))$ und Gleichheit, falls $w_p(a) \neq w_p(b)$.
		\item $w_p(ab) = w_p(a) + w_p(b)$
	\end{enumerate}
\end{falko}

\begin{satz}[Fundamentalsatz der elementaren Arithmetik] \label{satz_1.2}
	Für jedes $a \in \ZZ \setminus \setnull$ gilt $w_p(a) > 0$ nur für endlich viele $p$. Es ist
	\begin{equation}
		a = \sgn(a) \cdot \prod\limits_p^{} p^{w_p(a)} \label{eq_satz_1.2}
	\end{equation}
\end{satz}

\minisec{Bemerkung}
	\begin{enumerate}[1)]
		\item $w_p$ lässt sich eindeutig zu einer Abbildung $w_p \colon \QQ \rightarrow \ZZ \cup \{\infty\}$ fortsetzen, sodass (ii) für alle $a,b \in \QQ$ gilt. Es gilt dann auch (i). Für $a \in \QQ \setminus \setnull$ ist $w_p(a) \neq 0$ nur für endlich viele $p$, und die Formel \eqref{eq_satz_1.2} gilt entsprechend.  Ferner gilt: $a \in \ZZ \Leftrightarrow w_p(a) \geq 0$ für alle $p$.
		\item Sei
		\[ \NN_0^{(\PP)} := \{ (e_p)_{p \in \PP} : e_p \in \NN_0, e_p = 0 \text{ für fast alle } p \}. \]
		Nach Satz \ref{satz_1.2} sind $(\NN,\cdot)$ und $(\NN_0^{(\PP)},+)$ zwei zueinander isomorphe Halbgruppen. Nach Bemerkung 1) sind $\QQ^\times$ und $\{1,-1\} \times \ZZ^{(\PP)}$ sogar zwei zueinander isomorphe Gruppen.
	\end{enumerate}
	
\begin{defn}[faktorieller Ring, Vertretersystem für Primelemente] \label{def_1.8}
	Ein Integritätsring $R$ heißt \Index{faktoriell}, wenn jedes $a \in R \setminus \setnull$ eine eindeutige Zerlegung in unzerlegbare Faktoren hat. Man spricht dann auch von eindeutiger Primfaktorzerlegung in $R$. \\
	$P$ heißt \bet{Vertretersystem für die Primelemente $\neq 0$} von $R$, wenn:
	\begin{enumerate}[(1)]
		\item Zu jedem Primelement $q \neq 0$ von $R$ gibt es ein $p \in P$ mit $q \assoz p$. 
		\item Für $p, p' \in P$ mit $p \assoz p'$ gilt $p = p'$, d.h. $p$ in (1) ist eindeutig bestimmt durch $q$.
	\end{enumerate}
	Für $R = \ZZ$ nehme man stets $P = \PP$. Für $K$ Körper und $R = K[X]$ nimmt man $P = \{p \in K[X] : p \text{ irreduzibel und normiert}\}$. \index{Primelement} \index{Vertretersystem}
\end{defn}

\begin{falko} \label{F1.7}
	Sei $R$ faktoriell und $P$ ein Vertretersystem für Primelemente. Es gibt zu jedem $p \in P$ eine Funktion $w_p \colon R \rightarrow \NN_0 \cup \{\infty\}$ mit den Eigenschaften (i) und (ii) aus F\ref{F1.6}, sodass gilt: \begin{enumerate}[a)]
		\item Für jedes $a \in R \setminus \setnull$ ist $w_p(a) > 0$ nur für endlich viele $p \in P$.
		\item Für jedes $a \in R \setminus \setnull$ gilt
		\begin{equation}
			a = e \prod\limits_{p \in P}^{} p^{w_p(a)} \label{eq_F1.7}
		\end{equation}
		mit eindeutigem $e \in \RR^\times$.
	\end{enumerate}
\end{falko}

\begin{defn}[ggT und kgV] \label{def_1.9}
	Sei $R$ ein kommutativer Ring mit $1 \neq 0$. Gegeben $a_1, \dots, a_n \in R$.
	\begin{enumerate}[a)]
		\item Ein $d \in R$ heißt ein \Index{größter gemeinsamer Teiler} (ggT) von $a_1,\dots,a_n$, falls:
		\begin{center}
			1. \quad $d | a_i$ für alle $i$ \hspace{3cm} 2. \quad $t | a_i$ für alle $i \Rightarrow t | d$
		\end{center}
		\item Ein $m \in R$ heißt ein \Index{kleinstes gemeinsames Vielfaches} (kgV) von $a_1,\dots,a_n$, falls:
		\begin{center}
			1. \quad $a_i | m$ für alle $i$ \hspace{3cm} 2. \quad $a_i | c$ für alle $i \Rightarrow m | c$
		\end{center}
	\end{enumerate}
\end{defn}

\newpage
\minisec{Bemerkung}
	\begin{enumerate}[1)]
		\item $d, d'$ ggT von $a_1, \dots, a_n \Rightarrow d \assoz d'$ und $m, m'$ kgV von $a_1, \dots, a_n \Rightarrow m \assoz m'$
		\item Im Allgemeinen ist die Existenz eines ggT und kgV nicht gesichert. In faktoriellen Ringen existieren sie aber immer, siehe dazu folgende Feststellung.
	\end{enumerate}
	
\begin{falko} \label{F1.8}
	Sei $R$ faktoriell, $P$ wie oben. Es gelten: \marginnote{21.10.\\ \ [3]}
	\begin{enumerate}[(i)]
		\item $a | b \quad \Leftrightarrow \quad w_p(a) \leq w_p(b)$ für alle $p \in P$.
		\item Für $a_1, \dots, a_n \in R$ setze:
		\[d := \prod\limits_{p \in P} p^{\min(w_p(a_1),\dots,w_p(a_n))} =: (a_1,\dots,a_n) \]
		\[m := \prod\limits_{p \in P} p^{\max(w_p(a_1),\dots,w_p(a_n))} =: [a_1,\dots,a_n] \]
		Hierbei setze $p^\infty = 0$. Dann ist $d$ ein ggT von $a_1,\dots,a_n$ und $m$ ein kgV von $a_1,\dots,a_n$.
		\item $a,b \in R$. Dann ist $a,b \assoz [a,b] \cdot (a,b)$ und $m \assoz \frac{ab}{(a,b)}$, wenn $a,b$ nicht beide 0.
		\item $a_1,\dots,a_n$ paarweise teilerfremd, d.h. $(a_i,a_j) = 1$ für $i \neq j \Leftrightarrow [a_1,\dots,a_n] \simeq a_1a_2\dots a_n$.
		\item $(a_i,b)=1$ für $1 \leq i \leq n \Rightarrow (a_1a_2\dots a_n,b) = 1$
		\item $(a_1f,\dots,a_nf) \simeq (a_1,\dots,a_n)f, [a_1f,\dots,a_nf] \simeq [a_1,\dots,a_n]f$
		\item $((a_1,\dots,a_n),a_{n+1}) = (a_1,\dots,a_n,a_{n+1}), [[a_1,\dots,a_n],a_{n+1}] = [a_1,\dots,a_n,a_{n+1}]$
	\end{enumerate}
\end{falko}

\minisec{Bemerkung (Verallgemeinerung von (iii)}
	Seien $a_1,\dots,a_n \in R$ gegeben. Wähle $q_1,\dots,q_n$ und $c$ aus $R$ mit
	\[ a_1q_1 = a_2q_2 = \dots = a_n q_n = c \]
	(z.B. $c= a_1a_2 \dots a_n, q_i = \prod\limits_{j \neq i} a_j$). Dann gilt
	\[ c \assoz (a_1,\dots,a_n)[q_1,\dots,q_n] \]
	
\begin{falko} \label{F1.9}
	Sei $n \in \NN, a \in \ZZ$. Ist $X^n = a$ lösbar in $\QQ$, so ist $X^n = a$ auch lösbar in $\ZZ$. Anders ausgedrückt: Ist $a \in \ZZ$ keine $n$-te Potenz in $\ZZ$, so ist $a$ auch keine $n$-te Potenz in $\QQ$.
\end{falko}

\minisec{Anwendung}
	$\sqrt{2}$ ist irrational, denn $2$ ist kein Quadrat in $\ZZ$ aus Größengründen, also ist $2$ nach F\ref{F1.9} auch kein Quadrat in $\QQ$, d.h. $\sqrt{2} \in \QQ$.
	
\newpage
\minisec{Korollar}
	Sei $n \in \NN, a \in \NN$. Dann sind äquivalent:
	\begin{enumerate}[(i)]
		\item $a$ ist $n$-te Potenz in $\ZZ$.
		\item $n | w_p(a)$ für alle $p$.
		\item $a$ ist $n$-te Potenz in $\QQ$.
	\end{enumerate}

\begin{falko}[Verallgemeinerung von F\ref{F1.9}] \label{F1.10}
	Gegeben sei ein normiertes Polynom $f(X) \in \ZZ[X]$. Ist dann $b$ eine Nullstelle von $f$ mit $b \in \QQ$, so ist notwendigerweise $b \in \ZZ$ und außerdem ist $b$ ein Teiler des Absolutkoeffizienten $a_0$ von $f$.
\end{falko}
\newpage
\section{Der euklidische Algorithmus}
\label{sec:para2}
	Sei $R$ kommutativer Ring mit $1 \neq 0$. Für beliebiges $a \in R$ betrachte man die Menge der Vielfachen von $a \in R$, also
	\[ Ra := \{xa : x \in R\} = \{b \in R : a|b\} \]
	Die Teilmenge $I = Ra$ hat folgende Eigenschaften:
	\begin{enumerate}[(i)]
		\item $0 \in I$
		\item $b_1,b_2 \in I \Rightarrow b_1+b_2 \in I$
		\item $c \in R, b \in I \Rightarrow cb \in I$
	\end{enumerate}
	
\begin{defn}[Ideal, Hauptideal] \label{def_2.1}
	Eine Teilmenge $I$ von $R$ heißt ein \Index{Ideal} in $R$, falls die Eigenschaften (i), (ii), (iii) erfüllt sind. $I$ heißt \Index{Hauptideal}, wenn es ein $a \in R$ gibt mit $I = Ra$. Wir verwenden die Bezeichnung
	\[ (a) := Ra \]
	und nennen $(a)$ das von $a \in R$ erzeugte Hauptideal.
\end{defn}

\minisec{Bemerkung}
	\begin{enumerate}[(1)]
		\item $(b) \subseteq (a) \Leftrightarrow a | b$
		\item $a \assoz b \Leftrightarrow (a) = (b)$
		\item $c$ ist gemeinsames Vielfaches von $a_1,\dots,a_n \Leftrightarrow (c) \subseteq (a_1) \cap \dots \cap (a_n)$
		\item $m$ ist ein kgV von $a_1,\dots,a_n \Leftrightarrow (a_1) \cap \dots \cap (a_n) = (m)$
		\item $d$ ist ein gemeinsamer Teiler von $a_1,\dots,a_n \Leftrightarrow (a_i) \subseteq (d)$ für $1 \leq i \leq n$
		\item $d$ ist ein gemeinsamer Teiler von $a_1,\dots,a_n \Leftrightarrow Ra_1 + Ra_2 + \dots + Ra_n \subseteq (d)$
		\item $d$ ist ein ggT von $a_1,\dots,a_n \Leftrightarrow (d)$ ist das kleinste Hauptideal mit $Ra_1 + \dots + Ra_n \subseteq (d)$.
	\end{enumerate}
	Ein ggT lässt sich also idealtheoretisch nicht so einfach charakterisieren wie oben ein kgV durch (4). Am schönsten wäre es, wenn $Ra_1 + \dots + Ra_n$ ein Hauptideal wäre, dann würde (7) übergehen in:
	\[ d \text{ ist ein ggT von } a_1,\dots,a_n \Leftrightarrow Ra_1+Ra_2+ \dots + Ra_n = (d) \]
	
\begin{defn}[Hauptidealring] \label{def_2.2}
	Ein Integritätsring $R$ heißt ein \Index{Hauptidealring}, wenn jedes Ideal $I$ von $R$ ein Hauptideal ist.
\end{defn}

\minisec{Bezeichnung}
	Für Elemente $a_1,\dots,a_n$ in einem beliebigen kommutativen Ring $R$ mit $1 \neq 0$ setze
	\[ (a_1,\dots,a_n) := Ra_1 + \dots + Ra_n \]
	Man nennt $(a_1,\dots,a_n)$ das von $a_1,\dots,a_n$ erzeugte Ideal in $R$.
	
\begin{falko}[Satz vom größten gemeinsamen Teiler] \label{F2.1}
	Sei $R$ ein Hauptidealring.\marginnote{24.10.\\ \ [4]} Dann gilt: Zu jedem System $a_1,\dots,a_n$ von Elementen aus $R$ existiert ein ggT $d$ von $a_1,\dots,a_n$ und jedes solche $d$ besitzt eine Darstellung der Gestalt
	\begin{equation}
		d = x_1a_1+\dots+x_na_n \quad \text{mit } x_i \in R \label{eq_F2.1}
	\end{equation}
	Wir sagen, in $R$ gelte der \Index{Satz vom größten gemeinsamen Teiler}.
\end{falko}

\minisec{Bemerkung}
	Sei $R$ ein beliebiger Integritätsring. Ist $d$ ein gemeinsamer Teiler von $a_1,\dots,a_n$ aus $R$ und gibt es eine Darstellung der Form \eqref{eq_F2.1}, so ist $d$ ein ggT von $a_1,\dots,a_n$.

\begin{satz} \label{satz_2.1}
	$\ZZ$ ist ein Hauptidealring.
\end{satz}

\minisec{Definition (Gaußklammer)}
	Für $x \in \RR$ setze
	\[ [x] = \max \{g \in \ZZ : g \leq x\} \in \ZZ \]
	$[x]$ ist charakterisiert durch folgende zwei Eigenschaften: \begin{enumerate}[(1)]
		\item $[x] \in \ZZ$
		\item $[x] \leq x < [x] + 1$
	\end{enumerate}

\begin{falko}[Division mit Rest in $\ZZ$] \label{F2.2}
	Gegeben $a, b \in \ZZ$, $a \neq 0$. Dann gibt es eine Darstellung
	\begin{equation}
		b = qa +r \quad \text{mit } 0 \leq r < |a| \text{ und } q,r \in \ZZ \label{eq_F2.2}
	\end{equation}
\end{falko}

\minisec{Bemerkung}
	\begin{enumerate}[1)]
		\item Die Darstellung \eqref{eq_F2.2} ist eindeutig.
		\item Es gibt eine Darstellung
		\begin{equation}
			b = qa + r \quad \text{mit } |r| < |a|; q,r \in \ZZ, \label{eq_D'}
		\end{equation}
		doch diese ist nicht mehr eindeutig, z.B. $27 = 4 \cdot 6 + 3 = 5 \cdot 6 - 3$.
		\item Es gibt eine Darstellung
		\begin{equation}
			b = qa + r \quad \text{mit } -\frac{|a|}{2} < r \leq \frac{|a|}{2}; q,r \in \ZZ,
		\end{equation}
		und diese ist eindeutig.
		\item Es gibt eine Darstellung
		\[ b = qa + r \quad \text{mit } |r| \leq \frac{|a|}{2}; q,r \in \ZZ, \]
		doch diese ist nicht eindeutig, falls $a$ gerade.
	\end{enumerate}

\begin{defn}[euklidischer Ring] \label{def_2.3}
	Ein Integritätsring $R$ heißt ein \Index{euklidischer Ring}, falls eine Funktion $\nu\colon R \rightarrow \NN_0$ mit $\nu(0) = 0$ existiert, sodass gilt: Zu $a, b \in R$ mit $a \neq 0$ existieren $q,r \in R$ mit
	\[b = qa + r \text{ und } \nu(r) < \nu(a) \]
\end{defn}

\minisec{Beispiele}
	\begin{enumerate}[(1)]
		\item $R = \ZZ$ mit $\nu(a) = |a|$.
		\item $R = K[X], K$ Körper, mit $\nu(g) = \deg(g) + 1$ für $g \neq 0$, $\nu(0) = 0$.
		\item $R = \ZZ[i]$ mit $\nu(z) = N(z) = z \overline{z} = |z|^2$.
	\end{enumerate}
	
\begin{falko} \label{F2.3}
	Jeder euklidische Ring ist ein Hauptidealring.
\end{falko}

\begin{falko} \label{F2.4}
	Jeder Hauptidealring ist faktoriell.
\end{falko}

Im Folgenden sei $R$ ein euklidischer Ring mit euklidischer Normfunktion $\nu$. Allgemein gilt folgende elementare Umformung:
\begin{equation}
	(a_1,a_2,\dots,a_n) = (a_1, a_2-y_2a_1,\dots,a_n-y_na_1) \text{ für bel. } y_i \in R \label{eq_U}
\end{equation}

\minisec{Euklidischer Algorithmus}
	Gegeben $a_1,\dots,a_n \in R$. Wir wollen $d \in R$ bestimmen mit \index{Euklidischer Algorithmus}
	\[ (a_1,\dots,a_n) = (d) \]
	Sind alle $a_i = 0$, so ist $d = 0$ und wir sind fertig. Sei daher ohne Einschränkung
	\[ a_1 \neq 0 \text{ und } \nu(a_1) \leq \nu(a_i), \text{ falls } a_i \neq 0 \]
	Sei $a_i = q_ia_1 + r_i$ mit $\nu(r_i) < \nu(a_1)$ für $i \geq 2$. Dann ist
	\[ (a_1,\dots,a_n) \stackrel{\eqref{eq_U}}{=} (a_1,r_2,\dots,r_n) \]
	Fortsetzung des Verfahrens liefert
	\[ (d,0,0,\dots,0) = (d) \]
	
\minisec{Beispiel}
	\[\begin{array}{rl}
		(\textcolor{red}{27},63,114) & 63 = 2 \cdot 27 + 9, 114 = 4 \cdot 27 + 6 \\ 
		= (27,9,\textcolor{red}{6}) & 27 = 4 \cdot 6 + 3, 9 = 1 \cdot 6 + 3 \\ 
		= (\textcolor{red}{3},3,6) & 3 = 1 \cdot 3 + 0, 6 = 2 \cdot 3 + 0 \\ 
		= (3, 0, 0) = (3) & 
	\end{array}\]

\minisec{Beispiel im Fall $\boldsymbol{n=2}$}
	Sei $a,b \in R \setminus \setnull$. \marginnote{28.10.\\ \ [5]}
	\[\begin{array}{lll}
		b = q_0 a + r_1 & \nu(r_1) < \nu(a) & \text{Falls } r_1 = 0, \text{ dann Schluss. Sonst weiter:} \\ 
		a = q_1 r_1 + r_2 & \nu(r_2) < \nu(r_1) &  \\ 
		r_1 = q_2 r_2 + r_3 & \nu(r_3) < \nu(r_2) &  \\ 
		\qquad \vdots &  &  \\ 
		r_{n-2} = q_{n-1} r_{n-1} + r_n & \nu(r_n) < \nu(r_{n-1}) &  \\ 
		r_{n-1} = q_n r_n + 0 &  & 
	\end{array} \]
	Also:
	\[ (a,b) = (a,r_1) = (r_1,r_2) = \dots = (r_{n-1},r_n) = (r_n) \]
	
\begin{falko} \label{F2.5}
	$r_n$ ist ein größter gemeinsamer Teiler von $a$ und $b$. Es ist
	\[ r_n = xa + yb \text{ mit } x,y \in R, \]
	wobei $x$ und $y$ aus obiger Rechnung rekursiv bestimmbar sind.
\end{falko}

\minisec{Bemerkung}
\begin{enumerate}[1)]
	\item Von neuem erhalten wir für jeden euklidischen Ring also den Satz vom größten gemeinsamen Teiler. (Satz \ref{F2.1})
	\item Sei $R = \ZZ$. Verlangen wir $0 \leq r_i$ in obiger Rechnung, so sind $q_0,q_1,\dots,q_n$ sowie die $r_1,\dots,r_n$ eindeutig bestimmt.
\end{enumerate}

\minisec{Beispiel}
	Sei $a = 84, b=133$.
	\begin{equation}
	\begin{aligned}
		133 &= 1 \cdot 84 + 49 \\
		84 &= 1 \cdot 49 + 35 \\
		49 &= 1 \cdot 35 + 14 \\
		35 &= 2 \cdot 14 + 7 \\
		14 &= 2 \cdot 7	\qquad \Rightarrow n=4, r_4 = 7
	\end{aligned}
	\end{equation}
	Also ist $(133,84) = (7)$.

Wir können den euklidischen Algorithmus für $a,b$ auch wie folgt aufschreiben:
	\[\begin{array}{lll}
	\frac{b}{a} = q_0 + \frac{r_1}{a} & q_0 = \benbrace*{\frac{b}{a}} & 0 < \frac{r_1}{a} < 1, \text{ falls } r_1 \neq 0 \\ 
	\frac{a}{r_1} = q_1 + \frac{r_2}{r_1} & q_1 = \benbrace*{\frac{a}{r_1}} &  \\ 
	\frac{r_1}{r_2} = q_2 + \frac{r_3}{r_2} & &  \\ 
	\qquad \vdots &  &  \\ 
	\frac{r_{n-2}}{r_{n-1}} = q_{n-1} + \frac{r_n}{r_{n-1}} & &  \\ 
	\frac{r_{n-1}}{r_n} = q_n &  & 
	\end{array} \]
Zusammengefasst erhalten wir die \Index{Kettenbruchentwicklung} von $\frac{b}{a}$:
	\[ \frac{b}{a} = q_0 + \cfrac{1}{q_1 + \cfrac{1}{\textcolor{white}{q_2 + \cfrac{\textcolor{black}{\ddots}}{\textcolor{black}{+\cfrac{1}{q_{n-1}+\cfrac{1}{q_n}}}}}}} \]
	
Statt einer rationalen Zahl sei jetzt $\alpha$ allgemeiner eine beliebige reelle Zahl. \\
Es ist $\alpha = [\alpha] + \varepsilon$ mit $0 \leq \varepsilon < 1$. Falls $\alpha \notin \ZZ$, d.h. $\varepsilon > 0$, setze $q_0 := [\alpha]$ und $\rho_1 := \frac{1}{\varepsilon}$. Dann:
\[ \begin{array}{lll}
\alpha = q_0 + \frac{1}{\rho_1} & \text{ mit } \rho_1 > 1. & \text{Falls } \rho_1 \notin \ZZ, \text{ so setze } [\rho_1] =: q_1 \\ 
\rho_1 = q_1 + \frac{1}{\rho_2} & \text{ mit } \rho_2 > 1. & \text{usw.} \\ 
\qquad \vdots &  &  \\ 
\rho_k = q_k + \frac{1}{\rho_{k+1}} & \text{ mit } \rho_{k+1} > 1.  & 
\end{array} \]
Abbrechen, wenn $\rho_{n+1} \in \ZZ$, sonst weiter. Jedenfalls:
\[ \alpha = \frac{b}{a} = q_0 + \cfrac{1}{q_1 + \cfrac{1}{\textcolor{white}{q_2 + \cfrac{\textcolor{black}{\ddots}}{\textcolor{black}{+\cfrac{1}{q_k+\cfrac{1}{\rho_{k+1}}}}}}}} \]

\begin{defn}[Kettenbruch, $k$-ter Rest] \label{def_2.4}
	\begin{enumerate}[1)]
		\item $q_0,q_1,\dots,q_n$ seien reelle Zahlen mit $q_1,\dots,q_n > 0$. Unter dem \bet{endlichen Kettenbruch} \index{Kettenbruch}
		\begin{equation}
			[q_0;q_1,\dots,q_n] \label{eq_2.1}
		\end{equation}
		mit den Teilquotienten $q_i$ verstehen wir sowohl das $(n+1)$-Tupel $(q_0,q_1,\dots,q_n)$, als auch seinen wie folgt definierten Wert:
		\begin{equation}
			[q_0;q_1,\dots,q_n] = q_0 + \cfrac{1}{q_1 + \cfrac{1}{\textcolor{white}{q_2 + \cfrac{\textcolor{black}{\ddots}}{\textcolor{black}{+\cfrac{1}{q_n}}}}}} \label{eq_2.2}
		\end{equation}
		Für $0 \leq k \leq n$ nennen wir den Kettenbruch
		\begin{equation}
			\rho_k := [q_k;q_{k+1},\dots,q_n] \label{eq_2.3}
		\end{equation}
		den \bet{$k$-ten Rest} des Kettenbruchs \eqref{eq_2.1}. Für den Wert \eqref{eq_2.1} des Kettenbruchs \eqref{eq_2.1} gilt: \index{k-ter Rest@$k$-ter Rest}
		\begin{equation}
			[q_0;q_1,\dots,q_n] = [q_0;q_1,\dots,q_{k-1},\rho_k] \text{ für } 0 \leq k \leq n \label{eq_2.4}
		\end{equation}
		Man kann den Wert \eqref{eq_2.2} des Kettenbruchs \eqref{eq_2.1} durch \eqref{eq_2.4} mit \eqref{eq_2.3} rekursiv definieren: Es ist $[q_0] = q_0, [q_0;q_1] = q_0 + \frac{1}{q_1}$, also:
		\begin{equation}
			[q_0;q_1,\dots,q_n] = [q_0; \rho_1] = q_0 + \frac{1}{\rho_1} \text{ für } n \geq 1 \label{eq_2.5}
		\end{equation}
		\item Gegeben sei eine Folge $(q_k)_{k \geq 0}$ in $\RR$ mit $q_k > 0$ für $k \geq 1$. Unter dem \bet{unendlichen Kettenbruch} \index{Kettenbruch}
		\begin{equation}
			[q_0;q_1,q_2,\dots] \label{eq_2.6}
		\end{equation}
		verstehen wir die Folge der
		\begin{equation}
			[q_0;q_1,\dots,q_n] \qquad n=0,1,2,\dots \label{eq_2.7}
		\end{equation}
		Falls diese Folge in $\RR$ konvergiert, so bezeichnen wir auch deren Limes mit $[q_0;q_1,q_2,\dots]$. Der unendliche Kettenbruch
		\begin{equation}
			\rho_k := [q_k;q_{k+1},\dots] \qquad k = 0,1,2,\dots \label{eq_2.8}
		\end{equation}
		heißt der \bet{$k$-te Rest} von \eqref{eq_2.6}. Formal gilt: \index{k-ter Rest@$k$-ter Rest}
		\begin{equation}
			[q_0;q_1,q_2,\dots] = [q_0;q_1,\dots,q_{k-1},\rho_k] \label{eq_2.9}
		\end{equation}
		Später werden wir sehen, dass \eqref{eq_2.9} auch für die Werte der entsprechenden Kettenbrüche gilt, wenn \eqref{eq_2.8} konvergiert.
	\end{enumerate}
\end{defn}

\begin{defn}[Näherungsbruch] \label{def_2.5}
	Jedem endlichen Kettenbruch $[q_0;q_1,\dots,q_k]$ ordnen wir rekursiv ein Paar $\binom{c}{d} \in \RR \times \RR_{>0}$ reeller Zahlen zu mit
	\begin{equation}
		[q_0;q_1,\dots,q_k] = \frac{c}{d} \label{eq_2.10}
	\end{equation}
	$k = 0$: Für $[q_0]$ sei $\binom{c}{d} = \binom{q_0}{1}$. Es gilt dann in der Tat $[q_0] = q_0 = \frac{q_0}{1}$. \\
	$k \geq 1$: Zuerst Motivation (Heuristik):
	\begin{equation}
		[q_0;q_1,\dots,q_k] = [q_0;\rho_1] = q_0 + \frac{1}{\rho_1} \label{eq_2.11}
	\end{equation}
	mit $\rho_1 = [q_1;q_2,\dots,q_k]$. Gehöre $\binom{c'}{d'}$ zu $\rho_1$. Dann gilt
	\[ [q_0;q_1,\dots,q_k] = q_0 + \frac{d'}{c'} = \frac{q_0 c' + d'}{c'} \]
	Wir ordnen nun also $[q_0;q_1,\dots,q_k]$ das Tupel
	\begin{equation}
		\binom{c}{d} = \binom{q_0 c' + d'}{c'} = \underbrace{\begin{pmatrix}
		q_0 & 1 \\ 
		1 & 0
		\end{pmatrix}}_{=: M_1} \binom{c'}{d'} \label{eq_2.12}
	\end{equation}
	zu. Dann gilt \eqref{eq_2.10}. Sei jetzt
	\begin{equation}
		[q_0;q_1,\dots] \label{eq_2.13}
	\end{equation}
	ein endlicher oder unendlicher Kettenbruch. Das dem $k$-ten Abschnitt
	\begin{equation}
		[q_0;q_1,\dots,q_k] \label{eq_2.14}
	\end{equation}
	von \eqref{eq_2.13} zugeordnete 2-Tupel
	\begin{equation}
		\binom{c_k}{d_k} \label{eq_2.15}
	\end{equation}
	heißt der \bet{$k$-te Näherungsbruch} von \eqref{eq_2.13}. Auch $\frac{c_k}{d_k}$ heißt $k$-ter Näherungsbruch von \eqref{eq_2.13}. Ist \eqref{eq_2.13} der endliche Kettenbruch $[q_0,q_1,\dots,q_n]$, so ist der $n$-te Näherungsbruch $\frac{c_n}{d_n}$ gleich dem Wert dieses Kettenbruchs. Allgemein ist $\frac{c_k}{d_k}$ der Wert des Kettenbruchs \eqref{eq_2.14}. Aus formalen Gründen definieren wir noch
	\begin{equation}
		\binom{c_{-1}}{d_{-1}} = \binom{1}{0} \qquad \binom{c_{-2}}{d_{-2}} = \binom{0}{1} \label{eq_2.16}
	\end{equation}
\end{defn}

\begin{falko}[Rekursionsformeln für Näherungsbrüche] \label{F2.6}
	Mit den Bezeichnungen wie oben gilt:
	\begin{equation}
	\begin{aligned}
		c_k &= q_k c_{k-1} + c_{k-1} \\	\label{eq_2.17}
		d_k &= q_k d_{k-1} + d_{k-2}
	\end{aligned}
	\end{equation}
	Dies schreiben wir auch in Matrizenform:
	\begin{equation}
		\binom{c_k}{d_k} = \underbrace{\begin{pmatrix}
		c_{k-1} & c_{k-2} \\ 
		d_{k-1} & d_{k-2}
		\end{pmatrix}}_{=:M_k} \binom{q_k}{1}
	\end{equation}
\end{falko}

\minisec{Bemerkung}
	$d_k > 0$ für $k \geq 0$ (vgl. Definition \ref{def_2.5}, oder auch \eqref{eq_2.17}).
	
\begin{falko} \label{F2.7}
	Mit den obigen Bezeichnungen gilt: \marginnote{31.10.\\ \ [6]}
	\begin{enumerate}[(i)]
		\item $M_{k+1} = M_k \cdot \begin{pmatrix}
			q_k & 1 \\ 1 & 0 
		\end{pmatrix}$, also:
		\item $M_{k+1} = \begin{pmatrix} q_0 & 1 \\ 1 & 0 \end{pmatrix} \begin{pmatrix} q_1 & 1 \\ 1 & 0 \end{pmatrix} \dots \begin{pmatrix} q_k & 1 \\ 1 & 0 \end{pmatrix}$
		\item $d_kc_{k-1} - c_kd_{k-1} = (-1)^k$ für $k \geq -1$
		\item $\frac{c_{k-1}}{d_{k-1}} - \frac{c_k}{d_k} = \frac{(-1)^k}{d_k d_{k-1}}$ für $k \geq 1$ \hfill $d_{-1} = 0$
		\item $d_kc_{k-2} - c_kd_{k-2} = (-1)^{k-1} q_k$ für $k \geq 0$
		\item $\frac{c_{k-2}}{d_{k-2}} - \frac{c_k}{d_k} = \frac{(-1)^{k-1} q_k}{d_k d_{k-1}}$ für $k \geq 0$, aber $k \neq 1$
	\end{enumerate}
\end{falko}

\begin{falko} \label{F2.8}
	\begin{enumerate}[(i)]
		\item $\enbrace*{ \frac{c_{2m}}{d_{2m}}}_{m \geq 0}$ ist streng monoton steigend
		\item $\enbrace*{ \frac{c_{2m+1}}{d_{2m+1}}}_{m \geq 0}$ ist streng monoton fallend
		\item $\frac{c_{2m}}{d_{2m}} < \frac{c_{2n+1}}{d_{2n+1}}$ für alle $m \geq 0, n \geq 0$
	\end{enumerate}
\end{falko}

\begin{falko} \label{F2.9}
	\begin{enumerate}[(i)]
		\item $[q_0;q_1,\dots,q_n] = \frac{\rho_k c_{k-1} + c_{k-2}}{\rho_k d_{k-1} + d_{k-2}}$ für $1 \leq k \leq n$.
		\item $[q_k;q_{k-1}, \dots, q_1] = \frac{d_k}{d_{k-1}}$ für $k \geq 1$
	\end{enumerate}
\end{falko}

\begin{falko} \label{F2.10}
	Gegeben sei ein unendlicher Kettenbruch
	\begin{equation}
		\alpha = [q_0;q_1,\dots] \label{eq_2.24}
	\end{equation}
	Dann gelten: \begin{enumerate}[(i)]
		\item $\alpha$ konvergent $\Rightarrow$ jeder Rest $\rho_n = [q_n; q_{n+1}, \dots]$ ist konvergent.
		\item $\rho_n$ konvergent für ein $n \Rightarrow \alpha$ ist konvergent
		\item Ist $\alpha$ konvergent, so gilt für die Werte
		\begin{equation}
			\alpha = \frac{c_{n-1} \rho_n + c_{n-2}}{d_{n-1} \rho_n + d_{n-2}} \quad n \geq 1 \label{eq_F2.10}
		\end{equation}
		d.h. $[q_0;q_1,\dots] = [q_0; q_1, \dots, q_{n-1}, \rho_0]$.
		\item Ist $\alpha$ konvergent, so gilt $\frac{c_{2n}}{d_{2n}} < \alpha < \frac{c_{2m+1}}{d_{2m+1}}$ für alle $n,m \geq 0$.
	\end{enumerate}
\end{falko}

\begin{falko} \label{F2.11}
	Der Wert $\alpha$ eines konvergenten unendlichen Kettenbruchs genügt den Ungleichungen
	\begin{equation}
		\abs{\alpha - \frac{c_k}{d_k}} < \frac{1}{d_k d_{k+1}} \text{ für jedes } k \geq 0 \label{eq_2.26}
	\end{equation}
\end{falko}

\begin{defn}[natürlicher Kettenbruch]
	Ein Kettenbruch $[q_0;q_1,\dots]$ -- endlich oder unendlich -- heißt \Index{natürlicher Kettenbruch}, wenn $q_k \in \ZZ$ für alle $k \geq 0$. Nach wie vor setzen wir $q_k > 0$ für $k \geq 1$ voraus!
\end{defn}

Im weiteren betrachten wir nur natürliche Kettenbrüche und sprechen dann schlechthin von Kettenbrüchen. Nach F\ref{F2.6} ist dann
\[ c_k, d_k \in \ZZ \text{ für alle } k \geq -2, \quad d_k \in \NN \text{ für } k \geq 0, \quad d_k = q_kd_{k-1} + d_{k-2} \geq d_{k+1} + 1 > d_{k-1} \text{ für } k \geq 1 \]
\begin{equation}
	d_k > d_{k-1} \text{ für } k \geq 2, \quad d_k \geq k \text{ für } k \geq 1 \label{eq_2.28}
\end{equation} 
\eqref{eq_2.28} gilt im Allgemeinen nicht für $k = 1$. Denn $d_0 = 1$, und es ist $d_1 = 1$ möglich.

\minisec{Bemerkung}
	Induktiv folgt leicht $d_k > 2^{\frac{k-1}{2}}$ für $k \geq 2$.
	
\begin{falko} \label{F2.12}
	Jeder unendliche natürliche Kettenbruch ist konvergent.
\end{falko}

\begin{falko} \label{F2.13}
	Die Näherungsbrüche eines natürlichen Kettenbruchs lassen sich nicht kürzen, d.h. $c_k$ und $d_k$ sind teilerfremd für jedes $k \geq -2$. \\
	Wir können also wirklich $\binom{c_k}{d_k}$ mit $\frac{c_k}{d_k}$ identifizieren.
\end{falko}

\begin{falko} \label{F2.14}
	Jede rationale Zahl ist durch einen endlichen natürlichen Kettenbruch darstellbar.
\end{falko}

\begin{falko} \label{F2.15}
	Jede irrationale Zahl $\alpha$\marginnote{04.11.\\ \ [7]} ist auf genau eine Weise als natürlicher Kettenbruch darstellbar, und dieser Kettenbruch ist notwendigermaßen unendlich.
\end{falko}

\minisec{Bemerkung}
	Ist $\alpha \in \QQ$, so hat $\alpha$ eine Darstellung als endlicher Kettenbruch
	\begin{equation}
		\alpha = [q_0;q_1, \dots, q_n], \label{eq_2.34}
	\end{equation}
	der -- falls $n \geq 1$ ist -- mit einem $q_n \geq 2$ endet.

\begin{defn}[normierter Kettenbruch]
	Ein (natürlicher) Kettenbruch, der nicht mit 1 endet, falls er nicht von der Form $[q_0]$ ist, heißt ein \Index{normierter Kettenbruch}. Unendliche Kettenbrüche sind alle normiert.
\end{defn}

\minisec{Bemerkung}
	Sind $[q_0;q_1,\dots,q_n]$ und $[q_0';q_1',\dots,q_m']$ mit $n \geq m$ beide normiert vom selben Wert $\alpha$, so folgt $m = n$ und $q_i = q_i'$ für alle $i$.
	
\begin{satz} \label{satz_2.2}
	\begin{enumerate}[(i)]
		\item Ordnet man jeder reellen Zahl ihre Kettenbruchentwicklung zu, so erhält man eine Bijektion zwischen $\RR$ und der Menge aller normierten Kettenbrüche:
		\begin{equation}
			\RR \ni \alpha \longleftrightarrow [q_0;q_1,\dots] \label{eq_2.35}
		\end{equation}
		Die Umkehrabbildung ordnet jedem normierten Kettenbruch dessen Wert zu:
		\[ \alpha = [q_0;q_1,\dots] \]
		\item $\alpha$ rational $\Leftrightarrow$ Kettenbruchentwicklung von $\alpha$ ist endlich.
		\item Für die Näherungsbrüche $\frac{c_k}{d_k}$ des zu $\alpha$ gehörigen Kettenbruchs gilt\footnote{nur sinnvoll, falls es überhaupt noch ein $d_{k+1}$ gibt.}
		\begin{equation}
			\frac{1}{d_k(d_k+d_{k+1}} < \abs*{\alpha - \frac{c_k}{d_k}} \stackrel{(*)}{\leq} \frac{1}{d_k d_{k+1}} \qquad (k \geq 0), \label{eq_2.36}
		\end{equation}
		anders geschrieben:
		\begin{equation}
			\frac{1}{d_k + d_{k+1}} < \abs{d_k-\alpha - c_k} \stackrel{(*)}{\leq} \frac{1}{d_{k+1}} \qquad (k \geq 0) \label{eq_2.37}
		\end{equation}
	\end{enumerate}
\end{satz}

\minisec{Zusatz}
	Die Ungleichungen $(*)$ gelten mit $<$ bis auf den Fall $\alpha = [q_0;q_1,\dots,q_n]$ und $k = n-1$.

\minisec{Bemerkung}
	Aus (iii) folgt:
	\begin{equation}
		\abs*{ \alpha - \frac{c_{k-1}}{d_{k-1}}} < \alpha - \frac{c_k}{d_k} \label{eq_2.44}
	\end{equation}
	
\begin{falko} \label{F2.16}
	Für die Folge der Fibonacci-Zahlen $(u_n)_{n \in \ZZ}$\marginnote{07.11.\\ \ [8]} mit $u_0 = 0, u_1 = 1$ und $u_{n+1} = u_n + u_{n-1}$ gilt: \begin{enumerate}[(i)]
		\item $\frac{u_{n+2}}{u_{n+1}}$ ist der $n$-te Näherungsbruch von $\alpha = [1;1,1,\dots]$, $n \geq 2$.
		\item $\alpha = \frac{1}{2} + \frac{1}{2} \sqrt{5}$
		\item $u_n = \frac{\alpha^n-\beta^n}{\sqrt{5}}$ mit $\beta = \frac{1}{2} - \frac{1}{2} \sqrt{5}$, $n \in \ZZ$.
		\item $u_{m+n} = u_m u_{n+1} + u_{m-1} u_n$ mit $m,n \in \ZZ$, sowie für $m = n+1$: $u_{2m-1} = u_m^2 + u_{m-1}^2$.
		\item $u_{n+1} u_{n-1} - u_n^2 = (-1)^n$, bzw. $u_n^2 = u_{n-1} u_{n+1} + (-1)^{n+1}$, etc.
	\end{enumerate}
\end{falko}

\minisec{Bemerkung}
	$u_n$ ist die zu $\frac{\alpha^n}{\sqrt{5}}$ nächstgelegene ganze Zahl für $n \geq 0$.
	
\begin{falko} \label{F2.17}
	Für $a,b \in \ZZ$ gilt $(u_a,u_b) = u_{(a,b)}$, insbesondere $a | b \Rightarrow u_a | u_b$.
\end{falko}

\minisec{Bemerkung}
	$u_n$ Primzahl $\xLongrightarrow{\text{F\ref{F2.17}}} n$ Primzahl $\geq 3$, mit Ausnahme von $u_4 = 3$. Die Umkehrung gilt nicht.
\newpage
\section{Kongruenzrechnung}
\label{sec:para3}
	Zur Motivation:
	
\begin{satz}[Fermats kleiner Satz] \label{satz_3.1}
	Sei $p$ Primzahl. Für jede ganze Zahl mit $p \not | \ a$ gilt dann:
	\[ p | a^{p-1} - 1, \]
	d.h. $a^{p-1}$ lässt bei der Division durch $p$ stets den Rest 1. \index{Fermats kleiner Satz}
\end{satz}

\begin{defn}[Kongruenz in $\ZZ$] \label{def_3.1}
	Sei $m \in \NN$ fest. Für $x\in \ZZ$ sei $r_m(x)$ der eindeutige nichtnegative Rest von $x$ bei der Division von $x$ durch $m$. \marginnote{11.11. \\ \ [9]}
	\[ x = qm + r_m \quad 0 \leq r < m \]
	Wir definieren eine Relation
	\[ x \underset{m}{\sim} x' \qquad :\Leftrightarrow \qquad r_m(x) = r_m(x') \marginnote{Gleichheit bzgl. $m$}\]
	Statt $x \underset{m}{\sim} x'$ schreibt man nach Gauß:
	\[ x \equiv x' \modu m \]
	und sagt: $x$ ist kongruent zu $x'$ modulo $m$. \index{Kongruenz}
\end{defn}

\begin{falko} \label{F3.1} 
	\[ x \equiv x' \modu m \quad \Leftrightarrow \quad m | x-x' \]	
\end{falko}

\setcounter{countdef}{0}
\begin{defn}[Kongruenz allgemeiner] \label{def_3.1a}
	Sei $R$ ein kommutativer Ring und $m \in R$. Definiere:
	\[ x \equiv y \modu m :\Leftrightarrow m | x-y \]
	$x \equiv y \modu 0 \Rightarrow x=y$ \\
	$x \equiv y \modu 1$ gilt für alle $x,y \in R$. \\
	$x \equiv 0 \modu m \Leftrightarrow m | x$
\end{defn}

\begin{falko} \label{F3.2}
	\begin{enumerate}[(i)]
		\item $x \equiv x \modu m, \quad x \equiv y \modu m \Rightarrow y \equiv x \modu m, \quad x \equiv y \modu m, y \equiv z \modu m \Rightarrow x \equiv z \modu m,$ \\
		d.h. $\cdot \equiv \cdot \modu m$ ist eine Äquivalenzrelation auf $R$. Diese ist verträglich mit Addition und Multiplikation:
		\item $x \equiv x' \modu m, y \equiv y' \modu m \quad \Rightarrow x+y \equiv x'+y' \modu m, xy \equiv x'y' \modu m$
		\item $x \equiv y \modu m, m'|m \quad \Rightarrow \quad x \equiv y \modu m'$
		\item Für $R = \ZZ$: \qquad $x \equiv y \modu m_i, 1 \leq i \leq r \quad \Leftrightarrow \quad x \equiv y \modu \kgV(m_1,\dots,m_r)$
		\item $x \equiv y \modu m \quad \Rightarrow \quad cx \kon cy \modu xm \quad \Rightarrow \quad cx \equiv cy \modu m$
		\item Für einen Integritätsring $R$ gilt: $x \equiv y \modu m$ und $l | x, l|m, l \neq 0 \quad \Rightarrow \quad l|y$ und $\frac{x}{l} \equiv \frac{y}{l} \modu \frac{m}{l}$
		\item Für $R = \ZZ$: Ist $\ggT(c,m) = d$ mit $d \neq 0$, so gilt $ac \kon bc \modu m \quad \Rightarrow \quad a \kon b \modu \frac{m}{d} $
		\item $m | ac-bc \quad \Rightarrow \quad m | c(a-b) \quad \Rightarrow \quad \frac{m}{d} | \frac{c}{d}(a-b) \quad \xLongrightarrow{\enbrace*{\frac{m}{d},\frac{c}{d}}=1} \quad \frac{m}{d} | a-b$
	\end{enumerate}
\end{falko}

\begin{defn}[Restklasse] \label{def_3.2}
	Ein $m \in R$ teilt $R$ in disjunkte Mengen ein, die den zugehörigen Äquivalenzklassen entsprechen. Diese heißen die \bet{Restklassen} modulo $m$. \index{Restklasse}
\end{defn}

\begin{falko} \label{F3.3}
	Sei $n \in \NN$ ungerade. Dann:
	\[ (n-1)! \kon 1^2 \cdot 2^2 \dots \enbrace*{\frac{n-1}{2}}^2 \cdot (-1)^{\frac{n-1}{2}} \modu n \]
\end{falko}

\begin{falko} \label{F3.4}
	$a,b,c \in \ZZ$, $d:= (a,b)$. Die Gleichung
	\begin{equation}
		aX + bY = c \label{eq_3.1}
	\end{equation}
	ist genau dann lösbar über $\ZZ$, wenn $d | c$. Sei $d \neq 0$. Ist $(x_0,y_0)$ eine Lösung von \eqref{eq_3.1}, so gehört zu jeder Lösung $(x,y)$ von \eqref{eq_3.1} genau ein $t \in \ZZ$ mit
	\begin{equation}
		x = x_0 + t \frac{b}{d} \qquad y = y_0 - t \frac{a}{d}, \label{eq_3.2}
	\end{equation}
	und jedes $(x,y)$ wie in \eqref{eq_3.2} ist eine Lösung von \eqref{eq_3.1}.
\end{falko}

\begin{falko} \label{F3.5}
	Die Kongruenz
	\begin{equation}
		aX \kon c \modu m \label{eq_3.1a}
	\end{equation}
	ist genau dann lösbar über $\ZZ$, wenn
	\begin{equation}
		(a,m) | c. \label{eq_3.2a}
	\end{equation}
	Sei $d := (a,m) \neq 0$, und es gelte \eqref{eq_3.2a}. Die Lösungsmenge von \eqref{eq_3.1a} ist dann eine Restklasse modulo $\frac{m}{d}$. Die Kongruenz \eqref{eq_3.1a} besitzt genau $d = (a,m)$ viele Lösungen modulo $m$. Insbesondere gilt: Ist $(a,m) = 1$, so ist \eqref{eq_3.1a} für jedes $c$ lösbar und die Lösungen sind modulo $m$ eindeutig. 
\end{falko}

\setcounter{countdef}{1}
\begin{defn}[Restklassen allgemein] \label{def_3.2a}
	Sei $R$ ein kommutativer Ring, $m \in R$. Die \Index{Restklasse} modulo $m$, in der $a \in R$ liegt, hat die Gestalt
	\[ \{x \in R : x \kon a \modu m\} = a+mR = \{a + ym : y \in R\}. \]
	Die Menge aller Restklassen modulo $m$ bezeichnen wir mit $R/mR$, aber auch $R/m$. Der für uns wichtigste Fall ist $R = \ZZ$ und $m \in \NN$.
\end{defn}

\minisec{Beispiel}
	Sei $m \in \NN$. Betrachte
	\[R = \ZZ_{(m)} := \penbrace*{ \frac{b}{a} : a,b \in \ZZ, (a,m) = 1} \subseteq \QQ \]
	Die Inklusionsabbildung $\ZZ \rightarrow \ZZ_{(m)}$ vermittelt einen Ringisomorphismus
	\begin{equation}
		\ZZ/m\ZZ \longrightarrow \ZZ_{(m)}/m\ZZ_{(m)}
	\end{equation}
	
\minisec{Bemerkung}
	Sei $(a,m) = 1$. Wohlverstanden darf man also sagen: \\
	Die Kongruenz $aX \kon c \modu m$ besitzt die Lösung $\frac{c}{a} \modu m$. Es gibt ein $x \in \ZZ$ mit $x \kon \frac{c}{a} \modu m$, und für dieses ist $ax \kon c \modu m$.
	
\minisec{Beispiel}
	Die Kongruenz $7X \kon 1 \modu 123$ ist "eindeutig" lösbar:
	\[ 7x \kon 1 \modu 123 \quad \Rightarrow \quad x \kon \frac{1}{7} = \frac{4}{28} \kon \frac{-119}{28} = \frac{-17}{4} \kon \frac{-140}{4} \kon -35 \modu 123 \]
	Das funktioniert nicht immer so gut, aber allgemein kann man folgendes sagen:
	
\minisec{Bemerkung}
	Zur Lösung der Kongruenz
	\begin{equation}
		aX \kon 1 \modu m \quad \text{mit } (a,m) = 1 \text{ und } a \in \NN \label{eq_3.bem}
	\end{equation}
	Betrachte $\alpha = \frac{m}{a} \in \QQ$ und führe die Kettenbruchentwicklung durch. Diese endet mit $\frac{m}{a} = \frac{c_n}{d_n}$. Dann gilt (vgl. Beispiel nach Satz \ref{satz_2.2}):
	\[ (-1)^n c_{n-1}a - (-1)^n d_{n-1} m = 1 \quad \Rightarrow \quad a(-1)^n c_{n-1} \kon 1 \modu m \]
	Somit ist
	\[ x = (-1)^n c_{n-1} \]
	eine Lösung von \eqref{eq_3.bem}.
	
\begin{falko} \label{F3.6}
	Sei $m \in \NN$. \marginnote{14.11. \\ \ [10]}
	\begin{enumerate}[(i)]
		\item $\ZZ/m\ZZ$ ist auf natürliche Weise ein kommutativer Ring mit Eins ($\neq 0$, falls $m > 1$). Die Restklassenprojektion
		\begin{equation}
		\begin{aligned}
			\ZZ &\longrightarrow \ZZ/m\ZZ \\
			a &\longmapsto \overline{a} = a+m\ZZ =: a \mod m
		\end{aligned}
		\end{equation}
		ist ein Ringhomomorphismus. Für $m > 1$ ist $\ZZ/m\ZZ \rightarrow \ZZ_{(m)}/m\ZZ_{(m)}$ ein Ringisomorphismus, sodass man $\ZZ_{(m)}/m\ZZ_{(m)}$ mit $\ZZ/m\ZZ$ identifizieren kann.
		\item $\ZZ/m\ZZ$ hat genau $m$ Elemente.
		\item Für beliebige $c \in \ZZ$ ist $c, c+1, \dots, c+(m-1)$ ein \Index{Vertretersystem} modulo $m$. \\
		$S \subseteq \ZZ$ heißt ein Vertretersystem modulo $m$ bzw. von $\ZZ/m\ZZ$, wenn gilt: Zu jedem $x \in \ZZ$ existiert genau ein $a \in S$ mit $x \kon a \modu m$. Anders ausgedrückt: Die Einschränkung der Restklassenabbildung $\ZZ \rightarrow \ZZ/m\ZZ$ auf $S$ ist eine Bijektion. Äquivalent dazu ist, dass die Einschränkung injektiv oder surjektiv ist und $|S|=m$.
		\item $\ZZ/m\ZZ$ Integritätsring \quad $\Leftrightarrow$ \quad $\ZZ/m\ZZ$ Körper \quad $\Leftrightarrow$ \quad $m$ Primzahl. \\
		Für eine Primzahl $p$ heißt $\ZZ/p\ZZ$ der \Index{Restklassenkörper} modulo $p$.
		\item $(a,m) = d, x \kon a \modu m \quad \Rightarrow \quad (x,m) = d$
		\item $\overline{a} \in (\ZZ/m\ZZ)^\times \quad \Leftrightarrow \quad (a,m) = 1$. \\
		Die Elemente $\overline{a}$ von $(\ZZ/m\ZZ)^\times$ heißen prime Restklassen modulo $m$.
	\end{enumerate}
\end{falko}

\begin{lemma}
	Sei $R$ ein endlicher kommutativer Ring mit Eins. Dann ist
	\[ R^\times = \penbrace{a \in R : a \text{ ist kein Nullteiler von } R} \]
\end{lemma}

\begin{falko}[Satz von Wilson] \label{F3.7}
	Für $n \in \NN$ gilt: \index{Satz von Wilson}
	\[ n \text{ Primzahl} \quad \Leftrightarrow \quad (n-1)! \kon -1 \modu n \]
\end{falko}

\begin{falko} \label{F3.8}
	Sei $p \neq 2$ Primzahl. Dann ist die Kongruenz
	\[ X^2 \kon -1 \modu p \]
	genau dann lösbar in $\ZZ$, wenn $p \kon 1 \modu 4$, d.h. $p = 1+4k$ für ein $k \in \NN$.
\end{falko}

\minisec{Bemerkungen}
	\begin{enumerate}[1)]
		\item \ref{F3.8} anders formuliert. Setze $\FF_p := \ZZ/p\ZZ$ Körper.
		\[ \sqrt{-1} \in \FF_p \quad \Leftrightarrow \quad p \kon 1 \modu 4 \text{ oder } p = 2 \]
		\item Sei $p \neq 2$ Primzahl. Dann:
		\[ \enbrace*{\frac{p-1}{2}}!^2 \kon \begin{cases}
			-1 \modu p & \text{für } p \kon 1 \modu 4 \\
			1 \modu p & \text{für } p \kon 3 \modu 4
		\end{cases} \]
		Für $p \kon 3 \modu 4$ gilt also $\enbrace*{\frac{p-1}{2}}!^2 \kon \pm 1 \modu p$. Mehr dazu in Abschnitt 6.
	\end{enumerate}
	
\begin{defn}[Eulersche $\varphi$-Funktion]
	Für jede natürliche Zahl $m$ definiere
	\[ \varphi(m) := \#(\ZZ/m\ZZ)^\times \marginnote{$\varphi(1)=1$} \]
	Nach F\ref{F3.6} gilt $\varphi(m) = \# \{a \in \{0, 1, 2, \dots, m-1\} : a \text{ teilerfremd zu } m \}$. Für eine Primzahl $p$ ist daher $\varphi(p) = p-1$. $\varphi$ heißt \Index{Eulersche $\varphi$-Funktion}.
\end{defn}

\setcounter{countsatz}{0}
\begin{satz}[Satz von Euler-Fermat] \label{satz_3.1a}
	Aus $(a,m) = 1$ folgt $a^{\varphi(m)} \kon 1 \modu m$. \index{Satz von Euler-Fermat}
\end{satz}

\begin{lemma}
	Sei $G$ eine abelsche Gruppe der Ordnung $n$. Dann gilt $x^n = 1$ für alle $x \in G$.
\end{lemma}

\subsection{Simultane Kongruenzen}
\begin{satz}[Chinesischer Restsatz] \label{satz_3.2}
	Ist $m = m_1 m_2 \cdot m_r$ mit paarweise teilerfremden natürlichen Zahlen $m_1,\dots,m_r > 1$, so ist die Abbildung \index{Chinesischer Restsatz}
	\begin{equation}
	\begin{aligned}
		\ZZ/m &\longrightarrow \ZZ/m_1 \times \ZZ/m_2 \times \cdots \times \ZZ/m_r \\
		a \modu m &\longmapsto (a \modu m_1, a \modu m_2, \dots, a \modu m_r)	\label{eq_satz_3.2}
	\end{aligned}
	\end{equation}
	
	ein Isomorphismus von Ringen. Ist insbesondere $m = p_1^{e_1} p_2^{e_2} \cdots p_r^{e_r}$ die Primfaktorzerlegung einer natürlichen Zahl $m > 1$, so gilt
	\[ \ZZ/m\ZZ \simeq \ZZ/p_1^{e_1} \times \cdots \times \ZZ/p_r^{e_r} \]
	mit kanonischer Isomorphie. Der Isomorphismus $\eqref{eq_satz_3.2}$ vermittelt einen Isomorphismus
	\[ (\ZZ/m)^\times \simeq (\ZZ/m_1)^\times \times \dots \times (\ZZ/m_r)^\times \]
	der primen Restklassengruppen; insbesondere gilt
	\[ \varphi(m) = \varphi(m_1) \cdot \varphi(m_2) \cdots \varphi(m_r) \]
\end{satz}

\setcounter{countsatz}{1}
\begin{satz}[Chinesischer Restsatz für simultane Kongruenzen] \label{satz_3.2a}
	Sei $m = m_1m_2 \cdots m_r$ mit paarweise teilerfremden natürlichen Zahlen $m_1,\dots,m_r > 1$. Sind dann $a_1,\dots,a_r$ beliebige ganze Zahlen, so gibt es eine ganze Zahl $x$ mit \index{Chinesischer Restsatz} \index{simultane Kongruenz}
	\begin{equation}
	\begin{aligned}
		x &\kon a_1 \modu m_1 \\
		x &\kon a_2 \modu m_2 \\
		&\ \ \vdots \\
		x &\kon a_r \modu m_r \label{eq_satz_3.2a}
	\end{aligned}
	\end{equation}
	Durch \eqref{eq_satz_3.2a} ist $x$ modulo $m$ eindeutig bestimmt, ferner gilt:
	\[ x \text{ prim zu } m \quad \Leftrightarrow \quad a_i \text{ prim zu } m_i \text{ für alle } i \]
\end{satz}

\minisec{Bemerkung}
	Es genügt, sich $x_i$ zu verschaffen mit
	\begin{equation}
		x_1q_1 + x_2q_2 + \dots + x_rq_r \kon 1 \mod m \marginnote{$q_i = \frac{m}{m_i}$} \label{eq_satz_3.2_bem}
	\end{equation}
	Dann wird \eqref{eq_satz_3.2a} erfüllt von
	\[x = a_1(x_1q_1) + \dots + a_r(x_rq_r) \]
	Für jedes $1 \leq i \leq r$ bestimme (notfalls mit Kettenbruchentwicklung) ein $x_i \in \ZZ$ mit
	\[ q_i x_i \kon 1 \modu m_i \marginnote{$(q_i,m_i) =1$} \]
	Dann ist 
	\[ x_1q_1 + \dots + x_rq_r \kon 1 \modu m_i\]
	für alle $1 \leq i \leq r$, und es folgt \eqref{eq_satz_3.2_bem}.
	
\minisec{Korollar}
	Sei $f \in \ZZ[X]$, $m = m_1\cdot m_2 \cdot \dots \cdot m_r$ mit paarweise teilerfremden $m_i > 1$. Dann: \marginnote{18.11. \\ \ [11]}
	\[f(X) \kon 0 \modu m \text{ lösbar in } \ZZ \quad \Leftrightarrow \quad f(X) \kon 0 \modu m_i \text{ lösbar in } \ZZ \text{ für jedes } 1 \leq i \leq r \]
	Die natürliche Abbildung $\ZZ/m \rightarrow \prod_{i=1}^{r} \ZZ/m_i$ vermittelt eine Bijektion
	\[ \{\alpha \in \ZZ/m : f(\alpha) = 0\} \rightarrow \prod_{i=1}^{r} \{ \alpha_i \in \ZZ/m_i : f(\alpha_i) = 0 \} \]
	Für die Lösungsanzahlen $N_f(n) := \# \{\alpha \in \ZZ/n : f(\alpha) = 0 \}$ gilt also:
	\[N_f(m_1m_2\dots m_r) = N_f(m_1) N_f(m_2) \cdots N_f(m_r) \]
\newpage
\section{Die prime Restklassengruppe $\modu m$}
\label{sec:para4}

\begin{defn}[prime Restklassengruppe]
	Sei $m \in \NN$, $m > 1$. Dann heißt $(\ZZ/m\ZZ)^\times$ die \Index{prime Restklassengruppe} $\modu m$. Wir wissen: \begin{enumerate}[(1)]
		\item $\overline{a} \in (\ZZ/m\ZZ)^\times \Leftrightarrow (a,m)=1$
		\item $M := \{k \in \ZZ : 0 \leq k < m, (k,m) = 1 \}$ ist ein Vertretersystem von $(\ZZ/m\ZZ)^\times$. \\
		$(\ZZ/m\ZZ)^\times$ hat $\varphi(m) = \# M$ Elemente und ist eine abelsche Gruppe der Ordnung $\varphi(m)$.
		\item $\overline{a} = \alpha \in (\ZZ/m\ZZ)^\times \Rightarrow \alpha^{\varphi(m)} = 1 \Leftrightarrow a^{\varphi(m)} \kon 1 \modu m$ (Satz von Euler-Fermat)
	\end{enumerate}
\end{defn}

\begin{defn}[Primitivwurzel]
	Ein $\omega \in (\ZZ/m\ZZ)^\times$ heißt eine \Index{Primitivwurzel} von $(\ZZ/m\ZZ)^\times$, wenn sich jedes Element $\alpha \in (\ZZ/m\ZZ)^\times$ in der Form
	\[ \alpha = \omega^i \text{ für ein } i \in \NN_0 \]
	schreiben lässt; jedes $g \in \ZZ$ mit $\omega = \overline{g} = g \modu m$ heißt dann eine Primitivwurzel $\modu m$.
\end{defn}

\begin{satz}[Satz von Gauß] \label{satz_4.1}
	Ist $p$ eine Primzahl, so besitzt $(\ZZ/p\ZZ)^\times$ eine Primitivwurzel. Es gibt also ein $\omega \in (\ZZ/p\ZZ)^\times$, sodass sich jedes $\alpha \in (\ZZ/p\ZZ)^\times$ darstellen lässt in der Form
	\begin{equation}
		\alpha = \omega^i \text{ mit } 0 \leq i < p-1 \label{eq_4.1}
	\end{equation}
	Die Darstellung \eqref{eq_4.1} ist unter der Bedingung $0 \leq i < p-1$ eindeutig; $i = i(\alpha) = i_\omega(\alpha)$ heißt der \Index{Index} von $\alpha$ bezüglich $\omega$.
\end{satz}

Wählt man ein $g \in \ZZ$ mit $\omega = g \modu m$, so gilt also: Zu jedem $a \in \ZZ$ mit $p \not | a$ gibt es genau ein $i \in \ZZ$ mit
\[ a \kon gî \modu p, 0 \leq i < p-1 \]
$i = i(a) = i_g(a)$ heißt der Index von $a$ bzgl. $g$.

\minisec{Zusatz}
	Es gibt genau $\varphi(p-1)$ verschiedene Primitivwurzeln von $(\ZZ/p\ZZ)^\times$.

\subsection{Gruppentheoretische Vorbereitungen}
\begin{defn}[Ordnung eines Gruppenelements] \label{def_ordnung}
	Sei $G$ eine (abelsche) Gruppe der Ordnung $n$, d.h. $\#G = n$. Sei $\alpha \in G$. Wir wissen: $\alpha^n = 1$. Unter allen $m \in \NN$ mit $\alpha^m = 1$ sei nun $k$ das kleinste. Setze dann
	\[ \ord(\alpha) := k, \]
	die \Index{Ordnung} von $\alpha$. $\sprod{\alpha} := \penbrace{\alpha^j : j \in \ZZ}$ ist offenbar eine Untergruppe von $G$.
\end{defn}

\begin{lemma} \label{lemma_4.1}
	In der Situation von Definition \ref{def_ordnung} gelten: \begin{enumerate}[(1)]
		\item $\sprod{\alpha} = \{1, \alpha, \alpha^2, \dots, \alpha^{k-1}\}$, insbesondere $\ord(\alpha) = \ord(\sprod{\alpha})$.
		\item $\alpha^m = 1$ für $m \in \ZZ \Rightarrow \ord(\alpha) | m$
		\item Sei $\ord(\alpha) = k$ wie oben, dann vermittelt der Gruppenhomomorphismus
		\begin{equation}
		\begin{aligned}
			\ZZ &\longrightarrow \sprod{\alpha} \\
			j &\longmapsto \alpha^j
		\end{aligned}
		\end{equation}
		einen Gruppenisomorphismus $\ZZ/k\ZZ \rightarrow \sprod{\alpha}$, also $\sprod{\alpha} \simeq \ZZ/k\ZZ$.
		\item $G$ zyklisch $\Leftrightarrow \exists \alpha \in G$ mit $\ord(\alpha) = \ord(G)$. \\
		Eine Gruppe $G$ heißt \Index{zyklisch}, wenn es ein $\alpha \in G$ gibt mit $G = \sprod{\alpha}$. $\alpha$ heißt dann ein \Index{Erzeuger} von $G$.
	\end{enumerate}
\end{lemma}

\minisec{Bemerkung}
	Definitionsgemäß gilt:
	\[ (\ZZ/m\ZZ)^\times \text{ besitzt Primitivwurzel} \quad \Leftrightarrow \quad (Z/m\ZZ)^\times \text{ ist zyklisch}, \]
	und nach dem zuvor Gesagten:
	\[ \omega \text{ ist Primitivwurzel von } (\ZZ/m\ZZ)^\times \quad \Leftrightarrow \quad \ord(\omega) = \varphi(m)\]
	
\begin{defn}[Gruppenexponent]
	Sei $G$ eine endliche Gruppe. Das kgV aller $\ord(\alpha), \alpha \in G$ heißt der \Index{Exponent} $e = e(G)$ der Gruppe $G$.
\end{defn}

\minisec{Bemerkung}
	Ist $n = \ord(G), e=e(G)$, so gilt stets $e|n$, denn für jedes $\alpha \in G$ gilt $\alpha^n=1 \Rightarrow \ord(\alpha) | n \Rightarrow e|n$.
	
\begin{falko} \label{F4.1}
	Sei $G$ eine endliche abelsche Gruppe und sei $e$ ihr Exponent. Dann gibt es ein Element $\omega \in G$ mit $\ord(\omega) = e$.
\end{falko}

\setcounter{countsatz}{0}
\begin{satz}
	Sei $K$ ein Körper und $G$ eine endliche Untergruppe von $K^\times$. Dann ist $G$ zyklisch.
\end{satz}

\subsection{Restklassengruppen}
\begin{defn}[Restklassen, Restklassenabbildung]
	Sei $G$ eine Gruppe und $H \subseteq G$ eine Untergruppe. \marginnote{21.11. \\ \ [12]} Für $x,y \in G$ definiere eine Relation:
	\[ x \overset{H}{\sim} y :\Leftrightarrow yx^{-1} \in H (\Leftrightarrow y \in Hx), \]
	oder für eine abelsche Gruppe mit $+$ statt $\cdot$ als Verknüpfungssymbol:
	\[ x \overset{H}{\sim} y :\Leftrightarrow y-x \in H (\Leftrightarrow y \in H+x), \]
	$\overset{H}{\sim}$ ist eine Äquivalenzrelation. Mit $G/H$ bezeichnen wir die Menge der zugehörigen Äquivalenzklassen (\bet{Restklassen}). Die Abbildung \index{Restklasse}
	\begin{equation}
	\begin{aligned}
		G &\longrightarrow G/H \\
		x &\longmapsto \overline{x} := Hx
	\end{aligned}
	\end{equation}
	heißt \Index{Restklassenabbildung}.
\end{defn}

\minisec{Bemerkung}
	Die Relation $\overset{H}{\sim}$ ist verträglich mit der Multiplikation, falls $G$ abelsch ist. In diesem Fall ist $G/H$ eine Gruppe und die Restklassenabbildung ein Homomorphismus. \\
	Ist $G$ eine beliebige Gruppe, so gilt gleiches für $G/H$ genau dann, wenn für jedes $x \in G$ gilt: $Hx = xH$.
	
\begin{falko} \label{F4.2}
	Sei $G$ eine abelsche Gruppe der Ordnung $n$ und $n = p_1^{\nu_1} p_2^{\nu_2} \dots p_r^{\nu_r}$ die Primfaktorzerlegung von $n$. Für $1 \leq i \leq r$ sei
	\[ G_{p_i} := \{\alpha \in G : \alpha^{p_i^{\nu_i}} = 1 \} \leq G \]
	Dann ist die Abbildung
	\begin{equation}
	\begin{aligned}
		f\colon \prod_{i=1}^{r} G_{p_i} &\longrightarrow G \\
		(\alpha_1,\dots,\alpha_r) &\longmapsto \alpha_1 \alpha_2 \cdots \alpha_r \label{eq_F4.2}
	\end{aligned}
	\end{equation}
	ein Isomorphismus von Gruppen. Ferner gilt $\# G_{p_i} = p_i^{\nu_i}$ für $1 \leq i \leq r$.
\end{falko}

\minisec{Bemerkung 1}
	$G$ endliche Gruppe, $\alpha \in G$, $j \in \ZZ$. Dann gilt:
	\[ \ord(\alpha^j) = \frac{\ord(\alpha)}{(\ord(\alpha),j)} \]
	
\minisec{Bemerkung 2}
	Eine zyklische Gruppe der Ordnung $n$ hat genau $\varphi(n)$ Elemente der Ordnung $n$, also $\varphi(n)$ Erzeuger. \\
	
	
Wir werden jetzt die Struktur der primen Restklassengruppe modulo $p^\nu$ \marginnote{25.11. \\ \ [13]}
\[ G = G_\nu = (\ZZ/p^\nu\ZZ)^\times, p \text{ Primzahl}, \nu \in \NN, \nu > 1 \]
untersuchen. Es ist
\[ \ord(G) = \varphi(p^\nu) = \# \{0 \leq a < p^\nu : p \not | \ a \} = p^\nu - \# \{0 \leq a < p^\nu : p | a \} = p^\nu - p^{\nu-1} = (p-1) p^{\nu-1} \]
Damit gilt für jedes $n \in \NN$:
\[ \varphi(n) = \varphi \enbrace*{ \prod_{p|n} p^{w_p(n)}} = \prod_{p | n} \varphi \enbrace*{p^{w_p(n)}} = \prod_{p | n} \enbrace*{p^{w_p(n)}-p^{w_p(n)-1}} = \prod_{p | n} p^{w_p(n)} \enbrace*{1- \frac{1}{p}} = n \prod_{p | n} \enbrace*{1- \frac{1}{p}}\]

\begin{falko} \label{F4.3}
	Für jedes $n \in \NN$ gilt:
	\[ \varphi(n) = \prod_{p|n} \enbrace*{p^{w_p(n)} - p^{w_p(n)-1}} = n \prod_{p | n} \enbrace*{1 - \frac{1}{p}} \]
\end{falko}

\begin{defn}[1-Einheit und 1-Einheitengruppe]
	Sei $G_\nu = (\ZZ/p^\nu\ZZ)^\times$. Der Kern $G_\nu^{(1)}$ des Homomorphismus \index{1-Einheit}
	\begin{equation}
	\begin{aligned}
		(\ZZ/p^\nu\ZZ)^\times &\longrightarrow (\ZZ/p\ZZ)^\times \\
		a \modu p^\nu &\longmapsto a \modu p
	\end{aligned}
	\end{equation}
	heißt die \bet{Gruppe der 1-Einheiten} von $(\ZZ/p^\nu \ZZ)^\times$. Sie besteht aus den Elementen $a \modu p^\nu$ von $(\ZZ/p^\nu\ZZ)^\times$ mit $a \kon 1 \modu p$. Es ist $\ord(G_\nu^{(1)}) = p^{\nu-1}$.
\end{defn}

\newpage
\begin{lemma} \label{lemma_4.2}
	Sei $p$ Primzahl, $j \in \NN$, $a \in \ZZ$. Es gelte
	\begin{equation}
		a \kon 1 \modu p^j, \text{ aber } a \not\kon 1 \modu p^{j+1} \label{eq_lemma_4.2_1}
	\end{equation}
	Dann folgt -- außer für $p=2$ und $j = 1$:
	\begin{equation}
		a^p \kon 1 \modu p^{j+1}, \text{ aber } a^p \not\kon 1 \modu p^{j+2} \label{eq_lemma_4.2_2}
	\end{equation}
\end{lemma}

\begin{falko} \label{F4.4}
	Sei $\nu > 1$. \begin{enumerate}[(i)]
		\item Im Fall $p \neq 2$ ist für jedes $a$ der Gestalt $a = 1+cp$ mit $p \nmid c$ die Restklasse $a \modu p^\nu$ ein Element der Ordnung $p^{\nu - 1}$ in der 1-Einheitengruppe von $(\ZZ/p^\nu\ZZ)^\times$. Insbesondere gilt dies für $a = 1+p$. \\
		Die 1-Einheitengruppe von $(\ZZ/p^\nu\ZZ)^\times$ ist also für $p \neq 2$ zyklisch mit kanonischem Erzeuger $1+p \modu p^\nu$.
		\item Im Falle $p=2$ gilt: Für $\nu \geq 3$ ist $5 \modu 2^\nu$ ein Element der Ordnung $2^{\nu-2}$ in $(\ZZ/2^\nu\ZZ)^\times$. \\
		Für $\nu=2$: $(\ZZ/4\ZZ)^\times$ ist zyklisch mit $-1 \modu 4$ als Erzeuger.
	\end{enumerate}
\end{falko}

\begin{satz} \label{satz_4.2}
	Sei $p \neq 2$. Auch für $\nu \geq 2$ ist dann $(\ZZ/p^\nu\ZZ)^\times$ zyklisch. Mit anderen Worten: $(\ZZ/p^\nu\ZZ)^\times$ besitzt eine Primitivwurzel. Es existiert also ein $g \in \ZZ$, sodass es zu jedem $a \in \ZZ$ mit $(a,p) = 1$ genau ein $i \in \ZZ$ gibt mit
	\[ a \kon g^i \modu p^\nu \quad \text{und} \quad 0 \leq i < \varphi(p^\nu) \]
	Es gibt genau $\varphi(\varphi(p^\nu)) = \varphi((p-1)p^{\nu-1}) = \varphi(p-1) \varphi(p^{\nu-1})$ Primitivwurzeln von $(\ZZ/p^\nu)^\times$.
\end{satz}

\minisec{Zusatz}
	Ist schon eine Primitivwurzel $g_0 \modu p$ bekannt, so findet man eine Primitivwurzel $\modu p^\nu$ wie folgt: Ist $g_0^{p-1} \not\kon 1 \modu p^2$, so ist $g = g_0$ eine Primitivwurzel $\modu p^\nu$. Ist $g_0^{p-1} \kon 1 \modu p^2$, so ist $g = g_0 + p$ eine Primitivwurzel $\modu p^\nu$.
	
\minisec{Bemerkung}
	Folgende Aussagen sind für $p \neq 2$ und $g \in \ZZ$ äquivalent: \begin{enumerate}[(i)]
		\item $g$ ist Primitivwurzel $\modu p$ und $g^{p-1} \not \kon 1 \modu p^2$.
		\item $g$ ist Primitivwurzel $\modu p^n$ für alle $n \in \NN$.
		\item $g$ ist Primitivwurzel $\modu p^2$.
	\end{enumerate}
	
\begin{satz} \label{satz_4.3}
	Sei $\nu \in \NN, \nu \geq 3$. Zu jeder ungeraden Zahl $a \in \ZZ$ gibt es eindeutig bestimmte $k \in \{0,1\}$ und \linebreak $j \in \{0,1,\dots,2^{\nu-2}-1\}$ mit
	\[ a \kon (-1)^k 5^j \modu 2^\nu \]
	Mit anderen Worten: Die Abbildung
	\begin{equation}
	\begin{aligned}
		\ZZ/2\ZZ \times \ZZ/2^{\nu-2}\ZZ &\longrightarrow (\ZZ/2^\nu\ZZ)^\times \\
		(k \modu 2, j \modu 2^{\nu-2}) &\longmapsto (-1 \modu 2^\nu)^k \cdot (5 \modu 2^\nu)^j \label{eq_satz_4.3}
	\end{aligned}
	\end{equation}
	ist ein Isomorphismus von Gruppen. Es ist also
	\[ (\ZZ/2^\nu\ZZ)^\times = \sprod{-1 \modu 2^\nu} \times \sprod{5 \modu 2^\nu} \simeq \ZZ/2\ZZ \times \ZZ/2^{\nu-2}\ZZ \]
	Insbesondere ist $(\ZZ/2^\nu\ZZ)^\times$ nicht zyklisch.
\end{satz}

\setcounter{countsatz}{2}
\begin{satz} \label{satz_4.3a}
	Sei $p$ Primzahl mit $p \neq 2$, $\nu \in \NN$, $\nu \geq 2$. Dann existiert eine Primitivwurzel $g \modu p$, sodass die Abbildung
	\begin{equation}
	\begin{aligned}
		\ZZ/(p-1)\ZZ \times \ZZ/p^{\nu-1}\ZZ &\longrightarrow (\ZZ/p^\nu\ZZ)^\times \\
		(i \modu (p-1), j \modu p^{\nu-1}) &\longmapsto g^i(1+p)^j \modu p^\nu \label{eq_satz_4.3a}
	\end{aligned}
	\end{equation}
	wohldefiniert und ein Isomorphismus von Gruppen ist. Insbesondere ist $(\ZZ/p^\nu\ZZ)^\times$ zyklisch.
\end{satz}

\minisec{Bemerkung}
	Das direkte Produkt $G_1 \times G_2 \times \dots \times G_r$ \marginnote{28.11. \\ \ [14]} zyklischer Gruppen $G_i$ mit paarweise teilerfremden Ordnungen $m_i$ ist zyklisch von der Ordnung $m_1m_2\cdots m_r$. 
	
\begin{falko} \label{F4.5}
	Seien $G_1,G_2,\dots,G_r$ endliche abelsche Gruppen der Ordnungen $m_1,m_2,\dots,m_r$. Wenn $G:= G_1 \times \dots \times G_r$ zyklisch ist, so sind die $m_1,\dots,m_r$ paarweise teilerfremd und die $G_i$ sind zyklisch.
\end{falko}

\begin{falko} \label{F4.6}
	Sei $G$ eine zyklische Gruppe der Ordnung $n$. Dann ist jede Untergruppe $H$ von $G$ zyklisch mit $\ord(H) | n$. Die Abbildung $H \mapsto \ord(H)$ ist eine Bijektion zwischen der Menge aller Untergruppen $H$ von $G$ und der Menge aller natürlichen Teiler $d$ von $n$, und zwar ist $H_d := \{x \in G : x^d = 1\}$ \underline{die} Untergruppe der Ordnung $d$ von $G$. Es ist
	\[ H_{\frac{n}{d}} = \{x \in G : x^{\frac{n}{d}}=1 \} \overset{!}{=} \{y^d ; y \in G\} \]
	die Untergruppe der $d$-ten Potenzen in $G$.
\end{falko}

\minisec{Korollar}
	Für beliebige $n \in \NN$ gilt:
	\[ \sum_{d | n} \varphi(d) = n \]
	
\minisec{Bemerkung}
	Sei $G$ eine beliebige endliche Gruppe der Ordnung $n$. Für jedes $d | n, d \in \NN$ habe $G$ höchstens $d$ Elemente $x$ mit $x^d = 1$. Dann ist $G$ zyklisch.
	
\begin{falko} \label{F4.7}
	Sei $G$ eine beliebige endliche Gruppe der Ordnung $n$. Folgende Aussagen sind äquivalent:
	\begin{enumerate}[(i)]
		\item Für jedes $d \mid n$ gilt $\#\penbrace*{x \in G : x^d = 1} \leq d$.
		\item Für jedes $d \mid n$ hat $G$ höchstens eine Untergruppe der Ordnung $d$.
		\item $G$ ist zyklisch.
	\end{enumerate}
\end{falko}
	
\begin{satz} \label{satz_4.4}
	Sei $m \in \NN, m > 1$. Genau dann besitzt $(\ZZ/m\ZZ)^\times$ eine Primitivwurzel, wenn $m$ eine der Zahlen folgender Gestalt ist (mit einer Primzahl $p \neq 2$ und $\nu \geq 1$):
	\[ 2, \qquad 4, \qquad p^\nu, \qquad 2p^\nu \]
\end{satz}
\newpage
\section{Summen von zwei Quadraten in $\ZZ$ und der Gaußsche Zahlring $\ZZ[i]$}
\label{sec:para5}
	Ausgangspunkt ist F\ref{F3.8}:	\marginnote{02.12. \\ \ [15]}
	\[ p \kon 1 \modu 4 \quad \Rightarrow \quad \exists c \in \ZZ \text{ mit } c^2 \kon -1 \modu p, \]
	d.h. $c^2 + 1 = kp$ mit einem $k \in \ZZ$.
	
\begin{satz}[Fermat, Euler] \label{satz_5.1}
	Sei $p$ eine Primzahl. Ist $p \kon 1 \modu 4$, so gibt es $x,y \in \ZZ$ mit
	\begin{equation}
		p = x^2 + y^2 \label{eq_5.1}
	\end{equation}
	Ist umgekehrt $p$ in der Gestalt \eqref{eq_5.1} darstellbar, so ist $p \kon 1 \modu 4$ oder $p = 2$.
\end{satz}

\begin{defn}[Gaußscher Zahlring]
	\[ \ZZ[i] := \{a + bi : a,b \in \ZZ\} \]	
	heißt \Index{Gaußscher Zahlring}. Es ist $\ZZ[i]^\times = \{1, -1, i, -i\}$ und $(a+bi)(a-bi) = a^2+b^2$.
\end{defn}

\begin{satz}[Gaußscher Zahlring ist euklidisch] \label{satz_5.2}
	$\ZZ[i]$ ist ein euklidischer Ring mit euklidischer Normfunktion $\nu$ definiert durch \index{euklidischer Ring}
	\[ \nu(z) = z \cdot \overline{z} =: N(z), z \in \ZZ[i] \]
\end{satz}

\begin{falko} \label{F5.1}
	Sei $\pi$ ein Primelement $\neq 0$ von $\ZZ[i]$. Dann gibt es genau eine Primzahl $p$ mit $\pi | p$ in $\ZZ[i]$. Es gilt entweder $N(\pi)=p$ oder $N(\pi)=p^2$. Im ersten Fall nennen wir $\pi$ vom Grad 1, im zweiten Fall vom Grad 2.
\end{falko}

Um alle Primelemente $\pi$ von $\ZZ[i]$ zu finden, haben wir also die Primfaktorzerlegung aller $p \in \ZZ[i]$ zu untersuchen. Die $p$ heißen \bet{rationale Primzahlen}, die $\pi$ \bet{Gaußsche Primzahlen}. \index{rationale Primzahl} \index{Gaußsche Primzahl}

\begin{satz} \label{satz_5.3}
	Sei $p$ Primzahl sowie $\pi$ ein Primfaktor von $p$ in $\ZZ[i]$. Dann gibt es drei Fälle: \begin{enumerate}[(i)]
		\item $p \assoz \pi^2$ \quad ($p$ ist \bet{verzweigt} in $\ZZ[i]$)
		\item $p \assoz \pi$ \quad ($p$ ist \bet{träge} in $\ZZ[i]$, d.h. $p$ bleibt Primelement in $\ZZ[i]$)
		\item $p = \pi \overline{\pi}$ mit $pi \not \assoz \overline{\pi}$ \quad ($p$ \bet{zerfällt} in $\ZZ[i]$)
	\end{enumerate}
	Und zwar gilt: \index{Verzweigtheit} \index{Trägheit}
	\[\begin{array}{rclcl}
		\text{(i)} & \Leftrightarrow &  p = 2 &  & \\ 
		\text{(ii)} & \Leftrightarrow  & N(\pi) = p^2 & \Leftrightarrow &  p \kon 3 \modu 4 \\ 
		\text{(iii)} & \Leftrightarrow & N(\pi) = p & \Leftrightarrow &  p \kon 1 \modu 4 \\ 
	\end{array}\]
	Also ist z.B. $7$ auch in $\ZZ[i]$ ein Primelement, aber $5 = (2+i)(2-i)$ nicht.	
\end{satz}

\minisec{Korollar}
	Ist $p$ eine Primzahl mit $p \kon 1 \modu 4$, so ist $p$ in der Gestalt $p = a^2 + b^2$ mit $a,b \in \NN$ darstellbar. Bis auf Vertauschung von $a$ und $b$ ist diese Darstellung eindeutig. Ferner ist notwendigerweise $(a,b) = 1$.
	
\begin{satz} \label{satz_5.4}
	Sei $n \in \NN$. \begin{enumerate}[(i)]
		\item Genau dann ist $n$ eine Summe von zwei Quadraten in $\ZZ$, wenn für jede Primzahl $p \kon 3 \modu 4$ der Exponent $w_p(n)$ gerade ist.
		\item Besitzt $n$ eine primitive Darstellung als Summe von zwei Quadraten, d.h.
		\begin{equation}
			n = a^2 + b^2 \text{ mit teilerfremden } a,b \in \ZZ, \label{eq_satz_5.4_1}
		\end{equation}
		so folgt:
		\begin{equation}
			n \text{ hat keine Primteiler } p \kon 3 \modu 4, \text{ und es ist } 4 \not | \ n. \label{eq_satz_5.4_2}
		\end{equation}
		\item Umgekehrt: Gelte \eqref{eq_satz_5.4_2}, und bezeichne $s$ die Anzahl der ungeraden Primteiler von $n$. Für $n > 2$ hat dann $n$ genau $2^{s-1}$ primitive Darstellungen als Summe von zwei Quadraten, wenn nur wesentlich verschiedene Darstellungen gezählt werden. \\
		(Beachte: $n$ kann außerdem noch nicht-primitive Darstellungen haben, z.B. $50 = 7^2 + 1^2 = 5^2 + 5^2$.)
	\end{enumerate}
\end{satz}

\minisec{Korollar}
	Es sei $n$ \marginnote{05.12. \\ \ [16]} eine ungerade natürliche Zahl, $n > 1$. Besitzt $n$ im Wesentlichen nur eine einzige Darstellung als Summe von zwei Quadraten und ist diese Darstellung primitiv, so ist $n$ eine Primzahl (Umkehrung des Korollars von Satz \ref{satz_5.3}).
	
\minisec{Bemerkung}
	$45 = 6^2 + 3^2$ ist die einzige Darstellung von $45$ als Summe von zwei Quadraten, doch diese ist nicht primitiv. \\
	Im Übrigen ist die Voraussetzung, dass $n$ ungerade ist, wesentlich: Für $n = 10$ ist $10 = 3^2 + 1^2$ die im Wesentlichen einzige Darstellung von $10$ als Summe von zwei Quadraten und diese ist auch primitiv.
	
\subsection{Pythagoräische Tripel}
	\begin{equation}
		X^2 + Y^2 = Z^2 \label{eq_pyth_tripel}
	\end{equation}
	Eine Lösung $(a,b,c) \in \NN^3$ von \eqref{eq_pyth_tripel} heißt \Index{pythagoräisches Tripel}.\marginnote{16.12. \\ \ [18]} Wenn $\ggT(a,b,c) = 1$, heißt es \bet{primitiv}. Es genügt, primitive pythagoräische Tripel zu betrachten.
	
\begin{falko} \label{F5.2}
	Die Menge der primitiven pythagoräischen Tripel $(a,b,c)$, bei denen ohne Einschränkung $b$ gerade ist, wird geliefert durch:
	\begin{equation}
		\penbrace{ (u^2-v^2,2uv,u^2+v^2) \in \NN^3 : u,v \in \NN, \ggT(u,v) = 1, uv \text{ gerade, } u>v } \label{eq_pyth_tripel2}
	\end{equation}
\end{falko}

\begin{falko}
	Die Gleichung $X^4+Y^4 = Z^2$ besitzt keine Lösung $(a,b,c) \in \NN^3$. (In $\ZZ^3$ hat sie gewisse triviale Lösungen: $(0, \pm b, b^2), (\pm a, 0, a^2)$.)
\end{falko}

\minisec{Folgerung}
	Die Gleichung $X^4+Y^4=Z^4$ hat keine Lösung in $\NN^3$. (Fermat-Vermutung für den Exponenten $4$).
	
\subsection{Pythagoräische Quadrupel}
	Wir fragen nach allen Quadern mit ganzzaligen Kantenlängen $a,b,c > 0$, deren Raumdiagonale ebenfalls ganzzahlige Länge $d > 0$ hat. Wir suchen also nach allen Quadrupeln $(a,b,c,d) \in \NN^4$ mit
	\begin{equation}
		d^2 = a^2+b^2+c^2, \label{eq_pyth_quadr_1}
	\end{equation}
	und zwar ohne Einschränkungen nur nach primitiven, d.h. solchen mit $\ggT(a,b,c,d) = 1$, was gleichbedeutend mit $\ggT(a,b,c) = 1$ ist. Von den Zahlen $a,b,c$ kann höchstens eine ungerade sein, denn sonst ist $d^2 \kon 2 \modu 4$ oder $d^2 \kon 3 \modu 4$. Wegen $\ggT(a,b,c) = 1$ sind $a,b,c$ nicht alle gerade. Wir setzen dabei ohne Einschränkung voraus, dass $a$ ungerade ist, aber $b,c$ gerade sind. Nach \eqref{eq_pyth_quadr_1} ist $d$ ungerade, also sind $a+d$ und $a-d$ gerade. Daher schreibe \eqref{eq_pyth_quadr_1} in der Form:
	\begin{equation}
		\frac{d+a}{2} \frac{d-a}{2} = \frac{b^2+c^2}{4} = \enbrace*{\frac{b}{2}}^2 + \enbrace*{\frac{c}{2}}^2 \label{eq_pyth_quadr_2}
	\end{equation}
	Es liegt also nahe, nach einer Zerlegung
	\begin{equation}
		\frac{b}{2} + i \frac{c}{2} = (p+qi)(u+vi) \label{eq_pyth_quadr_3}
	\end{equation}
	in $\ZZ[i]$ zu suchen, die bei Normbildung in \eqref{eq_pyth_quadr_2} übergeht, die also
	\begin{equation}
		N(P+qi) = \frac{d+a}{2} \quad \text{und} \quad N(u+iv) = \frac{d-a}{2} \label{eq_pyth_quadr_4}
	\end{equation}
	erfüllt, d.h. $\frac{d+a}{2} = p^2 + q^2$ und $\frac{d-a}{2} = u^2+v^2$, mithin
	\begin{equation}
		d = p^2 + q^2 + u^2 + v^2 \quad a = p^2+q^2-u^2-v^2 \label{eq_pyth_quadr_5}.
	\end{equation}
	Wegen $(p+qi)(u+iv) = (pu-qv) + (qu+pv)i$ ist andererseits \eqref{eq_pyth_quadr_3} äquivalent mit
	\begin{equation}
		b = 2(pu-qv) \quad c = 2qu + pv. \label{eq_pyth_quadr_6}
	\end{equation}
	Um zu einer Zerlegung \eqref{eq_pyth_quadr_3}, die \eqref{eq_pyth_quadr_4} erfüllt, zu gelangen, gehen wir von der Primfaktorzerlegung von $\frac{b}{2} + i \frac{c}{2}$ in $\ZZ[i]$ aus. Diese notieren wir in der Gestalt
	\begin{equation}
		\frac{b}{2} + i \frac{c}{2} = \varepsilon n \pi_1^{\nu_1} \overline{\pi_1}^{\nu_1'} \cdots \pi_r^{\nu_r} \overline{\pi_r}^{\nu_r'} \pi_{r+1}^{\nu_{r+1}} \cdots \pi_s^{\nu_s}, \label{eq_pyth_quadr_7}
	\end{equation}
	wobei wir in $n$ alle Primfaktoren zusammenfassen, die zu Primzahlen $q \kon 3 \modu 4$ gehören, während zu jedem der Primelemente $\pi_i$ Primzahlen $p_i = N(\pi_i) = \pi_i \overline{\pi_i}$ mit $p_i \kon 1 \modu 4$ oder $p_i = 2$ gehören; im letzteren Fall ist $\overline{\pi_i} \assoz \pi_i \assoz 1+i$, während für $p_i \neq 2$ stets $\overline{\pi_i} \not\assoz \pi_i$ gilt. Die $\pi_i$ mit $1 \leq i \leq r$ bezeichnen genau die Primteiler $\pi_i$ mit $\overline{\pi_i} \not\assoz \pi_i$, für die auch $\overline{\pi_i}$ in $\frac{b}{2} + i \frac{c}{2}$ aufgeht. Per Normbildung geht \eqref{eq_pyth_quadr_7} über in
	\begin{equation}
		\enbrace*{\frac{b}{2}}^2 + \enbrace*{\frac{c}{2}}^2 = n^2 p_1^{\nu_1 + \nu_1'} \cdots p_r^{\nu_r + \nu_r'} p_{r+1}^{\nu_{r+1}} \cdots p_s^{\nu_s} \label{eq_pyth_quadr_8}
	\end{equation}

\begin{lemma}
	Jede Primzahl $q \kon 3 \modu 4$ geht in $d+a$ bzw. $d-a$ mit gerader Vielfachheit auf.
\end{lemma}

	Nach dem Lemma haben die Faktoren $\frac{d+a}{2}$ und $\frac{d-a}{2}$ in \eqref{eq_pyth_quadr_2} mit Blick auf \eqref{eq_pyth_quadr_8} die Darstellungen
	\begin{equation}
		\begin{array}{c}
			\frac{d+a}{2} = n_1^2 p_1^{\lambda_1} \cdots p_s^{\lambda_s} \quad \frac{d-a}{2} = n_2^2 p_1^{\mu_1} \cdots p_s^{\mu_s} \\
			\text{ mit } n_1,n_2 \in \NN, \lambda_i,\mu_i \geq 0 \text{ und } \lambda_i + \mu_i = \nu_i + \nu_i' \text{ für } 1 \leq i \leq r, \label{eq_pyth_quadr_9} \\
			\lambda_i + \mu_i = \nu_i \text{ für } r+1 \leq i \leq s; n_1 n_2 = n
		\end{array}
	\end{equation}
	Ausschlaggebend ist nun die folgende Feststellung: Für $1 \leq i \leq r$ geht $p_i$ nicht in beiden der Zahlen $\frac{d+a}{2}, \frac{d-a}{2}$ auf. Denn sonst wäre $p_i$ ein Teiler von $a$ und $d$, andererseits geht $p_i = \pi_i \overline{\pi_i}$ nach \eqref{eq_pyth_quadr_7} in $b$ und $c$ auf, im Widerspruch zu $\ggT(a,b,c)=1$. Unter den $p_i$ mit $1 \leq i \leq r$ bezeichnen nun ohne Einschränkung $p_1, \dots, p_m$ genau die $p_i$, die in $\frac{d+a}{2}$ aufgehen. Dann gilt in \eqref{eq_pyth_quadr_9} genauer:
	\begin{equation}
		\begin{array}{c}
			\frac{d+a}{2} = n_1^2 p_1^{\nu_1 + \nu_1'} \cdots p_m^{\nu_m + \nu_m'} p_{r+1}^{\lambda_{r+1}} \cdots p_s^{\lambda_s} \\
			\frac{d-a}{2} = n_2^2 p_{m+1}^{\nu_{m+1} + \nu_{m+1}'} \cdots p_r^{\nu_r + \nu_r'} p_{r+1}^{\mu_{r+1}} \cdots p_s^{\mu_s} \label{eq_pyth_quadr_10} \\
			\text{ mit } n_1 n_2 = n; \lambda_i, \mu_i \geq 0 \text{ und } \lambda_i + \mu_i = \nu_i \text{ für } r+1 \leq i \leq s
		\end{array}
	\end{equation}
	Setzen wir nun
	\begin{equation}
		\begin{array}{c}
			p+qi := \varepsilon n_1 \pi_1^{\nu_1} \overline{\pi_1}^{\nu_1'} \cdots \pi_m^{\nu_m} \overline{\pi_m}^{\nu_m'} \pi_{r+1}^{\lambda_{r+1}} \cdots \pi_s^{\nu_s} \\
			u+vi := n_2 \pi_{m+1}^{\nu_{m+1}} \overline{\pi_{m+1}}^{\nu_{m+1}'} \cdots \pi_r^{\nu_r} \overline{\pi_r}^{\nu_r'} \pi_{r+1}^{\mu_{r+1}} \cdots \pi_s^{\nu_s}
		\end{array}		
	\end{equation}
	(mit derselben Einheit $\varepsilon$ wie in \eqref{eq_pyth_quadr_7}), so sind in der Tat \eqref{eq_pyth_quadr_3} wie \eqref{eq_pyth_quadr_4} erfüllt, vgl. \eqref{eq_pyth_quadr_7}. Im übrigen erhält man auch folgende Eindeutigkeitsaussage:
	$p+qi, u+vi$ in \eqref{eq_pyth_quadr_3} lassen sich nur durch $\varepsilon_1(p+qi), \varepsilon_2(u+iv)$ mit Einheiten $\varepsilon_1, \varepsilon_2$ ersetzen, die $\varepsilon_1 \varepsilon_2 = 1$ erfüllen. Mit einem Parameter-Quadrupel $(p,q,u,v)$ sind so auch $(-p,-q,-u,-v), (-q,p,v,-u), (q,-p,-v,u)$ zulässige Quadrupel von Parametern, aber keine weiteren.
	
\minisec{Bemerkung 1}
	Für ein Parameter-Quadrupel zu $a,b,c,d$ muss offenbar gelten:
	\begin{equation}
		\begin{array}{c}
			p^2+q^2 > u^2+v^2 > 0 \quad pu - qv > 0 \quad qu+pv > 0 \quad \ggT(p,q,u,v)=1 \text{ und} \\
			\text{von den Parametern } p,q,u,v \text{ ist genau einer ungerade oder genau einer gerade.} \label{eq_pyth_quadr_12}
		\end{array}
	\end{equation}
	
\minisec{Bemerkung 2}
	Für jedes Quadrupel $(p,q,u,v) \in \ZZ^4$ mit \eqref{eq_pyth_quadr_12} liefern die Formeln \eqref{eq_pyth_quadr_5} und \eqref{eq_pyth_quadr_6} ein Quadrupel $(a,b,c,d) \in \NN^4$ mit $a^2+b^2+c^2 = d^2$. (Allerdings ist $(a,b,c,d)$ nicht notwendig primitiv, vgl. weiter unten.)
	
\minisec{Bemerkung 3}
	Für $(4,7,2,-1)$ und $(8,1,2,1)$ liefern \eqref{eq_pyth_quadr_5} und \eqref{eq_pyth_quadr_6} das gleiche Quadrupel $(a,b,c,d) = (60,30,20,70)$. Aber dieses ist nicht primitiv und man hat keinen Widerspruch zur obigen Eindeutigkeitsaussage. Außerdem verletzen diese Quadrupel die Paritätsbedingung in \eqref{eq_pyth_quadr_12}.
	
\minisec{Bemerkung 4}
	Die Quadrupel $(1,13,2,-1)$ und $(11,7,2,1)$ erfüllen jeweils alle Bedingungen in \eqref{eq_pyth_quadr_12}. Sie liefern das gleiche Quader-Quadrupel $(a,b,c,d) = (165, 30, 50, 175)$, aber dieses ist nicht primitiv. Kann es auch nicht sein, denn sonst widerspräche das der obigen Eindeutigkeitsaussage.
	
\minisec{Frage}
	Kann man \eqref{eq_pyth_quadr_12} so ergänzen, dass $(p,q,u,v)$ eine primitive Lösung $(a,b,c,d)$ liefert?
\newpage
\section{Schwache Lösungen}
\label{sec:para6}
Frage: Wie können auch nicht-glatte Funktionen $u$ als (verallgemeinerte) Lösungen aufgefasst werden? \marginnote{2. Mai} \\
Idee: Nutze partielle Integration, um Ableitungen loszuwerden.\\
Häufig kann die PDGL in \Index{Erhaltungsform} oder \Index{Divergenzform} geschrieben werden:
\begin{equation}
	\diver G(x,u(x)) = R(x,u(x)) \label{eq_0205_11}
\end{equation}
für Funktionen $G\colon \RR^n \times \RR \rightarrow \RR^n, R\colon \RR^n \times \RR \rightarrow \RR$. Wir wollen \eqref{eq_0205_11} auf einem Gebiet $\Omega$ lösen für Cauchy-Daten
\[u = g \text{ auf } \Gamma \subseteq \der \Omega \]

\begin{bsp} \label{bsp_15}
	$a(x)u_{x_1} + u u_{x_2} = c(x,u)$ ist äquivalent \marginnote{[15]} zu $(au)_{x_1} + \left( \frac{u^2}{2} \right)_{x_2} = c(x,u) + a_{x_1} u$ und kann in Erhaltungsform geschrieben werden für $G(x,u) = \begin{pmatrix} au \\ u^2/2 \end{pmatrix}, R(x,u) = c(x,u) + a_{x_1} u$. \\
	Wir multiplizieren \eqref{eq_0205_11} mit einer sogenannten \Index{Testfunktion}
	\[ \psi \colon \RR^n \rightarrow \RR, \psi \text{ ist glatt}, \psi = 0 \text{ auf } \der \Omega \setminus \Gamma \]
	und integrieren partiell:
	\begin{equation}
	\begin{aligned}
		&\int_\Omega \psi(x) (\diver G(x,u(x)) - R(x,u(x))) dx \\ \notag
		= &\int_{\der \Omega} \psi(x) G(x,u(x)) \cdot v(x) dx - \int_\Omega \nabla \psi(x) \cdot F(x,u(x)) + \psi(x) R(x,u(x)) dx \\
		= &\int_\Gamma \psi(x) G(x,g(x)) \cdot v(x) dx - \int_\Omega \nabla \psi(x) \cdot G(x,u(x)) + \psi(x) R(x,u(x)) dx
	\end{aligned}
	\end{equation}
	Die Gleichung macht auch für nicht-differenzierbares $u$ Sinn.
\end{bsp}

% % % % % % % % % % % % % 02. Mai	
\begin{defn}[schwache Lösung]
\label{def_schwache_lsg} \label{def_16}
	Eine messbare, beschränkte Funktion $u(x)$ auf $\Omega$ heißt schwache Lösung von \eqref{eq_0205_11}, wenn sie \marginnote{[16]}
	\begin{equation}
		0 = \int_\Gamma \psi(x) G(x,g(x)) \cdot v(x) dx - \int_\Gamma \nabla \psi(x) \cdot G(x,u(x)) + \psi (x) R(x,u(x)) dx \label{eq_0205_12}
	\end{equation}
	für alle Testfunktionen $\psi$ erfüllt.
\end{defn}

\minisec{Unstetige Lösungen und Schocks}
	Seien
	\begin{itemize}
		\item $G \subseteq \Omega$ stückweise glatte Hyperfläche, die $\Omega$ in $\Omega_1$ und $\Omega_2$ teilt.
		\item $u$ eine schwache Lösung von \eqref{eq_0205_11}, die auf $\Omega_1,\Omega_2$ stetig differenzierbar ist
		\item $\nu_i$ die äußere Einheitsnormale an $\der \Omega_i$,$i=1,2$
		\item $[f]_-^+$: Sprung einer Funktion $f$ über $C$ von $\Omega_1$ nach $\Omega_2$
	\end{itemize}
	Für alle Testfunktionen $\psi$ gilt:
	\begin{equation}
	\begin{aligned}
		0 &= \int_\Gamma \psi G(x,g(x)) \cdot \nu(x)dx - \int_{\Omega_1} \nabla \psi \cdot G(x,u(x)) + \psi R(x,u(x)) dx - \int_{\Omega_2} \nabla \psi \cdot G(x,u(x)) + \psi R(x,u(x)) dx \\ \notag
		&= \int_{\Gamma} \psi G(x,g(x)) \cdot \nu(x) dx - \int_{\der\Omega_1} \psi G(x,u(x)) \cdot \nu_1(x) dx + \int_{\Omega_1} \psi (\diver G(x,u(x)) - R(x,u(x))) dx \\
		&\quad - \int_{\der\Omega_2} \psi G(x,u(x)) \cdot \nu_2(x) dx + \int_{\Omega_2} \psi (\diver G(x,u(x)) - R(x,u(x))) dx \\
		&= \int_C \psi[G(x,u(x))]_-^+ \cdot \nu_1(x) dx,
	\end{aligned}
	\end{equation}
	da $\diver G(x,u(x)) - R(x,u(x)) = 0$ in $\Omega_1,\Omega_2$. Wir erhalten
	\begin{equation}
		[G(x,u(x))]_-^+ \cdot \nu(x) = 0 \label{eq_0205_13}
	\end{equation}
	für die Normale $\nu$ an $C$. Eine Unstetigkeit von $u$ bzw. $G$, ein sogenannter \Index{Schock}, kann also nur entlang einer Fläche $C$ normal zu $[G]_-^+$ auftreten. \\
	\mbox{} \\
	In 2D kann die Kurve $C$ lokal als Graph einer Funktion $x_2(x_1)$ (bzw. $x_1(x_2)$) geschrieben werden. Ein Tangentialvektor ist somit $\begin{pmatrix} 1 \\ dx_2/dx_1 \end{pmatrix}$, ein Normalvektor $\begin{pmatrix} dx_2/dx_1	\end{pmatrix}$, und \eqref{eq_0205_13} wird zur sogenannten \bet{Rankine-Hugoniot} \bet{-Bedingung} für einen Schock: \index{Rankine-Hugoniot-Bedingung}
	\begin{equation}
		\frac{dx_2}{dx_1} = \frac{[G_2]_-^+}{[G_1]_-^+} \label{eq_0205_14}
	\end{equation}
	
\begin{bem}[Rankine-Hugoniot-Bedingung für semilineare PDGL] \label{bem_17}
	Für semilineare Gleichungen \marginnote{[17]} $0 = b(x) \cdot \nabla u(x) + c(u(x),x)$ wird dies zu $\frac{dx_2}{dx_1} = \frac{b_2[u]_-^+}{b_1[u]_-^+} = \frac{b_1}{b_2}$, d.h. der Schock liegt entlang einer charakteristischen Kurve.
\end{bem}
	
\begin{bsp} \label{bsp_18}
	Betrachte das Cauchy-Problem \marginnote{[18]}
	\[ \begin{cases}
		u_t+uu_x = 0 \text{ auf } (0,\infty) \times \RR \\
		u(0,x) = g(x) := \begin{cases}
			1	& \text{ für } x \leq 0 \\
			1-x & \text{ für } 0 \leq x \leq 1 \\
			0	& \text{ für } 1 \leq x \end{cases} \end{cases} \]
	Entlang der charakteristischen Kurven $s \mapsto (s,y+g(y)s)$ für $y \in \RR$ ist $u = z(s) = g(y)$ konstant, somit
	\[ u(t,x) = \begin{cases}
		1	&	\text{ für } x \leq t \\
		\frac{1-x}{1-t} & \text{ für } t \leq x \leq 1 \\
		0	&	\text{ für } 1 \leq x \end{cases} \]
	solange $t < 1$. In $(t,x) = (1,1)$ schneiden sich die charakteristischen Kurven, und die Lösung mithilfe der Charpits Gleichungen bricht zusammen. An dieser Stelle setzen wir einen Schock entlang einer Kurve $x = X(t)$ ein und betrachten die Erhaltungsform $u_t + \left( \frac{u^2}{2} \right)_x = 0$. Sei $(0,y_\pm(t))$ der Anfangspunkt der charakteristischen Kurven, die $X(t)$ zur Zeit $t$ von links ($-$) und von rehts ($+$) schneiden. Es gilt somit
	\[ X(t) = y_-(t) + g(y_-(t))t = y_+(t) + g(y_+(t))t \]
	Weiterhin gilt die Rankine-Hugoniot-Bedingung
	\[ \frac{dX}{dt} = \frac{\frac{u_+^2}{2}-\frac{u_-^2}{2}}{u_+ - u_-} = \frac{u_+ + u_-}{2} = \frac{g(y_+(t)) + g(y_-(t))}{2} \]
	und $X(1)=1$. Eine Lösung dieser Gleichungen ist gegeben durch $X(t) = \frac{1+t}{2}, y_-(t) = \frac{1-t}{2}, y_+(t) = \frac{1+t}{2}$ (die Charakteristiken von links und rechts haben $u_-=1, u_+ = 0$), folglich ist
	\[ u(t,x) = \left. \begin{cases}
		1	&	\text{ für } x \leq t \\
		\frac{1-x}{1-t}	&	\text{ für } t \leq x \leq 1 \\
		0	&	\text{ für } 1 \leq x \end{cases} \right\} \text{ für } t<1 \text{ und } u(t,x) = \left. \begin{cases}
		1	&	\text{ für } x \leq \frac{1+t}{2} \\
		0	&	\text{ für } \frac{1+t}{2} \leq x \end{cases} \right\} \text{ für } t \geq 1 \]
	eine schwache Lösung des Problems.
\end{bsp}
	
\begin{bem} \label{bem_19}
	Schwache Lösungen sind im Allgemeinen nicht eindeutig. Betrachte z.B. \marginnote{[19]}
	\[ \begin{cases} 
		u_t + uu_x = 0 \text{ auf } (0,\infty) \times \RR \\
		u(0,x) = g(x) := \begin{cases}
			0 & \text{ für } x \leq 0 \\
			1 & \text{ sonst} \end{cases} \end{cases} \]
	mit den schwachen Lösungen:
	\begin{itemize}
		\item $u(t,x) = \begin{cases}
			0 & \text{ für } x \leq t/2 \\
			1 & \text{ sonst} \end{cases}$ \\
		(Schock zu Erhaltungsform $u_t + \left( \frac{u^2}{2} \right)_x = 0$; Rankine-Hugoniot: $\frac{dx}{dt} = \frac{\frac{u_+^2}{2}-\frac{u_-^2}{2}}{u_+-u_-} = \frac{u_++u_-}{2}$)
		\item $u(t,x) = \begin{cases}
			0 & \text{ für } x \leq 2t/3 \\
			1 & \text{ sonst} \end{cases}$ \\
		(Schock zu Erhaltungsform $\left( \frac{u^2}{2} \right)_t + \left( \frac{u^3}{2} \right)_x = 0$; Rankine-Hugoniot: $\frac{dx}{dt} = \frac{2}{3} \frac{u_+^3-u_-^3}{u_+^2-u_-^2} = \frac{2}{3} \frac{u_+^2+u_-u_++u_-^2}{u_++u_-}$)
		\item $u(t,x) = \begin{cases}
			0 & \text{ für } x\leq 0 \\
			x/t & \text{ für } 0 < x < t \\
			1 & \text{ sonst} \end{cases}$ \\
		(Verdünnungswelle)
	\end{itemize}
	Welche Erhaltungsform korrekt ist, hängt vom modellierten physikalischen Problem ab -- auch, ob ein Schock oder stattdessen eine Verdünnungswelle korrekt sind. Man kann für eine vorgegebene Erhaltungsform eine eindeutige Lösung selektieren durch Hinzunahme einer zusätzlichen Bedingung, zum Beispiel: \begin{enumerate}
		\item Entropiebedingung: Wir legen eine bestimmte Funnktion von $f(u)$ fest, die über den Schock zunehmen soll, d.h. $[f(u)]_-^+ \geq 0$. Die Wahl $f(u) = u$ verbietet obige Schocks.
		\item Viskosität: Wir betrachten $0 = F(\nabla u, u, x)$ als Grenzfall für $\varepsilon \rightarrow 0$ von $\varepsilon \Delta u = F(\nabla u,u,x)$, welches eine eindeutige Lösung besitzt.
		\item Kausalität: Ist eine Variable die Zeit, so fließt physikalisch Information entlang der charakteristischen Kurven in Richtung wachsender Zeit. Es soll Information in den Schock hinein und nicht aus ihm hinaus fließen (dies verbietet obige Schocks).
	\end{enumerate}
\end{bem}

% % % % % % % % % % % % 6. Mai	
\minisec{Viskositätslösungen}
	Betrachte die nichtlineare PDGL 1. Ordnung \marginnote{6. Mai}
	\begin{equation}
		H(x,u(x),\nabla u(x)) = 0 \text{ auf } \Omega \label{eq_0605_15}
	\end{equation}
	mit Randbedingungen auf $\der \Omega$, wobei $H\colon \Omega \times \RR \times \RR^n \rightarrow R$ stetig ist und konvex in $\nabla u$. \eqref{eq_0605_15} heißt auch \Index{Hamilton-Jacobi-Gleichung}. Lösungen sind typischerweise weder eindeutig, noch überall differenzierbar.
	
\begin{bsp}[Eikonalgleichung] \label{bsp_20}
	Die \Index{Eikonalgleichung} \marginnote{[20]}
		\[ 0 = H(x,u,p) = |p|-1 \]
	wird oft verwendet, um den Abstand $u(x) = \dist(x,\Gamma)$ einer Teilmenge $\Gamma \subset \RR^n$ zu berechnen. Allerdings existiert bereits für $n = 1$ und Randdaten $u = 0$ auf $\Gamma = \{0,1\}$ keine stetig differenzierbare Lösung $u$, aber viele Funktionen $u$, die \eqref{eq_0605_15} fast überall erfüllen. \\
	\mbox{} \\
	Idee, um Wohlgestelltheit zu erhalten: Füge eine \glqq Viskositätsterm\grqq \ $-\varepsilon \Delta u_\varepsilon$ hinzu und betrachte $\varepsilon \rightarrow 0$:
	\begin{equation}
		-\varepsilon \Delta u_\varepsilon + H(x,u_\varepsilon(x),\nabla u_\varepsilon (x)) = 0 \label{eq_0605_16}
	\end{equation}
\end{bsp}
	
\begin{lemma} \label{lemma_21}
	Sei $\Omega \subset \RR^n$ offen, \marginnote{[21]} $\overline{x} \in \Omega$, $v, v_k \in C^1(\Omega)$ mit $v_k \rightarrow v$ lokal gleichmäßig. Hat $v$ ein striktes lokales Minimum in $\overline{x}$, dann existiert eine Folge $x_k \rightarrow \overline{x}$ derart, dass $v_k$ ein lokales Minimum in $x_k$ hat.
\end{lemma}
	
\minisec{Beweis}
	Sei o.B.d.A. $\overline{x} = 0, v(\overline{x})=0$. \begin{itemize}
		\item Definiere $\omega_\rho := \min_{\der B_\rho(0)} v > 0$ für $0 < \rho < r$.
		\item Definiere $\rho_k := \inf\{ \rho : \frac{\omega_\rho}{4} \geq \sup_{B_r(0)} |v_k-v|\}$, dann gilt $\rho_k \rightarrow 0$.
	\end{itemize}
	Angenommen, es wäre $\rho_{k_j} \geq \delta > 0$ für eine Folge $k_j \rightarrow \infty$, dann wäre $\frac{\omega_\delta}{4} < \sup_{B_r(0)} |v_{k_j} - v| \rightarrow 0$. Widerspruch! \\
	1. Fall: $\rho_k = 0$. \\
	$\Rightarrow v_k = v \Rightarrow x_k = \overline{x}$ ist lokale Minimalstelle von $v_k$. \\
	2. Fall: $\rho_k > 0$. \\
	$\Rightarrow \min_{\der B_{\rho_k}(0)} v_k \geq \min_{\der B_{\rho_k}(0)} v - \frac{\omega_{\rho_k}}{4} = \frac{3 \omega_{\rho_k}}{4} \geq \frac{\omega_{\rho_k}}{4} \geq v_k(0) \Rightarrow v_k$ hat ein $x_k \in B_{\rho_k}(0)$ als Minimalstelle. \\
	Damit folgt $x_k \rightarrow 0$. \qed

\mbox{} \\
Sei $u_\varepsilon \in C^2(\Omega)$ eine Lösung von \eqref{eq_0605_15} mit $u_\varepsilon \rightarrow u$ lokal gleichmäßig. Betrachte eine Vergleichsfunktion $\phi \in C^2(\Omega)$ mit $\phi \leq u$, sodass $u-\phi$ ein striktes Minimum in $\overline{x} \in \Omega$. \\
Lemma \ref{lemma_21} impliziert, dass eine Folge $\varepsilon_k \rightarrow 0$, $x_k \rightarrow \overline{x}$ existiert, sodass $u_{\varepsilon_k} - \phi$ ein lokales Minimum in $x_k$ besitzt, das heißt:
\[ \begin{array}{rcl}
	\nabla(u_{\varepsilon_k}(x_k) - \phi(x_k)) = 0	&	\hspace*{1cm}	&	\text{(Ableitung = 0)} \\
	D^2(u_{\varepsilon_k}(x_k) - \phi(x_k)) \geq 0	&	&	\text{(Hesse-Matrix positiv semidefinit)} \\
	\Delta u_{\varepsilon_k}(x_k) - \Delta \phi(x_k) \geq 0	&	& \text{(s.o.)}
\end{array} \]
Damit folgt
\begin{equation}
\begin{aligned}
	0 &= -\varepsilon_k \Delta u_{\varepsilon_k}(x_k) + H(x_k,u_{\varepsilon_k}(x_k)) \\ \label{eq_0605_17}
	&\leq -\varepsilon_k \Delta\phi (x_k) + H(x_k,u_{\varepsilon_k}(x_k),\nabla \phi (x_k)) \rightarrow H(\overline{x},u_{\varepsilon_k}(\overline{x}),\nabla \phi (\overline{x}))
\end{aligned}
\end{equation}
das heißt
\[ H(\overline{x},u_{\varepsilon_k}(\overline{x})) \geq 0 \text{, falls } u-\phi \text{ ein striktes lokales Minimum in } \overline{x} \text{ besitzt.} \]
Dies motiviert zu folgendem Begriff:

\begin{defn}[Viskositätslösung] \label{def_22}
	Sei $\Omega \subset \RR^n$ offen, $H\colon \Omega \times \RR \times \RR^n$ stetig. \marginnote{[22]}
	\begin{enumerate}
		\item $u \in C^0(\Omega)$ heißt \bet{Viskositäts-Superlösung}, falls $H(x,u(x),\nabla\phi(x)) \geq 0$ für alle $\phi \in C^1(\Omega)$ derart, dass $u-\phi$ ein Minimum in $x$ besitzt.
		\item $u \in C^0(\Omega)$ heißt \bet{Viskositäts-Sublösung}, falls $H(x,u(x),\nabla\phi(x)) \leq 0$ für alle $\phi \in C^1(\Omega)$ derart, dass $u-\phi$ ein Maximum in $x$ besitzt.
		\item $u \in C^0(\Omega)$ heißt \Index{Viskositätslösung}, falls $u$ Viskositäts-Superlösung und Viskositäts-Sublösung ist.
	\end{enumerate}
\end{defn}

\begin{bsp}[Abstandsfunktion]\label{lemma_23}
	Für $\Gamma \subset \RR^n$ abgeschlossen ist $u(x) = \dist(x,\Gamma)$ eine Viskositätslösung von $|\nabla u| - 1 = 0$ auf $\RR^n \setminus \Gamma$. \marginnote{[23]}
\end{bsp}
	
\minisec{Beweis}
	\begin{enumerate}
		\item zur Sublösung: Sei $x$ ein lokales Maximum von $u - \phi$, das heißt
		\begin{equation}
		\begin{aligned}
			&(u-\phi)(x+z)-(u-\phi)(x) \leq 0 \\ \notag
			\Leftrightarrow \quad &\phi(x+z) - \phi(x) \geq u(x+z) - u(x) \geq -|z| \\
			\Rightarrow \quad &\frac{\phi(x+s\zeta) - \phi(x)}{s} \geq -|\zeta| \text{ und somit } \nabla\phi(x) \cdot \zeta \geq -|\zeta| \ \forall \zeta \in \RR^n 
		\end{aligned}
		\end{equation}
		Wähle $\zeta = - \frac{\nabla \phi(x)}{|\nabla \phi(x)|}$ und erhalte $|\nabla \phi(x)| \leq 1$ oder $|\nabla \phi(x)| - 1 \leq 0$.
		\item zur Superlösung:  Sei $x$ ein lokales Minimum von $u - \phi$, das heißt
		\begin{equation}
		\begin{aligned}
			&(u-\phi)(x+z)-(u-\phi)(x) \geq 0 \\ \notag
			\Leftrightarrow \quad &\phi(x+z) - \phi(x) \leq u(x+z) - u(x) \\
		\end{aligned}
		\end{equation}
		Sei $\overline{x} \in \Gamma$ mit $u(x) = |x-\overline{x}|$ und wähle $z = -s \frac{x-\overline{x}}{|x-\overline{x}}$.
		\begin{equation}
		\begin{aligned}
			&\Rightarrow \quad \frac{ \phi \left(x+s\frac{\overline{x}-x}{|\overline{x}-x|} \right) - \phi(x)}{s} \leq \frac{|\overline{x} - x| - s - |\overline{x}-x|}{s} = -1 \\ \notag
			&\Rightarrow \quad -\nabla \phi(x) \cdot \frac{\overline{x}-x}{|\overline{x}-x|} \geq 1
			\quad \Rightarrow \quad |\nabla\phi(x)| \geq 1 \text{ oder } |\nabla \phi(x)|-1 \geq 0. 
		\end{aligned}
		\end{equation}
		\qed
	\end{enumerate}
	
\begin{bem} \label{bem_24}
	Im Allgemeinen ist $u  = -\dist(x,\Gamma)$ \marginnote{[24]} keine Viskositätslösung von $|\nabla u|-1 = 0$: Sei $u_{\min} = u(\overline{x})$ ein lokales Minimum von $-\dist(x,\Gamma)$ und $\phi \equiv u_{\min}$, sodass $u  -\phi$ ein lokales Minimum in $\overline{x}$ hat. Dann ist $|\nabla \phi| - 1 = -1 < 0$ und $u$ damit keine Viskositäts-Superlösung. \\
	Allerdings ist $u$ eine Viskositätslösung von $1-|\nabla u| = 0$.
\end{bem}
	
	\mbox{} \\
	Wir betrachten nun Hamilton-Jacobi-Gleichungen der Gestalt
	\begin{equation}
		H(x,\nabla u(x)) = 0 \label{eq_0605_18} 
	\end{equation}
	mit $H$ stetig, $H(x,\cdot)\colon \RR^n \rightarrow \RR$ gleichmäßig konvex, $H(x,p) \xlongrightarrow{|p| \rightarrow \infty} \infty$ gleichmäßig in $x$ (Koerzivität) und $H(x,0) \leq 0$.
	
\begin{defn}[Stützfunktion]
	Eine Funktion \marginnote{[25]}
	\begin{equation}
	\begin{aligned}
		L_x\colon \RR^n &\longrightarrow \RR \\ \notag
		w &\longmapsto \sup_{H(x,p) \leq 0} w \cdot p
	\end{aligned}
	\end{equation}
	heißt \Index{Stützfunktion} der Menge $\{p \in \RR^n : H(x,p) \leq 0\}$.
\end{defn}
	
\begin{bsp} \label{bsp_26}
	$H(x,p) = |p| - \frac{1}{v(x)}$ für eine Geschwindigkeit $v(x)$ hat die Stützfunktion $L_x(w) = \frac{|w|}{v(x)}$. \marginnote{[26]}
\end{bsp}
	
\begin{defn}[optische Distanz] \label{def_27}
	\[ \delta(x,y) = \inf \left\{ \int_{0}^{1} L_{c(t)}(\dot{c}(t)) dt : c\colon [0,1] \rightarrow \RR^n, c(0) = x, c(1) = y\right\} \]
	heißt die \Index{optische D<istanz} zwischen $x$ und $y$. \marginnote{[27]}
\end{defn}
	
\begin{bsp} \label{bsp_28}
	Die optische Distanz von $H(x,p) = |p| - \frac{1}{v(x)}$ ist gegeben durch
	\[ \delta(x,y) = \inf_{\substack{c\colon[0,1] \rightarrow \RR^n \\ c(0) = x, c(1) = y}} \int_{0}^{1} \frac{|\dot{c}(t)|}{v(c(t))} dt \]
	und besc<<hreibt die Zeit, die eine seismische Welle mit der Ausbreitungsgeschwindigkeit $v(x)$ benötigt, um von $x$ nach $y$ zu gelangen. \marginnote{[28]}
\end{bsp}

% % % % % % % 09.05.
\begin{thm}[Hopf-Lax-Formel] \label{thm_29}
	Sei $\Omega \subset \RR^n$ offen und beschränkt \marginnote{9. Mai \\ \ [29]} und $g\colon \der \Omega \rightarrow \RR$ mit der Eigenschaft $g(x)-g(y) \leq \delta(y,x)$ für alle $x,y \in \der \Omega$. Dann ist
	\[ u(x) = \inf_{y \in \der \Omega} \{g(y) + \delta(y,x) \} \]
	eine lipschitzstetige Viskositätslösung von $H(x,\nabla u(x)) = 0$ auf $\Omega$ mit $u = g$ auf $\der \Omega$.
\end{thm}
	
\begin{bem} \label{bem_30}
	\begin{itemize}
		\item Dieser Satz impliziert die Existenz einer Viskositätslösung. \marginnote{[30]}
		\item Für $H(x,p) = |p| - \frac{1}{v(x)}$ und $g \equiv 0$ ist $u(x)$ die Ankunftszeit einer seismischen Welle, die auf $\der \Omega$ startet.
		\item $g(x) - g(y) \leq \delta(y,x)$ bedeutet, dass die Wellenfront zum Erreichen von $x$ nicht länger braucht als um von $y$ nach $x$ zu gelangen.
	\end{itemize}
\end{bem}
	
\begin{lemma} \label{bem_31}
	\begin{enumerate}[(i)]
		\item $L_x$ ist konvex. \marginnote{[31]}
		\item $L_x(w) \leq C|w|$ für eine Konstante $C$ unabhängig von $w$.
		\item $L_x(\lambda w) = \lambda L_x(w)$ für alle $\lambda > 0$.
	\end{enumerate}
\end{lemma}
	
\minisec{Beweis}
	\begin{enumerate}[(i)]
		\item $L_x(tq+(1-t)w) = \sup_{H(x,p) \leq 0} [tq + (1-t)w] \cdot p = \sup_{H(x,p) \leq 0} tq \cdot p + (1-t)w \cdot p$ \\
		$\leq t \sup_{H(x,p) \leq 0} q \cdot p + (1-t) \sup_{H(x,p) \leq 0} w \cdot p = tL_x(q) + (1-t)L_x(w)$ für alle $t \in [0,1]$.
		\item $L_x(w) = \sup_{H(x,p) \leq 0} w \cdot p \leq \sup_{H(x,p) \leq 0} |w| \cdot |p| \leq (\sup_{H(x,p) \leq 0} |p|) |w|$
		\item $L_x(\lambda w) = \sup_{H(x,p) \leq 0} \lambda w \cdot p = \lambda \sup_{H(x,p) \leq 0} w \cdot p$ \qed
	\end{enumerate}
	
\begin{bem} \label{bem_32}
	Das obige Lemma impliziert, dass $\delta$ eine Pseudometrik ist, falls $H(x,p) = H(x,-p)$. \marginnote{[32]}
\end{bem}
	
\minisec{Beweis von Theorem \ref{thm_29}}
	Wir betrachten den Fall, dass $H(x,P) \equiv H(p)$ und $H(p) = H(-p)$. In diesem Fall gilt $\delta(x,y) = L(y-x)$, da $L$ nicht von $x$ abhängig ist. \begin{enumerate}[a)]
		\item $u$ ist lipschitzstetig und $u(x) \leq u(y) + \delta(x,y)$ für alle $x,y \in \overline{\Omega}$:
		\begin{equation}
		\begin{aligned}
			u(x) - u(y) &= \inf_{z_1 \in \der\Omega} (g(z_1) + \delta(x,z_1)) - \inf_{z_2 \in \der \Omega} (g(z_2) + \delta(y,z_2)) \\ \notag
			&\leq \sup_{z_2 \in \der \Omega} g(z_2) + \der(x,z_2) - g(z_2) - \der(y,z_2) \\
			&\leq \der(x,y) = L(y-x) \overset{(ii)}{\leq} C|y-x|
		\end{aligned}
		\end{equation}
		\item Sublösung: \\
		$u-\phi$ habe ein lokales Maximum in $x$, d.h. $u(x')-\phi(x') \leq u(x) - \phi(x)$, bzw.
		\[ \phi(x') - \phi(x) \geq u(x') - u(x) \geq -\delta(x,x') \]
		Mit $\zeta$ beliebig und $x' = x + s \zeta$ folgt
		\[ \frac{\phi(x+s\zeta)-\phi(x)}{s} \geq \frac{-\delta(x,x+s\zeta)}{s} \overset{\substack{L \text{ unabh.} \\ \text{von } x}}{=} \frac{-L(s\zeta)}{s} \overset{(iii)}{=} -L(\zeta) \quad \Rightarrow \quad \nabla \phi \cdot \zeta \geq -L(\zeta) \]
		Analog folgt $\nabla \phi \cdot \zeta \leq L(\zeta) = \sup_{H(p) \leq 0} \zeta \cdot p$, also $|\nabla \phi \cdot \zeta| \leq L(\zeta)$ für alle $\zeta \in \RR^n$. Angenommen, $H(\nabla \phi) > 0$. Da $H$ konvex und stetig ist, existiert eine Hyperfläche mit Normale $\nu$, die $\{p : H(p) \leq 0\}$ und $\nabla \phi$ trennt. Damit wäre $|\nabla \phi \cdot \nu| > \sup_{H(p) \leq 0} p \cdot \nu = L(\nu)$. Widerspruch.
		\item Superlösung: \\
		Sei $u(x) = g(\overline{y}) + \delta(x,\overline{y}) = g(\overline{y}) + L(x-\overline{y})$ für ein $\overline{y} \in \der \Omega$ und definiere $c(t) = \overline{y} + t(x-\overline{y})$. Wir haben
		\begin{equation}
		\begin{aligned}
			g(\overline{y}) + L(x-\overline{y}) &= u(x) \leq u(c(t)) + \delta(c(t),x) ) u(c(t)) + (1-t)L(x-\overline{y}) \\ \notag
			&\leq g(\overline{y} + \delta(\overline{y},c(t)) + (1-t)L(x-\overline{y}) = g(\overline{y}) + L(x-\overline{y})
		\end{aligned}
		\end{equation}
		sodass $u(c(t)) = g(\overline{y}) + tL(x-\overline{y}) = g(\overline{y}) + L(c(t)-\overline{y})$. Nun besitze $u-\phi$ ein lokales Minimum in $x$, d.h. $\phi(x) - \phi(x') \geq u(x) - u(x')$, und setze $x=c(1), x'=c(1-s)$. Wir erhalten:
		\begin{equation}
		\begin{aligned}
			&\frac{\phi(x)-\phi(c(1-s))}{s} \geq \frac{u(x) - u(c(1-s))}{s} \\ \notag
			\Rightarrow \quad &\nabla \phi(x) \cdot (x-\overline{y}) = \nabla \phi(x) \cdot \dot{c}(1) \geq \nabla(u) \cdot \dot{c}(1) = L(x-\overline{y}) = \sup_{H(p) \leq 0} (x-\overline{y}) \cdot p
		\end{aligned}
		\end{equation}
		Daher folgt $(\nabla \phi(x) + \alpha(x-\overline{y})) \cdot (x-\overline{y}) > \sup_{H(p) \leq 0} (x-\overline{y})\cdot p$ für alle $\alpha > 0$ und somit $H(\nabla \phi(x) + \alpha(x- \overline{y})) > 0$ und $H(\nabla \phi) \geq 0$ wegen Stetigkeit.
		\item Randdaten: \\
		Sei $x \in \der\Omega$, dann folgt aus $g(x)-g(y) \leq \delta(y,x)$, dass $g(x) \leq g(y) + \delta(y,x)$ für alle $y \in \der\Omega$ und somit $g(x) \leq u(x)$. Des Weiteren ist $u(x) \leq g(x) + \delta(x,x) = g(x)$ und damit $u(x) = g(x)$. \qed
	\end{enumerate}

\mbox{} \\
Für Eindeutigkeit müssen wir zusätzliche Bedingungen an $H$ stellen, wie das folgende Beispiel zeigt. Anschließend formulieren wir zwei mögliche Resultate für die Eindeutigkeit.

\begin{bsp} \label{bsp_33}
	Sei $\Psi \in C^1(\Omega)$, $\Psi = 0$ auf $\der\Omega$,\marginnote{[33]} und betrachte $H(x,p) = |p|^2 - |\nabla \Psi(x)|^2$. Dann ist sowohl $u = \Psi$ als auch $u = -\Psi$ eine (Viskositäts-)Lösung von $0 = H(x,\nabla u (x))$.
\end{bsp}
	
\begin{thm}[Eindeutigkeit durch Vergleich] \label{thm_34}
	Sei $\Omega \subset \RR^n$ beschränkt \marginnote{[34]} und offen sowie $H\colon \overline{\Omega} \times \RR \times \RR^n$ stetig mit \begin{itemize}
		\item $H(x,u,p) - H(x,v,p) > \gamma(u-v)$ für ein $\gamma > 0$.
		\item $|H(x,u,p) - H(y,u,p)| \leq C|y-x|(1+|p|)$ für ein $C > 0$.
	\end{itemize}
	Falls $u$ eine Viskositäts-Sublösung und $v$ eine Viskositäts-Superlösung von $0 = H(x,u(x),\nabla u(x))$ ist mit $u \leq v$ auf $\der \Omega$, dann ist auch $u \leq v$ auf $\Omega$. Damit folgt, dass die Viskositätslösung eindeutig ist.
\end{thm}
	
\minisec{Beweisidee}
	Angenommen, $u$ und $v$ seien glatt und $u-v$ habe ein Maximum in $x_0 \in \overline{\Omega}$ mit $u(x_0) - v(x_0) > 0$. Nach Definition \ref{def_22} haben wir
	\[ \begin{array}{c}
		H(x_0,u(x_0),\nabla v(x_0)) \leq 0 \\
		H(x_0,v(x_0),\nabla u(x_0)) \geq 0
	\end{array}\]
	Wegen $\nabla(u-v)(x_0) = 0$ gilt $\nabla u(x_0) = \nabla v(x_0)$, und damit den Widerspruch
	\[ 0 \geq H(x_0,u(x_0),\nabla v(x_0)) = H(x_0,u(x_0),\nabla(u(x_0)) > H(x_0,v(x_0),\nabla u(x_0)) \geq 0 \]
	Für nichtglatte $u,v$ ist der Beweis aufwändiger. \qed
	
\begin{thm}[Eindeutigkeit für Hamilton-Jacobi-Bellman-Gleichung] \label{thm_35}
	Sei \marginnote{[35]}
	\begin{equation}
	\begin{aligned}
		H\colon (\RR \times \RR^{n-1}) \times \RR^n &\longrightarrow \RR \\ \notag
		((t,x),(p^t,p^x)) &\longmapsto p^t + \tilde{H}(x,p^x)
	\end{aligned}
	\end{equation}
	mit $\tilde{H}$ stetig und \begin{itemize}
		\item $|\tilde{H}(x,p) - \tilde{H}(x-q)| \leq C|p-q|$ für ein $C > 0$
		\item $|\tilde{H}(x,p) - \tilde{H}(y,p)| \leq C|y-x|(1+|p|)$
	\end{itemize}
	Dann existiert höchstens eine Viskositätslösung von $0 = H((t,x),(u_t(t,x),\nabla u(t,x))) = u_t + \tilde{H}(x,\nabla u)$ mit gegebenen Randdaten auf $t = 0$.
\end{thm}

\minisec{Beweis}
	siehe z.B. Evans, \glqq PDEs\grqq, S. 587 
\newpage
% % % % % % % % % % % 13. Mai
\section{Partielle Differentialgleichungen zweiter Ordnung -- Elliptische PDGL}
\label{sec:para7}
	In diesem Abschnitt betrachten wir semilineare Gleichungen der Gestalt \marginnote{13. Mai}
	\begin{equation}
		\sum\limits_{i,j=1}^{n} a_{ij}(x) u_{x_ix_j}(x) + c(\nabla u(x),u(x),x) = 0 \label{eq_0905_19}
	\end{equation}
	mit einer symmetrischen Matrix $A := (a_{ij})_{ij}$.
	
\begin{defn}[Klassifikation einer PDGL zweiter Ordnung] \label{def_36}
	Die Gleichung \eqref{eq_0905_19} heißt \marginnote{[36]} \begin{itemize}
		\item \Index{elliptisch}, falls $A$ genau $n$ positive Eigenwerte besitzt
		\item \Index{parabolisch}, falls $A$ den Eigenwert $0$ besitzt
		\item \Index{hyperbolisch}, falls $A$ einen positiven und $n-1$ negative Eigenwerte besitzt.
	\end{itemize}
\end{defn}
	
\mbox{} \\
Da eine Multiplikation mit $-1$ die PDGL nicht verändert, lassen sich in der obigen Definition die Begriffe positiv und negativ vertauschen. Nichtlineare PDGL können durch Linearisierung und Betrachtung der zweiten Ableitungen lokal klassifiziert werden. Kleine Änderungen in $A$ können dazu führen, dass die PDGL ihren Typ ändert.

\begin{defn}[Laplace-Gleichung und Harmonische Funktion]
\label{def_37} \label{sub:def_harmonische_fkt}
	Eine zweimal stetig differenzierbare Funktion, die die \Index{Laplace-Gleichung}
	\begin{equation}
		\Delta u = 0 \label{eq_1305_20}
	\end{equation}
	erfüllt, heißt \Index{harmonische Funktion}.
\end{defn}
	
\mbox{} \\
Harmonische Funktionen besitzen einige praktische Eigenschaften, die alle mehr oder weniger mit der Glattheit der Laplace-Gleichung zusammenhängen.

\begin{thm}[Mittelwertformel] \label{thm:mittelwertformel} \label{thm_38}
	$u$ ist genau dann harmonisch auf $\Omega$, falls $u \in C^2$ und \marginnote{[38]}
	\begin{equation}
		u(x) = \int_{\der B_r(x)} u(\tilde{x})d\tilde{x} \label{eq_1305_21}
	\end{equation}
	für jede Kugel $B_r(x) \subseteq \Omega$.
\end{thm}
	
\minisec{Beweis}
	\begin{itemize}
		\item[\glqq$\Rightarrow$\grqq:] Definiere $f(\tau) = \int_{\der B_r(x)} u(\tilde{x}) d\tilde{x} = \int_{B_1(0)} u(x+rz) dz.$ \\
		$f(0) = u(x)$ und $f'(\tau) = \int_{\der B_1(0)} \nabla u(x+rz)\cdot z dz = \int_{\der B_r(x)} \nabla u(\tilde{x}) \cdot \nu(\tilde{x}) d \tilde{x} = \frac{1}{|\der B_1(x)|} \int_{B_r(x)} \underbrace{\Delta u}_{=0} d \tilde{x}$
		\item[\glqq$\Leftarrow$\grqq:] Angenommen, $\Delta u(x) > 0$, d.h. es existiert $B_r(x)$ mit $0 < \int_{B_r(x)} \delta u d \tilde{x}$, aber \\ $0 = f'(r) = \frac{1}{|\der B_1(x)|} \int_{B_r(x)} \Delta u d \tilde{x} > 0$ $\lightning$ \qed
	\end{itemize}

\begin{bem} \label{bem_39}
	Benutzung \marginnote{[39]} von $\int_{B_r(x)} u(\tilde{x})d\tilde{x} = \int_{0}^{r} \left( \int_{\der B_s(x)} u(\tilde{x}) d \tilde{x} \right) ds = u(x) \int_0^r \left( \int_{\der B_s(x)} d \tilde{x} \right) ds = u(x)$ liefert eine zweite Mittelwertformel.
\end{bem}
	
\begin{thm}[Starkes Maximumprinzip] \label{thm:starkes_maximum} \label{thm_40}
	Eine nichtkonstante harmonische Funktion $u \in C^2(\Omega) \cap C^0(\overline{\Omega})$ hat kein lokales Maximum auf $\Omega$.\marginnote{[40]} Insbesondere gilt:
	\[ \max_{\overline{\Omega}} u = \max_{\der \Omega} u \]
	(Analoges gilt für Minima: ersetze $u$ durch $-u$)
\end{thm}
	
\minisec{Beweis}
	Sei $u(x_0) = M := \max_{\overline{\Omega}} u$ für ein $x_0 \in \Omega$. Nach Theorem \ref{thm_38} bzw. Bemerkung \ref{bem_39} gilt jedoch $M = u(x_0) = \int_{B_r(x_0)} u dx < M$, sofern nicht $u = M$ auf ganz $B_r(x_0)$. Benutze dieses Argument für jeden Punkt in $\der B_r(x_0)$, anschließend für die neuen Randpunkte usw. und erhalte $u = M$ auf ganz $\Omega$. \qed

\begin{thm}[Harnack-Ungleichung] \label{thm:harnack-ungl} \label{thm_41}
	Für jede zusammenhängende offene Menge $U \subset \Omega$ existiert eine Konstante $C = C(U)$, sodass \marginnote{[41]}
	\[ \sup_u u \leq C \inf_u u \]
	für jede nichtnegative harmonische Funktion $u$ auf $\Omega$.
\end{thm}
	
\minisec{Beweis}
	Sei $r := \frac{1}{4} \dist(u,\der \Omega)$ und $x,y \in U$ beliebig mit $|x-y| \leq r$. Dann:
	\[ u(x) = \int_{B_{2r}(x)} ud \tilde{x} \geq \frac{1}{2^n} \int_{B_r(y)} ud\tilde{x} = \frac{1}{2^n} u(y)\]
	Da $U$ zusammenhängend und $\overline{U}$ kompakt, können wir $\overline{U}$ überdecken durch eine endlich viele Umgebungen $\{B_i\}_{i = 0}^N$ mit Radius $\frac{r}{2}$ und $B_i \cap B_{i-1} \neq \emptyset$. Damit folgt $u(x) \geq \left( \frac{1}{2^n} \right)^{N+1} u(y)$ für alle $x,y \in U$. \qed
	
\begin{thm}[Glattheit harmonischer Funktionen] \label{thm:glattheit_harm_fkt} \label{thm_42}
	Eine harmonische Funktion $u \in C^0(\overline{\Omega})$ ist unendlich oft differenzierbar auf $\Omega$. \marginnote{[42]}
\end{thm}
	
\minisec{Beweis}
	Definiere 
	\[\eta_\varepsilon(x) = \widetilde{\eta_\varepsilon}(|x|) = \frac{1}{\varepsilon^n} \eta \left( \frac{|x|}{\varepsilon} \right) \text{ für } \eta(|x|)= \begin{cases}
		Ce^{\frac{1}{|x|^2-1}} & \text{ für } |x| < 1 \\
		0 & \text{ sonst} \end{cases} \]
	mit $C$ derart, dass $\int_{\RR^n} \eta(|x|) dx = 1$. Dann ist $\eta_\varepsilon \in C^\infty(\RR^n)$. \\
	Setze $u_\varepsilon(x) = (u * \eta_\varepsilon)(x) = \int_\Omega u(y) \eta_\varepsilon(x-y) dy \in C^\infty(\Omega_\varepsilon)$ für $\Omega_\varepsilon = \{x \in \Omega : \dist(x,\der \Omega) > \varepsilon\}$. Dann ist $u_\varepsilon = u$ auf $\Omega_\varepsilon$:
	\begin{equation}
	\begin{aligned}
		u_\varepsilon(x) &= \int_{\Omega} u(y) \widetilde{\eta_\varepsilon}(|x-y|)dy = \int_{B_\varepsilon(x)} \widetilde{\eta_\varepsilon} (|x-y|)u(y) dy \\ \notag
		&= \int_0^\varepsilon \widetilde{\eta_\varepsilon}(r) \underbrace{\left( \int_{\der B_r(x)} u(\tilde{x}) d \tilde{x} \right)}_{u(x) |\der B_r(x)|} dr = u(x) \underbrace{\int_{B_\varepsilon(0)} \eta_\varepsilon(\tilde{x}) d\tilde{x}}_{= \int_{B_1(0)} \eta(|x|) dx = 1} = u(x) 
	\end{aligned}
	\end{equation}
	\qed
	
% % % % 16. Mai % % % %
\begin{defn}[Dirichlet-Problem, Neumann-Problem]
\label{def:dirichlet_neumann}
	Das \Index{Dirichlet-Problem} ist gegeben durch \marginnote{16. Mai}
	\begin{equation}
		\begin{cases} -\Delta u(x) = f(x) \\
		u=g \text{ auf } \der \Omega \end{cases} \label{dirichlet_problem}
	\end{equation}
	Das \Index{Neumann-Problem} ist gegeben durch
	\begin{equation}
		\begin{cases} -\Delta u(x) = f(x) \\
		\frac{\der u}{\der \nu} \text{ auf } \der \Omega \end{cases} \label{neumann_problem}
	\end{equation}
\end{defn}
	
\begin{thm}[Eindeutige Lösbarkeit des Dirichlet- und Neumann-Problems] \label{thm_43}
	Die Lösung des Dirichlet-Problems \eqref{dirichlet_problem} ist eindeutig, \marginnote{[43]} falls existent. Das Neumann-Problem \eqref{neumann_problem} besitzt nur dann eine Lösung, falls gilt:
	\[ \int_\Omega f(x) dx = - \int_{\der \Omega} g(x) dx \]
	In diesem Fall ist sie eindeutig bis auf eine konstanten Term.
\end{thm}
	
\minisec{Beweis}
	Seien $u, v$ zwei Lösungen, dann ist $w := u-v$ eine Lösung von $\Delta w = 0$. Nach Theorem \ref{thm_40} nimmt $w$ sein Maximum und Minimum auf $\der \Omega$ an. Im Fall von \eqref{dirichlet_problem} ist ferner $w = 0$ auf $\der \Omega$, also folgt $w \equiv 0$. Im Fall \eqref{neumann_problem} ist $\der w / \der \nu = 0$ auf $\der \Omega$, also ist $w$ konstant. \\
	Beachte: Im zweiten Fall haben wir
	\[ \int_{\Omega} f(x) dx = \int_\Omega - \Delta u(x) dx = -\int_{\der \Omega} \nabla u \cdot \nu dx = - \int_{\der \Omega} g(x) dx\]
	\qed
	
\mbox{} \\
Wir erarbeiten nun eine praktische Formel für die Lösung der \Index{Poisson-Gleichung}
\begin{equation}
	-\Delta u(x) = f(x) \label{poissongl}
\end{equation}
Dafür betrachten wir zunächst Lösungen der Laplace-Gleichung. Sei $r := |x|$ und $v(r) = u(x)$ eine Lösung der Laplace-Gleichung auf $\RR^n \setminus \setnull$. Durch Benutzung von $\frac{\der r}{\der x_i} = \frac{x_i}{r}$, also $u_{x_i} = \frac{v'(r) x_i}{r}$ und $u_{x_ix_i} = \frac{v''(r)x_i^2}{r^2} + v'(r) \left( \frac{1}{r}-\frac{x_i^2}{r^3} \right)$, erhalten wir
\[ v''(r) + \frac{n-1}{r} v'(r) = 0 \]
Diese Gleichung besitzt die Lösung $v''(r) = \frac{a}{r^{n-1}}$ für eine Konstante $a$.

\begin{defn}[Fundamentallösung der Laplace-Gleichung] \label{def:fundamentallsg_laplace} \label{def_44}
	Die Funktion \marginnote{[44]}
	\[ \Phi(x) = \begin{cases}
		-\frac{1}{2}|x| & (n=1) \\
		-\frac{1}{2\pi} \log |x| & (n=2) \\
		-\frac{1}{(n-2)|\der B_1(0)|}\frac{1}{|x|^{n-2}} & (n \geq 3) \end{cases} \]
	löst die Laplace-Gleichung auf $\RR^n \setminus \setnull$ und heißt \Index{Fundamentallösung} der Laplace-Gleichung.
\end{defn}
	
\begin{defn}[Delta-Verteilung] \label{def:delta_vert} \label{def_45}
	Der lineare Operator $\widehat{\delta}\colon C^0(\RR^n) \rightarrow \RR$ mit $\widehat{\delta}(u) = u(0)$ heißt \bet{$\delta$-Verteilung}. Üblich ist auch die Notation \marginnote{[45]}
	\[ \widehat{\delta}(u) = \int_{\RR^n} \delta(x) u(x) dx \]
	wobei $\delta$ eine Funktion mit $\int_{\RR^n} \delta(x) dx = 1$ ist, die im Punkt 0 unendlich ist und an allen anderen Stellen den Wert 0 annimmt.
\end{defn}
	
\mbox{} \\
Multiplizieren wir $\Delta \Phi$ mit einer glatten Funktion $\psi$ mit $\psi = 0$ auf $\der \Omega$ und integrieren wir zweimal partiell, erhalten wir:
\[ \int_{\Omega} \psi(x) \Delta \Phi(x) dx = \int_{\Omega} \Phi(x) \Delta \psi(x) dx \]

\begin{thm}[Fundamentallösung] \label{thm_46} \label{fundamentallsg}
	Für die $\delta$-Verteilung gilt: \marginnote{[46]}
	\[ \Delta \Phi(x) = \delta(x) \]
	d.h. es gilt $\Delta \Phi(x) = 0$ auf $\RR^n \setminus \setnull$ und $\int_{\RR^n} \Phi(x) \Delta \psi(x) dx = \psi(0)$ für alle glatten Funktionen $\psi$ mit kompaktem Träger.
\end{thm}
	
\minisec{Beweis}
	\todo{fehlt bisher noch!}
	
\mbox{} \\	
Sei nun $\der \Omega$ Lipschitz. So wie zuvor betrachte die Lösung des folgenden Problems:
\begin{equation}
	\begin{cases}
		- \Delta G^y(x) = \delta(x-y) & \text{ auf } \Omega \\
		G^y = 0 & \text{ auf } \der \Omega \label{eq_25}
	\end{cases}
\end{equation}
Motivation: Finden wir $G^y$ für alle $y \in \Omega$, dann genügt
\begin{equation}
	u(x) = \int_{\Omega} G^x(y) f(y) dy \label{eq_26}
\end{equation}
der Gleichung
\begin{equation}
	-\Delta u(x) = \int_{\Omega} -\Delta G^x(y) f(y) dy = f(y)
\end{equation}

\begin{bem} \label{bem_47}
	Offensichtlich ist $G^y(x) = \Phi(x-y)-\phi^y(x)$, wobei $\phi^x$ eine Lösung ist von \marginnote{[47]}
	\begin{equation} \begin{cases}
		\Delta \phi^y = 0 & \text{ auf } \Omega \\
		\phi^y = \Phi(x-y) & \text{ auf } \der \Omega \end{cases} \label{eq_27}
	\end{equation}
\end{bem}

\begin{thm}[Greensche Formel] \label{thm_48} \label{green_formel}
	Falls $u \in C^2(\overline{\Omega})$ eine Lösung des Dirichlet-Problems \eqref{dirichlet_problem} ist, dann gilt \marginnote{[48]}
	\[ u(x) = -\int_{\der \Omega} g(y) \frac{\der G^x(y)}{\der \nu} dy + \int_{\Omega} f(y) G^x(y) dy \]
\end{thm}
	
\minisec{Beweis}
	Nach Theorem \ref{thm_46} gilt für $x \in \Omega$:
	\begin{equation}
	\begin{aligned}
	 u(x) 	&= \int_{\Omega} u(y) \Delta \Phi(y-x) dy \\ \notag
			&= \int_{\Omega} u(y) \Delta G^x(y) dy \\
			&= \int_{\der \Omega} u(y) \nabla G^x(y) \cdot \nu dy - \int_\Omega \nabla u(y) \cdot \nabla G^x(y) dy \\
			&= \int_{\der \Omega} g(y) \frac{\der G^x(y)}{\der \nu} dy - \int_{\der \Omega} \nabla u(y) \cdot \nu \underbrace{G^x(y)}_{=0} dy + \int_{\Omega} \underbrace{\Delta u(y)}_{=f(y)} G^x(y) dy
	\end{aligned}
	\end{equation} \qed
	
\begin{bsp}[Greensche Funktion für Halbräume] \label{bsp_49}
	Die Greensche Funktion $G^y$ \marginnote{[49]} für $y \in \Omega = \{x \in \RR^n : x_n > 0\}$ kann mit Hilfe der \Index{Spiegelungsmethode} (method of images) gefunden werden: $\phi^y(x) = \Phi(x-(-y))$ erfüllt Gleichung $\eqref{eq_27}$, sodass
	\[ G^y(x) = \Phi(x-y) - \Phi(x+y) \]
	Speziell in 2D:
	\[ G^y(x) = \frac{1}{4\pi} \log \enbrace*{ \frac{|x-y|^2}{|x+y|^2}}\]
\end{bsp}
	
\begin{bsp}[Greensche Funktion für eine Scheibe] \label{bsp_50}
	Die Greensche Funktion $G^y$ \marginnote{[50]} für $y \in \Omega = B_r(0)$ kann ähnlich gefunden werden: Für $\widetilde{y} = \frac{r^2}{|y|^2}y$ ist $\frac{|x-y|}{|x+\widetilde{y}|} = \frac{|y|}{r}$ konstant auf $x \in \der \Omega$. Damit erhalten wir:
	\[ G^y(x) = \Phi(x-y) - \Phi \enbrace*{(x-\widetilde{y}) \frac{|y|}{r}} \]
\end{bsp}

\mbox{} \\
Ein ähnliches Vorgehen ist für das Neumann-Problem möglich.
\newpage
\section{Exkursion: Hölder- und Sobolevräume}
\label{sec:para8}
	Um die Existenz und Regularität von elliptischen Gleichungen\marginnote{20. Mai} zu verstehen, müssen wir einige Funktionenräume einführen. \\
	Für eine stetige Funktion $u \in C^0(\overline{\Omega})$ auf $\Omega \subset \RR^n$ offen und beschränkt und $\gamma \in [0,1]$ definiere:
	\[ \lbrack u \rbrack_\gamma = \sup_{x,y \in \overline{\Omega}, x \neq y} \frac{|u(x)-u(y)|}{|x-y|^\gamma} \]
	
\begin{defn}[Hölderraum] \label{def_51} \label{hoelderraum}
	Für $u \in C^k(\overline{\Omega})$ definiere die \Index{Höldernorm}: \marginnote{[51]}
	\[ ||u||_{C^{k,\gamma}(\overline{\Omega})} := \sum_{|\alpha| \leq k} ||D^\alpha u||_{C^0(\overline{\Omega})} + \sum_{|\alpha| = k} \lbrack D^\alpha u \rbrack_\gamma \]
	Der Funktionenraum
	\[ C^{k,\gamma}(\overline{\Omega}) := \{ u \in C^k(\overline{\Omega}) : ||u||_{C^{k,\gamma}(\overline{\Omega})} < \infty \} \]
	heißt \Index{Hölderraum} mit Exponent $\gamma$.
\end{defn}
	
\begin{thm}[Hölderräume sind Banachräume] \label{thm_52}
	Hölderräume mit der Höldernorm sind Banachräume, d.h. jede Cauchyfolge in einem Hölderraum konvergiert. \marginnote{[52]}
\end{thm}
	
\minisec{Beweis}
	Übung.
	
\mbox{} \\
Man stelle fest, dass $C^{k,0} = C^k$. Ferner ist $C^{0,1}$ der Raum der lipschitzstetigen Funktionen. \\
Als nächstes führen wir einen abgeschwächten Begriff der Differenzierbarkeit ein.

\begin{defn}[Schwache Ableitung] \label{def:schwache_abl} \label{def_53}
	Sei $u, v \in L^1_{loc}(\Omega)$ und $\alpha$ ein Multiindex. $v$ heißt die $\alpha$-te \Index{schwache Ableitung} von $u$, falls \marginnote{[53]}
	\begin{equation}
		\int_{\Omega} uD^\alpha \psi dx = (-1) |\alpha| \int_{\Omega} v \psi dx \label{eq_28}
	\end{equation}
	für alle Testfunktionen $\psi \in C_c^\infty(\Omega)$ (d.h. für alle unendlich oft differenzierbaren Funktionen mit kompaktem Träger in $\Omega$). Für die schwache Ableitung schreiben wir auch
	\[ D^\alpha u = v \]
\end{defn}

\begin{bem} \label{bem_54}
	Falls $u$ glatt ist, \marginnote{[54]} ist \eqref{eq_28} das exakte Ergebnis, welches man durch $k$-fache partielle Integration erhält, und $v$ ist die gewöhnliche Ableitung von $u$.
\end{bem}
	
\begin{bsp} \label{bsp_55}
	Sei $\Omega = (0,2)$. \marginnote{[55]}
	\begin{itemize}
		\item Sei $u(x) = \begin{cases}
		x & \text{ falls } 0 < x \leq 1 \\
		1 & \text{ falls } 1 < x < 2
		\end{cases}$ und $v(x) = \begin{cases}
		1 & \text{ falls } 0 < x \leq 1 \\
		0 & \text{ falls } 1 < x < 2
		\end{cases}$, dann ist $v = Du$, da für alle $\psi \in C_c^\infty(\Omega)$ gilt:
		\[ \int_0^2 u\psi' dx = \int_0^1 x \psi' dx + \int_1^2 \psi' dx = -\int_0^1 \psi dx + \psi(1) - \psi(1) = -\int_0^2 v\psi dx \]
		\item Sei $u(x) = \begin{cases}
		 x & \text{ falls} 0 < x \leq 1 \\
		 2 & \text{ falls} 1 < x < 2 \end{cases}$, dann besitzt $u$ keine schwache Ableitung, denn kein $\psi \in C_c^\infty(\Omega)$ erfüllt
		 \[ - \int_0^2 v \psi dx = \int_0^2 u \psi' dx = \int_0^1 x \psi'dx + 2 \int_1^2 \psi' dx = -\int_0^1 \psi dx - \psi(1) \]
		 für beliebiges $v \in L_{loc}^1(\Omega)$.
	\end{itemize}
\end{bsp}
	
\begin{defn}[Lebesgue-Raum] \label{def:lebesgue_raum}
	Sei $p \in [1,\infty]$. Dann definieren wir die \Index{Lebesgue-Norm} durch
	\[ ||u||_{L^p(\Omega)} = \begin{cases}
		\enbrace*{\int_{\Omega} |u|^p dx}^{1/p} & (p < \infty) \\
		\esssup_\Omega |u| & (p = \infty)
	\end{cases} \]
	Der Funktionenraum
	\[ L^p(\Omega) := \{u \colon \Omega \rightarrow \RR : u \text{ ist messbar mit } ||u||_{L^p(\Omega)} < \infty\} \]
	heißt \index{Lebesgue-Raum} mit Exponent $p$.
\end{defn}
	
\begin{thm}[Lebesgue-Räume sind Banachräume]
	Lebesgue-Räume mit der Lebesgue-Norm sind Banachräume.
\end{thm}
	
\begin{defn}[Sobolev-Raum] \label{def:sobolevraum} \label{def_56}
	Der Funktionenraum \marginnote{[56]}
	\[ W^{k,p}(\Omega) := \{ u \in L_{loc}^1(\Omega) : \text{ für alle } |\alpha| \leq k \text{ existiert  die schwache Ableitung } D^\alpha u \text{ mit } D^\alpha u \in L^p(\Omega) \} \]
	mit der \Index{Sobolevnorm}
	\[ ||u||_W^{k,p}(\Omega) = \begin{cases}
		\enbrace*{ \sum_{|\alpha| \leq k} \int_{\Omega} |D^\alpha u|^p dx}^{1/p} & 1 \leq p < \infty \\
		\sum_{|\alpha| = k } \esssup_\Omega |D^\alpha u| & p = \infty \end{cases} \]
	heißt \Index{Sobolevraum}. \\
	$W_0^{k,p}(\Omega)$ ist der Abschluss von $C_c^\infty(\Omega)$ in $W^{k,p}(\Omega)$. Beachte: $W^{0,p}(\Omega) = L^p(\Omega)$.
\end{defn}

\begin{thm}[Sobolevräume sind Banachräume]\label{thm_57}
	Sobolevräume mit der Sobolevnorm sind Banachräume. \marginnote{[57]}
\end{thm}
	
\minisec{Beweis}
	siehe z.B. Evans, S. 262.
	
\begin{bem} \label{bem_58}
	Die Räume \marginnote{[58]}
	\[ H^k(\Omega) = W^{k,2}(\Omega) \]
	sind Hilberträume. Ihre Norm wird induziert durch das Skalarprodukt
	\[ (u,v)_{H^k(\Omega)} = \sum_{|\alpha| \leq k} \int_{\Omega} D^\alpha u D^\alpha v dx \] 
\end{bem}

\begin{thm}[Hölder-Ungleichung] \label{hoelder_ungl} \label{thm_59}
	Sei $p, p^* \in [1,\infty], f \in L^p, g \in L^{p^*}$ und $\frac{1}{p} + \frac{1}{p^*} = 1$. Dann gilt: \marginnote{[59]}
	\[ \int_{\Omega} |fg|dx \leq ||f||_{L^p(\Omega)} ||g||_{L^{p^*}(\Omega)} \]
\end{thm}
	
\minisec{Beweis}
	siehe Alt, \glqq Lineare Funktionalanalysis\grqq, S. 52
	
\begin{thm}[Spursatz] \label{thm:spursatz} \label{thm_60}
	Sei $\Omega$ beschränkt und mit Lipschitzrand.\marginnote{[60]} Dann existiert ein stetiger linearer Operator $T\colon W^{1,p}(\Omega) \rightarrow L^p(\der \Omega)$, die \Index{Spur}, mit
	\begin{enumerate}[(i)]
		\item $Tu = u \big|_{\der \Omega}$, falls $u \in W^{1,p}(\Omega) \cap C^0(\overline{\Omega})$
		\item $||Tu||_{L^p(\der \Omega)} \leq C||u||_{W^{1,p}(\Omega)}$
		\item $Tu=0 \quad \Leftrightarrow \quad u \in W_0^{1,p}(\Omega)$
	\end{enumerate}
	wobei $C \geq 0$ nur von $p$ und $\Omega$ abhängt. Aus Gründen der Einfachheit meinen wir einfach  $u$ auf $\der \Omega$, wenn wir von der Spur von $u$ sprechen.
\end{thm}
	
\minisec{Beweis}
	siehe Evans, S. 272.
	
\begin{thm}[Poincaré-Ungleichung] \label{thm:poincare_ungl} \label{thm_61}
	Sei $\Omega \subset \RR^n$ beschränkt, offen, \marginnote{[61]} zusammenhängend und mit Lipschitzrand. Dann existiert eine Konstante $C = C(n,p,\Omega)$ mit 
	\[ ||u- \frac{1}{|\Omega|} \int_{\Omega} u dx||_{L^p(\Omega)} \leq C \cdot ||\nabla u||_{L^p(\Omega)} \]
	für alle $u \in W^{1,p}(\Omega)$ und
	\[ ||u||_{L^p}(\Omega) \leq C \cdot ||\nabla u||_{L^p(\Omega)} \]
	für alle $u \in W_0^{1,p}(\Omega)$. \index{Poincaré-Ungleichung}
\end{thm}
	
\minisec{Beweis}
	siehe Alt, \glqq Lineare Funktionalanalysis\grqq, S. 171.
	
\begin{thm}[Einbettungssatz von Sobolev] \label{thm:sobolev_einbettung} \label{thm_62}
	Sei $\Omega \subset \RR^n$ offen, beschränkt und mit Lipschitzrand, $m_1,m_2 \in \NN_0$ und $p_1,p_2 \in [1,\infty)$. Falls \marginnote{[62]}
	\[ m_1 \geq m_2 \text{ und } m_1 - \frac{n}{p_1} \geq m_2 - \frac{n}{p_2}, \]
	dann ist $W^{m_1,p_1}(\Omega) \subset W^{m_2,p_2}(\Omega)$ und es existiert eine Konstante $C > 0$, sodass für alle $u$ gilt:
	\[ ||u||_{W^{m_1,p_1}(\Omega)} \leq C \cdot ||u||_{W^{m_2,p_2}(\Omega)} \]
	Im dem Fall, dass die Ungleichung strikt ist, ist $W^{m_1,p_1}(\Omega)$ kompakte Teilmenge von $W^{m_2,p_2}(\Omega)$.
\end{thm}
	
\minisec{Beweis}
	siehe Alt, \glqq Lineare Funktionalanalysis\grqq, S. 328.
	
\begin{thm}[Einbettungssatz von Hölder] \label{thm:hoelder_einbettung} \label{thm_63}
	Sei $\Omega \subset \RR^n$ offen, beschränkt und mit Lipschitzrand, $m, k \in \NN_0, p \in [1,\infty)$ und $\alpha \in [0,1]$. Falls \marginnote{[63]}
	\[ m - \frac{n}{p} \geq k + \alpha \text{ und } \alpha \neq 0,1 \]
	dann ist $W^{m,p}(\Omega) \subset C^{k,\alpha}(\overline{\Omega})$, und es existiert eine Konstante $C > 0$, sodass für alle $u$ gilt:
	\[ ||u||_{W^{m,p}(\Omega)} \leq C \cdot ||u||_{C^{k,\alpha}(\overline{\Omega})} \]
	Falls $m - \frac{n}{p} < k + \alpha$, so ist $W^{m,p}(\Omega)$ kompakte Teilmenge von $C^{k,\alpha}(\overline{\Omega})$.
\end{thm}
	
\minisec{Beweis}
	siehe Alt, \glqq Lineare Funktionalanalysis\grqq, S. 333.

\newpage
\section{Der Vier-Quadrate-Satz}
\label{sec:para9}

\begin{satz}[Satz von Lagrange] \label{satz_langrange}
	Jede natürliche Zahl ist Summe von vier Quadraten, d.h. jedes $n \in \NN$ hat eine Darstellung der Gestalt \marginnote{23.01. \\ \ [26]}
	\begin{equation}
		n = x_1^2 + x_2^2 + x_3^2 + x_4^2 \text{ mit } x_i \in \ZZ \label{eq_4quad}
	\end{equation}
\end{satz}

\begin{defn}
	Für jedes $n \in \NN$ sei $r_4(n)$ die Anzahl der Quadrupel $(x_1,x_2,x_3,x_4) \in \ZZ^4$ mit \eqref{eq_4quad}.
\end{defn}

\begin{lemma}[Formel von Jacobi]
	Es gilt
	\begin{equation}
		r_4(n) = \begin{cases}
			8\sigma(n), & \text{falls } n \text{ ungerade} \\
			3\cdot 8 \sigma(u), & \text{falls } n \text{ gerade}
		\end{cases} \label{eq_jacobiformel}
	\end{equation}
	wobei $u$ den ungeraden Teil von $n$ bezeichnet, also $u = n \cdot 2^{-w_2(n)}$, und $\sigma$ die Teilersummenfunktion.
\end{lemma}

\subsection{Der Drei-Quadrate-Satz}

\begin{falko} \label{F10.1}
	Ist $n$ eine natürliche Zahl der Gestalt
	\begin{equation}
		n = 4^j \cdot (8k + 7) \text{ mit } j,k \in \NN_0, \label{eq_F10.1}
	\end{equation}
	so ist $n$ keine Summe von drei Quadraten in $\ZZ$.
\end{falko}

\minisec{Beweis}
Die Quadrate in $\ZZ/8\ZZ$ sind $0, 1$ und $4$, also ist $x^2+y^2+z^2 \not\kon 7 \modu 8$ für alle $x,y,z \in \ZZ$. Somit ist kein $n$ von der Gestalt $n = 8k+7$ eine Summe von drei Quadraten in $\ZZ$. 
Jetzt Induktion nach $j$: \\
Sei $j \geq 1$, dann ist $n = 4m$, und nach Induktionsanfang ist $m$ keine Summe von drei Quadraten. \\
Angenommen, es ist $4m = n = x^2+y^2+z^2$, dann folgt $x^2+y^2+z^2 \kon 0 \modu 4$. Dann sind $x,y$ und $z$ gerade, und es folgt $m = \leg{x}{2}^2 + \leg{y}{2}^2 + \leg{z}{2}^2 \quad \lightning$

\begin{satz}[Drei-Quadrate-Satz] \label{3quadsatz}
	Jede natürliche Zahl $n$, die nicht von der in F\ref{F10.1} genannten Gestalt \eqref{eq_F10.1} ist, ist eine Summe von drei Quadraten in $\ZZ$, d.h. $n$ hat eine Darstellung der Gestalt
	\begin{equation}
		n = x_1^2 + x_2^2 + x_3^2 \text{ mit } x_i \in \ZZ \label{eq_10_1}
	\end{equation}
	Ist zudem $n$ nicht durch $4$ teilbar, also $m \kon 1,2,3,5$ oder $6 \modu 8$, so existiert stets auch eine primitive Darstellung der Gestalt \eqref{eq_10_1}, d.h. eine mit
	\begin{equation}
		\ggT(x_1,x_2,x_3) = 1 \label{eq_10_2}
	\end{equation}
\end{satz}

Der Drei-Quadrate-Satz liegt anscheinend viel tiefer als der Vier-Quadrate-Satz, der aus dem Drei-Quadrate-Satz unmittelbar folgt: Sei $n \in \NN$ beliebig. Wir können annehmen, dass $n$ die Gestalt $n = 4^jm$ mit $m \kon 7 \modu 8$ hat (denn andernfalls ist $n$ ja schon eine Summe von drei Quadraten in $\ZZ$). Dann ist $m - 1 \kon 6 \modu 8$, also ist $m-1$ eine Summe von drei Quadraten und somit $m$ eine Summe von vier Quadraten. Aber dann ist auch $n = 4^jm$ eine Summe von vier Quadraten. \\
Der Drei-Quadrate-Satz folgt -- allerdings ohne den Zusatz \eqref{eq_10_2} -- aus einem allgemeinen Prinzip der Algebraischen Zahlentheorie ($\rightarrow$ Lokal-Global-Prinzip von Hasse). Was eine möglichst direkte Begründung des Drei-Quadrate-Satzes (ohne den Zusatz \eqref{eq_10_2}) angeht, so sind derzeit kursierende Beweise insofern "nicht elementar", als sie (anders als Gauß) den nachstehenden Satz von Dirichlet benutzen, dessen Beweis essentiell auf Methoden der komplexen Analysis beruht.

\begin{satz}[Satz von Dirichlet] \label{satz_dirichlet}
	Für beliebige $m > 1$ aus $\NN$ und jedes zu $m$ prime $a \in \ZZ$ gibt es unendlich viele Primzahlen $p$ mit $p \kon a \modu m$.
\end{satz}

\minisec{Bemerkungen}
\begin{enumerate}[1)]
	\item Auf den Zusatz \eqref{eq_10_2} im Drei-Quadrate-Satz legte schon Legendre Wert. Man beachte: Während z.B. $45 = 6^2+3^2$ im Wesentlichen die einzige Darstellung von $45$ als Summe von zwei Quadraten ist (und diese nicht primitiv ist), hat $45$ neben $45 = 6^2+3^2+0^2$ auch die primitive Darstellung $45 = 5^2+4^2+2^2$ als Summe von drei Quadraten.
	\item Gauß hat auch über die Anzahl $r_3(n)$ der Darstellungen $n = x_1^2 + x_2^2+x_3^2$ bzw. die Anzahl $R_3(n)$ der entsprechenden primitiven Darstellungen interessante Aussagen gemacht.
	\item Unabhängig von Gauß erhält man durch einfaches Abzählen der $(x_1,x_2,x_3,x_4)$ mit $n = x_1^2+x_2^2+x_3^2+x_4^2$ für $x_1 = 0, x_1=1, x_2 =2, \dots$ die banale Formel
	\begin{equation}
		r_4(n) = r_3(n) + 2r_3(n-1^2) + 2r_3(n-2^2) + 2r_3(n-3^2)+ \cdots \label{eq_10_3}
	\end{equation}
	Mit der Formel von Jacobi für $r_4(n)$ liefert \eqref{eq_10_3} eine gute Rekursionsformel für die $r_3(n)$. Aber dieser Rekursionsformel lässt sich nicht entnehmen -- jedenfalls nicht auf einfache Weise --, dass stets $r_3(n) \neq 0$, wenn $n$ nicht von der Gestalt $n = 4^j(8k+7)$ ist.
	\item In einer Extra-Vorlesung am 03. Februar 2015, 12:00 Uhr, wird Prof. Lorenz die Idee des Gaußschen Beweises für den Drei-Quadrate-Satzes vorstellen.
\end{enumerate}
\newpage

%\appendix							% Anhang beginnt hier (Nummerierung mit Buchstaben)
\printindex							% Index ausgeben
\subsection*{Liste der Sätze und Definitionen}
\listtheorems{satz,falko,defn,lemma}					% listet angegebene Sätze mit Namen und Seitenzahl auf
\end{document}